% Cálculo de Malliavin: de los procesos estocásticos y la integral de Itô al estado del arte.
% Fuente LaTeX RECONSTRUIDA a partir del PDF (docs/pdf/Calculo_de_Malliavin_corregido.pdf).
% Compilar:  latexmk -pdf main.tex
\documentclass[11pt,a4paper,twoside,openany]{book}
% Preámbulo reconstruido a partir del PDF compilado (pdfTeX, TeX Live 2023, Latin Modern).
\usepackage[utf8]{inputenc}
\usepackage[T1]{fontenc}
\usepackage{lmodern}
\usepackage[spanish,es-minimal,es-lcroman]{babel}
\usepackage{amsmath,amssymb,amsthm,mathtools,bm}
\usepackage{geometry}
\geometry{a4paper,margin=2.6cm,headheight=14pt}
\usepackage{enumitem}
\usepackage{booktabs}
\usepackage{fancyhdr}
\usepackage{xcolor}
\usepackage{tikz}
\usetikzlibrary{arrows.meta,decorations.pathmorphing}
\usepackage[most]{tcolorbox}
\usepackage[colorlinks=true,linkcolor=blue!60!black,citecolor=blue!60!black,urlcolor=blue!60!black]{hyperref}

\addto\captionsspanish{\renewcommand{\contentsname}{Índice}}

% ---- Cabeceras ----
\pagestyle{fancy}
\fancyhf{}
\fancyhead[LE]{\small\leftmark}
\fancyhead[RO]{\small\rightmark}
\fancyfoot[C]{\thepage}
\renewcommand{\headrulewidth}{0.4pt}
\fancypagestyle{plain}{\fancyhf{}\fancyfoot[C]{\thepage}\renewcommand{\headrulewidth}{0pt}}

% ---- Teoremas ----
\theoremstyle{plain}
\newtheorem{teorema}{Teorema}[chapter]
\newtheorem{proposicion}[teorema]{Proposición}
\newtheorem{lema}[teorema]{Lema}
\newtheorem{corolario}[teorema]{Corolario}
\theoremstyle{definition}
\newtheorem{definicion}[teorema]{Definición}
\newtheorem{ejemplo}[teorema]{Ejemplo}
\newtheorem{ejercicio}[teorema]{Ejercicio}
\theoremstyle{remark}
\newtheorem{observacion}[teorema]{Observación}

\definecolor{teoremaborde}{RGB}{60,60,140}
\definecolor{teoremafondo}{RGB}{246,246,252}
\definecolor{defborde}{RGB}{160,160,160}
\definecolor{deffondo}{RGB}{245,245,245}
\definecolor{panoramaborde}{RGB}{0,100,0}
\definecolor{panoramafondo}{RGB}{245,252,245}
\definecolor{intuicionborde}{RGB}{150,75,0}
\definecolor{intuicionfondo}{RGB}{255,250,240}

\tcolorboxenvironment{teorema}{enhanced,breakable,colback=teoremafondo,colframe=teoremaborde,
  boxrule=0.6pt,arc=1pt,left=4pt,right=4pt,top=2pt,bottom=2pt,before skip=8pt,after skip=8pt}
\tcolorboxenvironment{definicion}{enhanced,breakable,colback=deffondo,colframe=defborde,
  boxrule=0.4pt,arc=1pt,left=4pt,right=4pt,top=2pt,bottom=2pt,before skip=8pt,after skip=8pt}

\newtcolorbox{panorama}{enhanced,colback=panoramafondo,colframe=panoramaborde,boxrule=0.6pt,arc=1pt,
  title={Panorama del capítulo},fonttitle=\bfseries\small,left=6pt,right=6pt,before skip=10pt,after skip=12pt}
\newtcolorbox{intuicion}{enhanced,breakable,colback=intuicionfondo,colframe=intuicionborde,boxrule=0.6pt,arc=1pt,
  title={Intuición},fonttitle=\bfseries\small,left=6pt,right=6pt,before skip=10pt,after skip=10pt}

% ---- Notación ----
\newcommand{\R}{\mathbb{R}}
\newcommand{\N}{\mathbb{N}}
\newcommand{\C}{\mathbb{C}}
\newcommand{\Z}{\mathbb{Z}}
\newcommand{\E}{\mathbb{E}}
\renewcommand{\P}{\mathbb{P}}
\newcommand{\Q}{\mathbb{Q}}
\newcommand{\F}{\mathcal{F}}
\newcommand{\G}{\mathcal{G}}
\newcommand{\Hf}{\mathfrak{H}}
\newcommand{\Hc}{\mathcal{H}}
\newcommand{\Sc}{\mathcal{S}}
\newcommand{\D}{\mathbb{D}}
\newcommand{\ind}{\mathbf{1}}
\newcommand{\dd}{\mathrm{d}}
\newcommand{\Dom}{\operatorname{Dom}}
\newcommand{\Var}{\operatorname{Var}}
\newcommand{\Cov}{\operatorname{Cov}}
\newcommand{\Ent}{\operatorname{Ent}}
\newcommand{\tr}{\operatorname{tr}}
\newcommand{\sgn}{\operatorname{sgn}}
\newcommand{\Lip}{\operatorname{Lip}}
\newcommand{\ip}[2]{\left\langle #1, #2\right\rangle}
\newcommand{\norm}[1]{\left\lVert #1\right\rVert}
\newcommand{\abs}[1]{\left\lvert #1\right\rvert}
\newcommand{\cond}[2]{\E\left[#1\,\middle|\,#2\right]}
\newcommand{\law}{\overset{d}{=}}
\newcommand{\tolaw}{\xrightarrow{\ \mathrm{ley}\ }}


\begin{document}
\frontmatter
\begin{titlepage}
\centering
\vspace*{3.5cm}
{\Huge\bfseries Cálculo de Malliavin\par}
\vspace{1cm}
{\Large de los procesos estocásticos y la integral de Itô\\ al estado del arte\par}
\vspace{1.2cm}
\begin{tikzpicture}[x=0.75cm,y=0.75cm]
  \draw[->,gray] (0,0) -- (12,0) node[right] {\small $t$};
  \draw[->,gray] (0,-2.4) -- (0,2.6);
  \pgfmathsetseed{7}
  \draw[blue!70!black,thick] (0,0)
    \foreach \x in {1,...,47} { -- (\x/4, {0.9*sin(\x*37) + 0.6*sin(\x*11) - 0.03*\x + 0.04*(\x-20)*(\x>20)}) };
  \draw[red!80!black,thick,dashed] (0,0)
    \foreach \x in {1,...,47} { -- (\x/4, {0.9*sin(\x*37) + 0.6*sin(\x*11) - 0.03*\x + 0.04*(\x-20)*(\x>20) + 0.025*\x + 0.4*sin(\x*4)}) };
  \node[blue!70!black] at (2.8,-1.6) {$\omega$};
  \node[red!80!black] at (8.4,2.2) {$\omega+\varepsilon h$};
\end{tikzpicture}
\vspace{1.5cm}

{\large Una guía con demostraciones\par}
\vfill
{\large Noviembre de 2026\par}
\end{titlepage}

\chapter*{Prefacio}
\addcontentsline{toc}{chapter}{Prefacio}
\markboth{PREFACIO}{PREFACIO}

El cálculo de Malliavin es un cálculo diferencial en el espacio de trayectorias del movimiento
browniano. Nació en 1976, cuando Paul Malliavin dio una prueba probabilística del teorema
de hipoelipticidad de Hörmander, y ha crecido hasta convertirse en una de las herramientas
centrales del análisis estocástico: permite demostrar que las leyes de soluciones de ecuaciones
estocásticas tienen densidad regular, calcular sensibilidades de precios de derivados financieros
sin diferenciar funciones de pago discontinuas, obtener teoremas centrales del límite cuantitativos
(el método de Malliavin--Stein) y, más recientemente, construir y renormalizar soluciones de
ecuaciones en derivadas parciales estocásticas singulares.

Esta guía recorre ese camino completo. La Parte~I establece los fundamentos: martingalas,
movimiento browniano, integral de Itô y ecuaciones diferenciales estocásticas. La Parte~II
desarrolla el núcleo de la teoría en el marco de los procesos gaussianos isonormales: caos de
Wiener, la derivada $D$, la divergencia $\delta$, el operador de Ornstein--Uhlenbeck $L$ y las desigualdades
de Meyer. La Parte~III contiene la aplicación fundacional: criterios de existencia y regularidad
de densidades y el teorema de Hörmander. Las Partes~IV y~V desarrollan con detalle dos
grandes campos de aplicación, las finanzas matemáticas y las EDPEs. La Parte~VI presenta
extensiones (procesos gaussianos no markovianos y espacios de Poisson) y la Parte~VII describe
el estado del arte.

\paragraph{Convenciones.} Las demostraciones son completas para los resultados estructurales; cuando
un argumento es muy técnico se da un esbozo detallado y una referencia precisa. Los recuadros
\emph{Intuición} pueden saltarse en una primera lectura, pero recomendamos no hacerlo. Las constantes
$C$ pueden cambiar de línea a línea. Trabajamos sobre un espacio de probabilidad completo
$(\Omega,\F,\P)$.

\paragraph{Lecturas de referencia.} Las fuentes principales son Nualart~\cite{Nualart06}, Nualart--Nualart~\cite{NualartNualart18},
Nourdin--Peccati~\cite{NourdinPeccati12}, Karatzas--Shreve~\cite{KaratzasShreve91}, Revuz--Yor~\cite{RevuzYor99}, Da Prato--Zabczyk~\cite{DaPratoZabczyk14}, Walsh~\cite{Walsh86}
y Hairer~\cite{Hairer11}.

{\hypersetup{linkcolor=black}\tableofcontents}

\mainmatter
\part{Fundamentos: procesos estocásticos y cálculo de Itô}
\chapter{Preliminares de probabilidad}\label{cap:preliminares}

\begin{panorama}
Fijamos el lenguaje: esperanza condicional, filtraciones, tiempos de parada y martingalas
en tiempo discreto y continuo. Probamos las desigualdades de Doob y el teorema de parada
opcional, que se usarán constantemente. Cerramos con vectores y espacios gaussianos, el
objeto sobre el que se construirá todo el cálculo de Malliavin.
\end{panorama}

\section{Esperanza condicional}
Sea $(\Omega,\F,\P)$ un espacio de probabilidad y $\G\subset\F$ una sub-$\sigma$-álgebra.

\begin{teorema}[Existencia y unicidad]\label{teo:esp-cond}
Para cada $X\in L^1(\Omega,\F,\P)$ existe una variable
$Y\in L^1(\Omega,\G,\P)$, única salvo igualdad c.s., tal que
\[
  \E[X\ind_G]=\E[Y\ind_G]\qquad\text{para todo } G\in\G.
\]
Se denota $Y=\cond{X}{\G}$.
\end{teorema}

\begin{proof}
\emph{Caso $L^2$.} $L^2(\G)$ es un subespacio cerrado de $L^2(\F)$; sea $Y$ la proyección ortogonal
de $X$. Entonces $X-Y\perp\ind_G$ para todo $G\in\G$, que es la propiedad pedida.

\emph{Caso $X\ge0$.} Sea $X_n=X\wedge n\in L^2$ e $Y_n$ la proyección. La proyección conserva el orden
(si $Z\ge0$ y $Y=\cond{Z}{\G}$, tomando $G=\{Y<0\}$ obtenemos $\E[Y\ind_G]=\E[Z\ind_G]\ge0$, luego
$\P(G)=0$), así que $Y_n$ es creciente; $Y=\lim Y_n$ cumple la propiedad por convergencia monótona
e $\E Y=\E X<\infty$.

\emph{Caso general.} $X=X^+-X^-$ y linealidad. La unicidad: si $Y,Y'$ cumplen la propiedad, con
$G=\{Y>Y'\}\in\G$ resulta $\E[(Y-Y')\ind_G]=0$, luego $Y\le Y'$ c.s., y por simetría $Y=Y'$.
\end{proof}

Las propiedades habituales (linealidad, monotonía, torre $\cond{\cond{X}{\G}}{\mathcal K}=\cond{X}{\mathcal K}$ si $\mathcal K\subset\G$,
sacar lo conocido $\cond{ZX}{\G}=Z\cond{X}{\G}$ para $Z$ $\G$-medible acotada, Jensen condicional) se
deducen directamente de la definición. Destacamos que $\cond{\cdot}{\G}$ es una contracción en todo $L^p$,
$1\le p\le\infty$, por Jensen.

\section{Filtraciones, procesos y tiempos de parada}
Una \emph{filtración} es una familia creciente $(\F_t)_{t\ge0}$ de sub-$\sigma$-álgebras. Supondremos las \emph{condiciones
habituales}: $\F_0$ contiene los $\P$-nulos y $\F_t=\F_{t+}:=\bigcap_{s>t}\F_s$. Un proceso $X=(X_t)$ es \emph{adaptado}
si $X_t$ es $\F_t$-medible para todo $t$, y \emph{progresivamente medible} si para cada $t$ la aplicación $(s,\omega)\mapsto
X_s(\omega)$ en $[0,t]\times\Omega$ es $\mathcal B([0,t])\otimes\F_t$-medible. Todo proceso adaptado con trayectorias continuas
por la derecha es progresivo.

\begin{definicion}
Una variable $\tau:\Omega\to[0,\infty]$ es un \emph{tiempo de parada} si $\{\tau\le t\}\in\F_t$ para
todo $t$. La $\sigma$-álgebra de los sucesos anteriores a $\tau$ es $\F_\tau=\{A\in\F:\ A\cap\{\tau\le t\}\in\F_t\ \forall t\}$.
\end{definicion}

Si $X$ es progresivo y $\tau$ es de parada, $X_\tau\ind_{\{\tau<\infty\}}$ es $\F_\tau$-medible y el proceso parado $X^\tau_t=X_{t\wedge\tau}$
es progresivo. El tiempo de entrada de un proceso continuo adaptado en un cerrado es un
tiempo de parada; en un abierto, lo es bajo las condiciones habituales.

\section{Martingalas en tiempo discreto}

\begin{definicion}
Un proceso adaptado e integrable $(M_n)$ es una \emph{martingala} (resp.\ sub-,
super-) si $\cond{M_{n+1}}{\F_n}=M_n$ (resp.\ $\ge$, $\le$).
\end{definicion}

\begin{lema}[Transformada de martingala]\label{lem:transformada}
Si $(M_n)$ es martingala y $(H_n)$ es \textup{predecible} ($H_n$ es
$\F_{n-1}$-medible) y acotado, entonces $(H\cdot M)_n=\sum_{k=1}^nH_k(M_k-M_{k-1})$ es martingala.
\end{lema}

\begin{proof}
$\cond{H_{n+1}(M_{n+1}-M_n)}{\F_n}=H_{n+1}\cond{M_{n+1}-M_n}{\F_n}=0$.
\end{proof}

Este lema es el embrión de la integral estocástica: ``no se puede ganar a un juego justo con
una estrategia que no mira al futuro''.

\begin{teorema}[Parada opcional, versión acotada]\label{teo:parada}
Si $(M_n)$ es martingala y $\sigma\le\tau\le N$ son
tiempos de parada acotados, entonces $\cond{M_\tau}{\F_\sigma}=M_\sigma$.
\end{teorema}

\begin{proof}
Sea $A\in\F_\sigma$ y $H_k=\ind_A\ind_{\{\sigma<k\le\tau\}}$, que es predecible porque $A\cap\{\sigma<k\}\in\F_{k-1}$ y
$\{k\le\tau\}=\{\tau\le k-1\}^c\in\F_{k-1}$. Por el Lema~\ref{lem:transformada}, $\E[(H\cdot M)_N]=0$, es decir, $\E[\ind_A(M_\tau-M_\sigma)]=0$.
\end{proof}

\begin{teorema}[Desigualdades de Doob]\label{teo:doob}
Sea $(M_n)_{n\le N}$ una martingala o una submartingala
no negativa, y $M_N^*=\max_{n\le N}\abs{M_n}$. Entonces
\begin{enumerate}[label=(\roman*)]
  \item $\lambda\,\P(M_N^*\ge\lambda)\le\E\big[\abs{M_N}\ind_{\{M_N^*\ge\lambda\}}\big]\le\E\abs{M_N}$ para $\lambda>0$;
  \item $\norm{M_N^*}_p\le\frac{p}{p-1}\norm{M_N}_p$ para $1<p<\infty$.
\end{enumerate}
\end{teorema}

\begin{proof}
(i) $X_n=\abs{M_n}$ es submartingala no negativa. Sea $\tau=\min\{n:X_n\ge\lambda\}\wedge N$.
Por la versión submartingala de la parada opcional, $\E[X_\tau]\le\E[X_N]$. En $\{M_N^*\ge\lambda\}$, $X_\tau\ge\lambda$; en
el complementario, $\tau=N$. Así $\lambda\P(M_N^*\ge\lambda)+\E[X_N\ind_{\{M^*<\lambda\}}]\le\E X_\tau\le\E X_N$, y despejamos.

(ii) Por Fubini y (i),
\[
  \E[(M_N^*)^p]=\int_0^\infty p\lambda^{p-1}\P(M_N^*\ge\lambda)\,\dd\lambda\le\int_0^\infty p\lambda^{p-2}\E[X_N\ind_{\{M_N^*\ge\lambda\}}]\,\dd\lambda=\tfrac{p}{p-1}\E[X_N(M_N^*)^{p-1}].
\]
Por Hölder, $\E[X_N(M_N^*)^{p-1}]\le\norm{X_N}_p\norm{M_N^*}_p^{p-1}$. Si $\norm{M_N^*}_p<\infty$ (lo es, al ser un máximo finito
de variables en $L^p$), dividimos.
\end{proof}

\begin{teorema}[Convergencia de martingalas]
Una submartingala $(X_n)$ con $\sup_n\E X_n^+<\infty$
converge c.s.\ a un límite integrable. Si $(M_n)$ es martingala acotada en $L^p$, $p>1$, converge
c.s.\ y en $L^p$. Si es uniformemente integrable, converge c.s.\ y en $L^1$, y $M_n=\cond{M_\infty}{\F_n}$.
\end{teorema}

\begin{proof}[Esbozo]
La desigualdad de cruces de Doob, $(b-a)\E[U_N(a,b)]\le\E[(X_N-a)^+]$, donde $U_N(a,b)$
cuenta los cruces ascendentes de $[a,b]$, se prueba con la transformada $H$ que ``compra'' bajo $a$
y ``vende'' sobre $b$. Si el número de cruces es finito para todo par racional $a<b$, la sucesión
no puede oscilar, luego converge. La convergencia en $L^p$ se sigue de Doob (ii) y convergencia
dominada; la de $L^1$ de la integrabilidad uniforme y Vitali.
\end{proof}

\section{Martingalas en tiempo continuo}
Sea $(M_t)_{t\ge0}$ una martingala continua por la derecha. Aplicando el Teorema~\ref{teo:doob} en conjuntos
finitos $D_n\uparrow D=(\mathbb Q\cap[0,T])\cup\{T\}$ y la continuidad por la derecha:
\[
  \P\Big(\sup_{t\le T}\abs{M_t}\ge\lambda\Big)\le\frac{\E\abs{M_T}}{\lambda},\qquad
  \E\Big[\sup_{t\le T}\abs{M_t}^p\Big]\le\Big(\frac{p}{p-1}\Big)^p\E\abs{M_T}^p.
\]
La parada opcional (Teorema~\ref{teo:parada}) se extiende a tiempos acotados aproximando $\tau$ por $\tau_n=
2^{-n}\lceil2^n\tau\rceil$, usando continuidad por la derecha e integrabilidad uniforme de $\{M_{\tau_n}\}$ (que son
esperanzas condicionales de $M_T$).

\begin{definicion}
$M$ es una \emph{martingala local} si existen tiempos de parada $\tau_n\uparrow\infty$ tales
que cada $M^{\tau_n}$ es martingala. Denotamos $\mathcal M^2_c$ al espacio de martingalas continuas con
$\sup_t\E M_t^2<\infty$ y $M_0=0$.
\end{definicion}

\begin{proposicion}\label{prop:M2c}
$\mathcal M^2_c([0,T])$ con la norma $\norm{M}=\norm{M_T}_{L^2}$ es un espacio de Hilbert.
\end{proposicion}

\begin{proof}
Si $(M^n)$ es de Cauchy, por Doob $\E\sup_{t\le T}\abs{M^n_t-M^m_t}^2\le4\E\abs{M^n_T-M^m_T}^2$.
Extrayendo una subsucesión con $\sum_k\norm{M^{n_{k+1}}-M^{n_k}}<\infty$, la serie converge uniformemente en
$[0,T]$ c.s., luego el límite es continuo; es martingala por convergencia en $L^2$ de las esperanzas
condicionales.
\end{proof}

\section{Vectores y espacios gaussianos}

\begin{definicion}
Un vector $X=(X_1,\dots,X_d)$ es \emph{gaussiano centrado} con covarianza $\Sigma$ si
$\E e^{i\ip{u}{X}}=e^{-\frac12\ip{u}{\Sigma u}}$. Una familia $(X_i)_{i\in I}$ es \emph{gaussiana} si toda combinación lineal finita es
gaussiana. Un \emph{espacio gaussiano} es un subespacio cerrado de $L^2(\Omega)$ formado por variables
gaussianas centradas.
\end{definicion}

\begin{proposicion}\label{prop:gaussianos}
En un espacio gaussiano, ortogonalidad equivale a independencia: si
$(X_1,\dots,X_n,Y_1,\dots,Y_m)$ es gaussiano y $\Cov(X_i,Y_j)=0$ para todo $i,j$, los vectores $X$ e $Y$ son
independientes. Además, los límites en $L^2$ de variables gaussianas centradas son gaussianos
centrados, con convergencia en todo $L^p$.
\end{proposicion}

\begin{proof}
La función característica conjunta factoriza porque la forma cuadrática
$\ip{(u,v)}{\Sigma(u,v)}$ no tiene términos cruzados. Para el límite: si $X_n\to X$ en $L^2$ con
$X_n\sim N(0,\sigma_n^2)$, entonces $\sigma_n\to\sigma=\norm{X}_2$ y $e^{-\sigma_n^2u^2/2}\to\E e^{iuX}$. Los momentos de orden $p$
de una gaussiana son $c_p\sigma^p$, luego la sucesión está acotada en todo $L^p$ y hay integrabilidad
uniforme.
\end{proof}

\begin{lema}[Integración por partes gaussiana, o lema de Stein]\label{lem:stein}
Si $X\sim N(0,\sigma^2)$ y $f\in C^1$
con $f'$ de crecimiento polinomial, entonces $\E[Xf(X)]=\sigma^2\E[f'(X)]$. En dimensión $d$, si
$X\sim N(0,\Sigma)$, $\E[X_if(X)]=\sum_j\Sigma_{ij}\E[\partial_jf(X)]$.
\end{lema}

\begin{proof}
Con $\varphi_\sigma(x)=(2\pi\sigma^2)^{-1/2}e^{-x^2/2\sigma^2}$ se tiene $x\varphi_\sigma(x)=-\sigma^2\varphi_\sigma'(x)$; se integra por
partes. El caso multidimensional es análogo usando $\Sigma^{-1}$ si es invertible, y por aproximación en
general.
\end{proof}

\begin{intuicion}
El Lema~\ref{lem:stein} es, literalmente, el germen del cálculo de Malliavin. Lo que hará Malliavin es
desarrollar esta fórmula en el espacio de dimensión infinita de las trayectorias brownianas:
la ``variable gaussiana'' será la integral de Wiener $W(h)$ y la ``derivada'' será la derivada
de Malliavin. La identidad $\E[Xf(X)]=\E[f'(X)]$ se convertirá en la dualidad $\E[F\,\delta(u)]=
\E\ip{DF}{u}$, y también en la identidad de Stein que caracteriza la ley normal y que está en
la base del método de Malliavin--Stein (Capítulo~\ref{cap:malliavin-stein}).
\end{intuicion}

\begin{ejercicio}
Prueba la recíproca del Lema~\ref{lem:stein} (caracterización de Stein): si $\E[f'(X)-
Xf(X)]=0$ para toda $f\in C_b^1$, entonces $X\sim N(0,1)$. (Indicación: toma $f(x)=e^{iux}$ y obtén
una EDO para la función característica.)
\end{ejercicio}

\chapter{El movimiento browniano}\label{cap:browniano}

\begin{panorama}
Construimos el movimiento browniano de dos maneras (Lévy--Ciesielski y Kolmogorov),
estudiamos la regularidad de sus trayectorias (Hölder de orden $<1/2$, en ningún punto
diferenciable), su variación cuadrática y la propiedad de Markov fuerte. La construcción
de Lévy--Ciesielski es especialmente importante para nosotros: muestra que el browniano
es una \emph{serie aleatoria gaussiana} en una base ortonormal de $L^2$, que es exactamente el
punto de vista del espacio de Wiener abstracto.
\end{panorama}

\section{Definición}

\begin{definicion}
Un \emph{movimiento browniano} (estándar, unidimensional) es un proceso
$W=(W_t)_{t\ge0}$ tal que $W_0=0$; para $0\le t_0<\dots<t_n$ los incrementos $W_{t_k}-W_{t_{k-1}}$
son independientes con ley $N(0,t_k-t_{k-1})$; y casi todas las trayectorias son continuas. Equivalentemente,
$W$ es un proceso gaussiano centrado continuo con $\E[W_sW_t]=s\wedge t$. Es un
\emph{$(\F_t)$-browniano} si es adaptado y $W_t-W_s$ es independiente de $\F_s$ para $s<t$.
\end{definicion}

\section{Construcción de Lévy--Ciesielski}
Sea $\{e_n\}_{n\ge0}$ una base ortonormal de $L^2([0,1])$ y $\{\xi_n\}$ i.i.d.\ $N(0,1)$. Formalmente, el ``ruido
blanco'' es $\dot W=\sum_n\xi_ne_n$ y $W_t=\ip{\dot W}{\ind_{[0,t]}}=\sum_n\xi_n\int_0^te_n(s)\,\dd s$. La serie converge en $L^2(\Omega)$
para cada $t$, pues $\sum_n(\int_0^te_n)^2=\norm{\ind_{[0,t]}}^2=t$ por Parseval, y la covarianza es
\[
  \E[W_sW_t]=\sum_n\ip{e_n}{\ind_{[0,s]}}\ip{e_n}{\ind_{[0,t]}}=\ip{\ind_{[0,s]}}{\ind_{[0,t]}}=s\wedge t.
\]
Para obtener continuidad de trayectorias se elige la \emph{base de Haar}: $h_0\equiv1$, $h_{n,k}=
2^{n/2}(\ind_{[k2^{-n},(k+\frac12)2^{-n})}-\ind_{[(k+\frac12)2^{-n},(k+1)2^{-n})})$, $0\le k<2^n$. Sus primitivas $s_{n,k}(t)=\int_0^th_{n,k}$ son
las funciones de Schauder: ``tiendas'' de altura $2^{-n/2-1}$ con soportes disjuntos para $n$ fijo.

\begin{teorema}[Lévy--Ciesielski]
La serie $W_t=\xi_0t+\sum_{n\ge0}\sum_{k<2^n}\xi_{n,k}s_{n,k}(t)$ converge
uniformemente en $[0,1]$ c.s.; su suma es un movimiento browniano en $[0,1]$.
\end{teorema}

\begin{proof}
Sea $a_n=\max_k\abs{\xi_{n,k}}$. Como $\P(\abs{\xi}>x)\le e^{-x^2/2}$ para $x\ge1$, $\P(a_n>2\sqrt n)\le
2^ne^{-2n}$, sumable. Por Borel--Cantelli, c.s.\ $a_n\le2\sqrt n$ para $n$ grande. Como las tiendas de
nivel $n$ tienen soportes disjuntos, $\sup_t\abs{\sum_k\xi_{n,k}s_{n,k}(t)}\le a_n2^{-n/2-1}\le\sqrt n2^{-n/2}$, sumable. La
convergencia uniforme da continuidad; la convergencia en $L^2$ y el cálculo de covarianza anterior
dan la ley. Uniendo copias independientes en $[k,k+1]$ se obtiene el browniano en $[0,\infty)$.
\end{proof}

\begin{intuicion}
El browniano es una \emph{función lineal del ruido}: $W_t=W(\ind_{[0,t]})$, donde $h\mapsto W(h)=
\sum_n\xi_n\ip{e_n}{h}$ es una isometría de $L^2([0,1])$ en un espacio gaussiano. Esta isometría, la
\emph{integral de Wiener} $W(h)=\int_0^1h(s)\,\dd W_s$, será el objeto primitivo de la Parte~II, bajo el
nombre de \emph{proceso gaussiano isonormal}.
\end{intuicion}

\section{Regularidad de las trayectorias}

\begin{teorema}[Criterio de continuidad de Kolmogorov--Chentsov]\label{teo:kolmogorov}
Sea $(X_t)_{t\in[0,1]^d}$ un proceso
con valores en un espacio métrico completo tal que $\E[\rho(X_s,X_t)^p]\le C\abs{s-t}^{d+\beta}$ para ciertos
$p,\beta>0$. Entonces $X$ admite una modificación cuyas trayectorias son $\gamma$-Hölder para todo
$\gamma<\beta/p$.
\end{teorema}

\begin{proof}
Lo hacemos para $d=1$. Sea $D_n=\{k2^{-n}\}$. Por Chebyshev,
\[
  \P\Big(\max_k\rho(X_{(k+1)2^{-n}},X_{k2^{-n}})>2^{-\gamma n}\Big)\le2^n\cdot C2^{-n(1+\beta)}2^{\gamma pn}=C2^{-n(\beta-\gamma p)},
\]
sumable si $\gamma<\beta/p$. Por Borel--Cantelli existe $N(\omega)$ con $\rho(X_{(k+1)2^{-n}},X_{k2^{-n}})\le2^{-\gamma n}$ para
$n\ge N$. Si $s<t$ son diádicos con $2^{-m-1}<t-s\le2^{-m}$, $m\ge N$, se puede ir de $s$ a $t$ con a lo
sumo dos pasos de cada escala $n>m$ (expansión binaria), luego $\rho(X_s,X_t)\le2\sum_{n>m}2^{-\gamma n}\le
C'\abs{t-s}^\gamma$. Así $X|_D$ es uniformemente Hölder y se extiende por continuidad; la extensión es
una modificación por continuidad en probabilidad.
\end{proof}

Para el browniano, $\E\abs{W_t-W_s}^{2m}=c_m\abs{t-s}^m$, luego $\beta/p=(m-1)/2m\to1/2$.

\begin{corolario}
Las trayectorias brownianas son localmente $\gamma$-Hölder para todo $\gamma<1/2$.
\end{corolario}

\begin{teorema}[Paley--Wiener--Zygmund]
Casi seguramente, $t\mapsto W_t$ no es diferenciable en
ningún punto.
\end{teorema}

\begin{proof}
Si $W$ fuera diferenciable en $s\in[0,1]$, existirían $M,n_0$ tales que $\abs{W_t-W_s}\le
M\abs{t-s}$ para $\abs{t-s}\le4/n$, $n\ge n_0$. Con $k=\lfloor ns\rfloor$, los tres incrementos $\abs{W_{(j+1)/n}-W_{j/n}}$,
$j=k+1,k+2,k+3$, serían $\le7M/n$. Pero
\[
  \P\Big(\exists k\le n:\ \max_{j=1,2,3}\abs{W_{(k+j+1)/n}-W_{(k+j)/n}}\le\tfrac{7M}{n}\Big)\le n\,\P\big(\abs{W_{1/n}}\le\tfrac{7M}{n}\big)^3\le n\Big(\frac{C}{\sqrt n}\Big)^3\to0.
  \qedhere
\]
\end{proof}

\section{Variación cuadrática}

\begin{teorema}
Sea $\pi_n=\{0=t_0<\dots<t_{k_n}=t\}$ una sucesión de particiones con $\abs{\pi_n}\to0$.
Entonces $\sum_i(W_{t_{i+1}}-W_{t_i})^2\to t$ en $L^2$. Si $\sum_n\abs{\pi_n}<\infty$, la convergencia es casi segura.
\end{teorema}

\begin{proof}
Sea $\Delta_i=W_{t_{i+1}}-W_{t_i}$ y $\delta_i=t_{i+1}-t_i$. Las variables $\Delta_i^2-\delta_i$ son independientes,
centradas y $\Var(\Delta_i^2)=2\delta_i^2$. Así
\[
  \E\Big(\sum_i\Delta_i^2-t\Big)^2=2\sum_i\delta_i^2\le2t\abs{\pi_n}\to0.
\]
La convergencia casi segura sigue de Chebyshev y Borel--Cantelli.
\end{proof}

\begin{intuicion}
``$(\dd W_t)^2=\dd t$''. Esto es lo que distingue el cálculo de Itô del clásico: el término de
segundo orden en el desarrollo de Taylor no desaparece. Consecuencia: el browniano
tiene variación total infinita en todo intervalo, y la integral $\int H\,\dd W$ no puede definirse
trayectoria a trayectoria como integral de Stieltjes.
\end{intuicion}

\section{Propiedad de Markov fuerte y principio de reflexión}

\begin{teorema}[Markov fuerte]
Si $\tau$ es un tiempo de parada finito c.s., $B_t=W_{\tau+t}-W_\tau$ es
un movimiento browniano independiente de $\F_\tau$.
\end{teorema}

\begin{proof}[Esbozo]
Para $\tau$ con valores en un conjunto numerable se descompone según $\{\tau=t_k\}$ y se usa
la propiedad de Markov simple. Para $\tau$ general se aproxima por $\tau_n=2^{-n}\lceil2^n\tau\rceil\downarrow\tau$ y se usa la
continuidad de trayectorias.
\end{proof}

\begin{corolario}[Principio de reflexión]
Para $a>0$, $\P(\sup_{s\le t}W_s\ge a)=2\P(W_t\ge a)=
\P(\abs{W_t}\ge a)$.
\end{corolario}

\begin{proof}
Sea $\tau_a=\inf\{s:W_s=a\}$. El proceso reflejado en $\tau_a$ es también browniano
(Markov fuerte), luego $\P(\tau_a\le t,W_t<a)=\P(\tau_a\le t,W_t>a)=\P(W_t>a)$.
\end{proof}

\section{El espacio de Wiener clásico}
Sea $\Omega=C_0([0,T];\R^d)$ con la norma del supremo, $\F$ su $\sigma$-álgebra boreliana (completada) y $\P$ la
\emph{medida de Wiener}, ley del browniano $d$-dimensional; el proceso canónico $W_t(\omega)=\omega(t)$ es un
browniano. Este es el escenario original de Malliavin.

\begin{definicion}
El \emph{espacio de Cameron--Martin} es $H=\{h\in\Omega:\ h(t)=\int_0^t\dot h(s)\,\dd s,\ \dot h\in
L^2([0,T];\R^d)\}$ con $\ip{h}{g}_H=\ip{\dot h}{\dot g}_{L^2}$.
\end{definicion}

\begin{teorema}[Cameron--Martin]\label{teo:cameron-martin}
Para $h\in H$, la ley de $\omega+h$ bajo $\P$ es equivalente a $\P$
con densidad
\[
  \frac{\dd\P\circ\tau_h^{-1}}{\dd\P}(\omega)=\exp\Big(\int_0^T\dot h\,\dd W-\tfrac12\norm{h}_H^2\Big),\qquad\tau_h(\omega)=\omega+h.
\]
Si $h\notin H$, las leyes son mutuamente singulares.
\end{teorema}

\begin{proof}[Demostración de la primera parte]
Basta comparar las leyes de los vectores $(W_{t_1},\dots,W_{t_n})$, o
mejor, las transformadas de Laplace de $W(g)=\int\dot g\,\dd W$ para $g\in H$. Con $Z=\exp(W(h)-
\frac12\norm{h}^2)$:
\[
  \E[Ze^{W(g)}]=e^{-\frac12\norm{h}^2}\E e^{W(h+g)}=e^{-\frac12\norm{h}^2+\frac12\norm{h+g}^2}=e^{\ip{h}{g}+\frac12\norm{g}^2}=\E[e^{W(g)+\ip{g}{h}}],
\]
que es la transformada de $W(g)$ evaluada en $\omega+h$. Como las exponenciales $e^{W(g)}$ determinan
la medida, $Z\,\dd\P=\P\circ\tau_h^{-1}$. La singularidad para $h\notin H$ es la dicotomía de Feldman--Hájek; no
la usaremos.
\end{proof}

\begin{intuicion}
El espacio de Cameron--Martin $H$ es el conjunto de \emph{direcciones admisibles}: desplazar el
browniano en una dirección $h\in H$ cambia la medida de manera absolutamente continua.
Es, por tanto, la dirección natural en la que derivar: la derivada de Malliavin de $F$ en la
dirección $h$ será
\[
  D_hF(\omega)=\lim_{\varepsilon\to0}\frac{F(\omega+\varepsilon h)-F(\omega)}{\varepsilon}.
\]
Notemos que $\P(H)=0$: las trayectorias brownianas no son absolutamente continuas.
Derivamos en direcciones que el proceso casi nunca visita.
\end{intuicion}

\chapter{El cálculo de Itô}\label{cap:ito}

\begin{panorama}
Construimos la integral estocástica $\int_0^tH_s\,\dd W_s$ para integrandos progresivos, demostramos
la isometría de Itô, la fórmula de Itô, la desigualdad de Burkholder--Davis--Gundy, el
teorema de Girsanov y el teorema de representación de martingalas. Este último plantea la
pregunta que el cálculo de Malliavin responderá en el Capítulo~\ref{cap:divergencia}: \emph{¿cuál es explícitamente
el integrando?}
\end{panorama}

Fijamos un $(\F_t)$-browniano $W$ en $[0,T]$.

\section{La integral para procesos elementales}

\begin{definicion}
Un proceso \emph{elemental} es $H_t=\sum_{i=0}^{n-1}\xi_i\ind_{(t_i,t_{i+1}]}(t)$ con $0=t_0<\dots<t_n=T$
y $\xi_i\in L^2(\F_{t_i})$. Su integral es
\[
  I(H)_t=\int_0^tH_s\,\dd W_s:=\sum_i\xi_i(W_{t_{i+1}\wedge t}-W_{t_i\wedge t}).
\]
\end{definicion}

\begin{lema}[Isometría de Itô para elementales]\label{lem:isometria-elemental}
$I(H)$ es una martingala continua de cuadrado
integrable y $\E[I(H)_T^2]=\E\int_0^TH_s^2\,\dd s$.
\end{lema}

\begin{proof}
La propiedad de martingala es el Lema~\ref{lem:transformada} en tiempo continuo:
\[
  \cond{\xi_i(W_{t_{i+1}}-W_{t_i})}{\F_{t_i}}=\xi_i\cond{W_{t_{i+1}}-W_{t_i}}{\F_{t_i}}=0.
\]
Para la isometría, desarrollamos el cuadrado. Los términos cruzados, $i<j$, se anulan:
condicionando a $\F_{t_j}$, $\E[\xi_i\Delta_i\xi_j\Delta_j]=\E[\xi_i\Delta_i\xi_j\cond{\Delta_j}{\F_{t_j}}]=0$. Los diagonales: $\E[\xi_i^2\Delta_i^2]=
\E[\xi_i^2\cond{\Delta_i^2}{\F_{t_i}}]=\E[\xi_i^2](t_{i+1}-t_i)$.
\end{proof}

\section{Extensión por densidad}
Sea $L^2_{\mathrm{prog}}=L^2_{\mathrm{prog}}(\Omega\times[0,T])$ el espacio de procesos progresivos con $\E\int_0^TH_s^2\,\dd s<\infty$.

\begin{lema}\label{lem:densidad-elementales}
Los procesos elementales son densos en $L^2_{\mathrm{prog}}$.
\end{lema}

\begin{proof}
(a) Si $H$ es continuo y acotado, $H^n=\sum_iH_{t_i}\ind_{(t_i,t_{i+1}]}$ converge por convergencia
dominada. (b) Si $H$ es acotado y progresivo, $\tilde H^m_t=m\int_{(t-1/m)\vee0}^tH_s\,\dd s$ es continuo, adaptado y
converge a $H$ en $L^2(\dd t\otimes\dd\P)$ por el teorema de diferenciación de Lebesgue. (c) Truncación:
$H\ind_{\abs{H}\le n}\to H$.
\end{proof}

\begin{teorema}[Integral de Itô]\label{teo:integral-ito}
La aplicación $H\mapsto I(H)$ se extiende de manera única a
una isometría
\[
  I:L^2_{\mathrm{prog}}\to\mathcal M^2_c([0,T]),\qquad\E\Big[\Big(\int_0^TH\,\dd W\Big)^2\Big]=\E\int_0^TH_s^2\,\dd s.
\]
\end{teorema}

\begin{proof}
Lemas~\ref{lem:isometria-elemental} y~\ref{lem:densidad-elementales} y la completitud de $\mathcal M^2_c$ (Proposición~\ref{prop:M2c}).
\end{proof}

\paragraph{Localización.} Si $H$ es progresivo con $\int_0^TH_s^2\,\dd s<\infty$ c.s., sea $\tau_n=\inf\{t:\int_0^tH^2\ge n\}\wedge T$.
Entonces $H\ind_{[0,\tau_n]}\in L^2_{\mathrm{prog}}$ y las integrales son compatibles, lo que define $\int_0^tH\,\dd W$ como
\emph{martingala local} continua.

\begin{proposicion}[Variación cuadrática de la integral]
Si $M_t=\int_0^tH\,\dd W$ con $H\in L^2_{\mathrm{prog}}$, el
proceso $M_t^2-\int_0^tH_s^2\,\dd s$ es martingala, es decir, $\langle M\rangle_t=\int_0^tH_s^2\,\dd s$. Además $\sum_i(M_{t_{i+1}}-M_{t_i})^2\to
\langle M\rangle_t$ en probabilidad.
\end{proposicion}

\section{Semimartingalas y la fórmula de Itô}
Un \emph{proceso de Itô} es $X_t=X_0+\int_0^tb_s\,\dd s+\int_0^t\sigma_s\,\dd W_s$ con $b,\sigma$ progresivos, $\int\abs{b}+\int\sigma^2<\infty$ c.s.
Escribimos $\dd X=b\,\dd t+\sigma\,\dd W$ y $\dd\langle X\rangle=\sigma^2\,\dd t$.

\begin{teorema}[Fórmula de Itô]\label{teo:formula-ito}
Si $f\in C^{1,2}([0,T]\times\R)$ y $X$ es un proceso de Itô,
\[
  f(t,X_t)=f(0,X_0)+\int_0^t\partial_sf(s,X_s)\,\dd s+\int_0^t\partial_xf(s,X_s)\,\dd X_s+\frac12\int_0^t\partial_{xx}f(s,X_s)\sigma_s^2\,\dd s.
\]
En dimensión $d$, con $\dd X^i=b^i\,\dd t+\sum_k\sigma^{ik}\,\dd W^k$ y $\dd\langle X^i,X^j\rangle=(\sigma\sigma^\top)^{ij}\,\dd t$,
\[
  \dd f(t,X_t)=\partial_tf\,\dd t+\sum_i\partial_if\,\dd X^i_t+\frac12\sum_{i,j}\partial_{ij}f\,\dd\langle X^i,X^j\rangle_t.
\]
\end{teorema}

\begin{proof}
Por localización podemos suponer $X$, $b$, $\sigma$ acotados y $f$ con soporte compacto;
además omitimos la dependencia en $t$, que añade un término de primer orden inocuo. Sea
$\pi=\{t_i\}$ una partición de $[0,t]$. Por Taylor,
\[
  f(X_t)-f(X_0)=\sum_if'(X_{t_i})\Delta X_i+\frac12\sum_if''(X_{t_i})(\Delta X_i)^2+\sum_iR_i,
\]
con $\abs{R_i}\le\omega_{f''}(\abs{\Delta X_i})(\Delta X_i)^2$, siendo $\omega_{f''}$ el módulo de continuidad de $f''$.

\emph{Primer término.} $\sum_if'(X_{t_i})\Delta X_i\to\int_0^tf'(X_s)\,\dd X_s$ en probabilidad: la parte $\dd s$ por convergencia
dominada; la parte $\dd W$ porque el integrando elemental $\sum f'(X_{t_i})\ind_{(t_i,t_{i+1}]}$ converge en
$L^2_{\mathrm{prog}}$, por continuidad, y por la isometría.

\emph{Segundo término.} $(\Delta X_i)^2=(\int b)^2+2\int b\int\sigma\,\dd W+(\int\sigma\,\dd W)^2$. Los dos primeros suman
$O(\abs{\pi})\cdot(\text{acotado})\to0$. Para el tercero, $\sum_if''(X_{t_i})(\Delta M_i)^2-\sum_if''(X_{t_i})\Delta\langle M\rangle_i\to0$ en
$L^2$: sus sumandos son incrementos de martingala ortogonales, y $\E[((\Delta M_i)^2-\Delta\langle M\rangle_i)^2]\le
C\E[(\Delta\langle M\rangle_i)^2]\le C\abs{\pi}\delta_i$, por BDG o un cálculo directo con $\sigma$ acotado. Finalmente
$\sum f''(X_{t_i})\Delta\langle M\rangle_i\to\int f''(X_s)\sigma_s^2\,\dd s$.

\emph{Resto.} $\sum\abs{R_i}\le\max_i\omega_{f''}(\abs{\Delta X_i})\sum(\Delta X_i)^2\to0$ en probabilidad, pues el máximo tiende a $0$
por continuidad uniforme y la suma está acotada en probabilidad.

Pasando a una subsucesión con convergencia casi segura y usando la continuidad de ambos
lados en $t$, la igualdad vale para todo $t$ c.s.
\end{proof}

\begin{ejemplo}
$W_t^2=2\int_0^tW_s\,\dd W_s+t$. La exponencial estocástica $\mathcal E(h)_t=\exp(\int_0^th\,\dd W-
\frac12\int_0^th^2\,\dd s)$ satisface $\dd\mathcal E_t=h_t\mathcal E_t\,\dd W_t$.
\end{ejemplo}

\section{Desigualdades de Burkholder--Davis--Gundy}

\begin{teorema}[BDG]\label{teo:bdg}
Para cada $p>0$ existen constantes $c_p,C_p$ tales que, para toda
martingala local continua $M$ con $M_0=0$ y todo tiempo de parada $\tau$,
\[
  c_p\E[\langle M\rangle_\tau^{p/2}]\le\E\Big[\sup_{t\le\tau}\abs{M_t}^p\Big]\le C_p\E[\langle M\rangle_\tau^{p/2}].
\]
\end{teorema}

\begin{proof}[Demostración para $p\ge2$, cota superior]
Por localización, $M$ y $\langle M\rangle$ acotados. Por Itô,
$\abs{M_t}^p=\int_0^tp\abs{M_s}^{p-1}\sgn(M_s)\,\dd M_s+\frac{p(p-1)}{2}\int_0^t\abs{M_s}^{p-2}\,\dd\langle M\rangle_s$. Tomando esperanzas y con
$M^*=\sup_{s\le\tau}\abs{M_s}$,
\[
  \E\abs{M_\tau}^p\le\tfrac{p(p-1)}{2}\E\big[(M^*)^{p-2}\langle M\rangle_\tau\big]\le\tfrac{p(p-1)}{2}\norm{M^*}_p^{p-2}\norm{\langle M\rangle_\tau}_{p/2},
\]
por Hölder con exponentes $\frac{p}{p-2}$ y $\frac p2$. Por Doob, $\norm{M^*}_p^p\le(\frac{p}{p-1})^p\E\abs{M_\tau}^p$. Combinando y
dividiendo entre $\norm{M^*}_p^{p-2}$ se obtiene $\norm{M^*}_p^2\le C\norm{\langle M\rangle_\tau}_{p/2}$. Los demás casos se prueban con
argumentos de ``buen $\lambda$''; véase \cite[Th.~IV.4.1]{RevuzYor99}.
\end{proof}

\section{Teorema de Girsanov}

\begin{teorema}[Girsanov]\label{teo:girsanov}
Sea $\theta$ progresivo tal que $Z_t=\mathcal E(\theta)_t$ es martingala en $[0,T]$, lo cual
ocurre por ejemplo bajo la condición de Novikov, $\E\exp(\frac12\int_0^T\theta^2)<\infty$. Bajo $\dd\Q=Z_T\,\dd\P$, el
proceso $\tilde W_t=W_t-\int_0^t\theta_s\,\dd s$ es un browniano.
\end{teorema}

\begin{proof}
Por la caracterización de Lévy (una martingala local continua $M$ con $\langle M\rangle_t=t$
es browniana), basta ver que $\tilde W$ y $\tilde W^2-t$ son $\Q$-martingalas locales. Por la regla de Bayes, $X$ es
$\Q$-martingala local si $XZ$ es $\P$-martingala local. Por Itô, $\dd(\tilde WZ)=\tilde W\,\dd Z+Z\,\dd\tilde W+\dd\langle\tilde W,Z\rangle=
\tilde WZ\theta\,\dd W+Z\,\dd W-Z\theta\,\dd t+Z\theta\,\dd t$, que no tiene término en $\dd t$. Para $\tilde W^2-t$ el cálculo es
análogo.
\end{proof}

El Teorema de Cameron--Martin (\ref{teo:cameron-martin}) es el caso particular $\theta=\dot h$ determinista.

\section{Representación de martingalas}

\begin{teorema}[Representación de Itô]\label{teo:representacion}
Sea $\F_t=\F^W_t$ la filtración (completada) del
browniano. Toda $F\in L^2(\Omega,\F_T,\P)$ se escribe de manera única como
\[
  F=\E F+\int_0^Tu_s\,\dd W_s,\qquad u\in L^2_{\mathrm{prog}}.
\]
En consecuencia, toda $(\F^W_t)$-martingala de cuadrado integrable tiene una versión continua y
es una integral estocástica.
\end{teorema}

\begin{proof}
Sea $\mathcal R=\{\E F+\int_0^Tu\,\dd W:\ u\in L^2_{\mathrm{prog}}\}$. Por la isometría, $\mathcal R$ es cerrado en
$L^2$. Contiene las exponenciales $\mathcal E(h)_T=e^{W(h)-\frac12\norm{h}^2}$, con $h\in L^2([0,T])$ determinista, porque
$\mathcal E(h)_T=1+\int_0^Th_s\mathcal E(h)_s\,\dd W_s$. Basta ver que estas exponenciales generan un subespacio denso. Si
$G\in L^2(\F_T)$ es ortogonal a todas, en particular a $\exp(\sum_j\lambda_jW_{t_j})$ para todo $\lambda\in\R^n$, entonces la
medida (con signo) $G\,\dd\P$ restringida a $\sigma(W_{t_1},\dots,W_{t_n})$ tiene transformada de Laplace nula, que
es analítica en $\lambda$; por tanto se anula y $\cond{G}{W_{t_1},\dots,W_{t_n}}=0$. Por el teorema de convergencia
de martingalas, $G=0$. La unicidad es la isometría.
\end{proof}

\begin{intuicion}
El teorema es \emph{no constructivo}: prueba que $u$ existe por un argumento de densidad, pero
no dice cuál es. La \textbf{fórmula de Clark--Ocone} (Teorema~\ref{teo:clark-ocone}) dará la respuesta:
\[
  u_t=\cond{D_tF}{\F_t},
\]
donde $D_tF$ es la derivada de Malliavin. En finanzas, $u$ es la cartera de cobertura
(Capítulo~\ref{cap:cobertura}).
\end{intuicion}

\chapter{Ecuaciones diferenciales estocásticas}\label{cap:ede}

\begin{panorama}
Probamos existencia y unicidad fuerte bajo condiciones de Lipschitz, estimaciones de
momentos, la propiedad de Markov y la relación con el generador infinitesimal (Feynman--Kac).
Introducimos el \emph{proceso tangente} o primera variación $Y_t=\partial X_t/\partial x$, que reaparecerá
en la derivada de Malliavin de la solución (Capítulo~\ref{cap:hormander}). Usaremos indistintamente las
formulaciones de Itô y Stratonovich.
\end{panorama}

Sea $W=(W^1,\dots,W^m)$ un browniano $m$-dimensional, $b:\R^d\to\R^d$ y $\sigma=(\sigma_1,\dots,\sigma_m):
\R^d\to\R^{d\times m}$. Consideramos
\begin{equation}\label{eq:ede}
  X_t=x+\int_0^tb(X_s)\,\dd s+\sum_{k=1}^m\int_0^t\sigma_k(X_s)\,\dd W^k_s.
\end{equation}

\section{Existencia y unicidad}

\begin{teorema}\label{teo:existencia-ede}
Si $b,\sigma$ son globalmente Lipschitz, $\abs{b(x)-b(y)}+\abs{\sigma(x)-\sigma(y)}\le K\abs{x-y}$,
la ecuación~\eqref{eq:ede} tiene una única solución fuerte continua y adaptada. Además, para todo
$p\ge2$,
\[
  \E\Big[\sup_{t\le T}\abs{X_t}^p\Big]\le C(1+\abs{x}^p),\qquad\E\abs{X_t-X_s}^p\le C\abs{t-s}^{p/2}.
\]
\end{teorema}

\begin{proof}
Sea $\mathcal S^p$ el espacio de procesos continuos adaptados con $\norm{X}_{\mathcal S^p}^p=\E\sup_{t\le T}\abs{X_t}^p<
\infty$, y $\Phi(X)_t=x+\int_0^tb(X)\,\dd s+\int_0^t\sigma(X)\,\dd W$. Por Hölder y BDG (Teorema~\ref{teo:bdg}),
\[
  \E\sup_{s\le t}\abs{\Phi(X)_s-\Phi(Y)_s}^p\le C_{p,T}K^p\int_0^t\E\sup_{r\le s}\abs{X_r-Y_r}^p\,\dd s.
\]
Iterando, $\norm{\Phi^n(X)-\Phi^n(Y)}_{\mathcal S^p}^p\le\frac{(CT)^n}{n!}\norm{X-Y}_{\mathcal S^p}^p$, luego alguna iterada de $\Phi$ es contractiva y
hay un único punto fijo (iteración de Picard). Las cotas de momentos se obtienen con la misma
desigualdad aplicada a $\Phi(X)-\Phi(0)$ y el lema de Gronwall.
\end{proof}

\begin{lema}[Gronwall]
Si $\phi\ge0$ es acotada y $\phi(t)\le a+b\int_0^t\phi(s)\,\dd s$, entonces $\phi(t)\le ae^{bt}$.
\end{lema}

\section{Flujos y diferenciabilidad respecto al dato inicial}
Sea $X_t^x$ la solución que parte de $x$. Por el Teorema~\ref{teo:kolmogorov} aplicado en $(t,x)$, existe una versión
continua en $(t,x)$. Si $b,\sigma\in C^1_b$, formalmente $Y_t=\nabla_xX_t^x$ ($d\times d$) satisface la \emph{ecuación de primera
variación}
\begin{equation}\label{eq:primera-variacion}
  Y_t=I+\int_0^t\nabla b(X_s)Y_s\,\dd s+\sum_k\int_0^t\nabla\sigma_k(X_s)Y_s\,\dd W^k_s.
\end{equation}

\begin{teorema}[Kunita]\label{teo:kunita}
Si $b,\sigma\in C^\infty_b$ (derivadas de todo orden acotadas, la propia función
con crecimiento lineal), existe una versión de $(t,x)\mapsto X_t^x$ que es, c.s., un flujo de
difeomorfismos $C^\infty$ de $\R^d$. Su jacobiano $Y_t$ satisface~\eqref{eq:primera-variacion}, es invertible, y $Z_t=Y_t^{-1}$ satisface
\[
  Z_t=I-\int_0^tZ_s\Big(\nabla b-\sum_k(\nabla\sigma_k)^2\Big)(X_s)\,\dd s-\sum_k\int_0^tZ_s\nabla\sigma_k(X_s)\,\dd W^k_s.
\]
Además $\sup_{t\le T}(\abs{Y_t}+\abs{Z_t})\in L^p$ para todo $p$.
\end{teorema}

\begin{proof}[Comprobación de la ecuación de $Z$]
Por Itô, $\dd(Z_tY_t)=\dd Z\,Y+Z\,\dd Y+\dd\langle Z,Y\rangle$. Sustituyendo,
los términos en $\dd W^k$ se cancelan, $-Z\nabla\sigma_kY+Z\nabla\sigma_kY=0$, y los términos en $\dd t$ son $-Z(\nabla b-
\sum_k(\nabla\sigma_k)^2)Y+Z\nabla bY-\sum_kZ(\nabla\sigma_k)^2Y=0$. Así $ZY\equiv I$. La diferenciabilidad y la propiedad
de flujo se prueban derivando en $L^p$ la ecuación respecto a $x$ y aplicando Kolmogorov; véase
\cite{Kunita90}.
\end{proof}

\section{Formulación de Stratonovich}
La integral de Stratonovich $\int H\circ\dd W=\int H\,\dd W+\frac12\langle H,W\rangle$ obedece la regla de la cadena
clásica. La ecuación~\eqref{eq:ede} se reescribe
\[
  \dd X_t=V_0(X_t)\,\dd t+\sum_{k=1}^mV_k(X_t)\circ\dd W^k_t,\qquad V_k=\sigma_k,\quad V_0=b-\frac12\sum_k(\nabla\sigma_k)\sigma_k.
\]
Escribiremos los campos vectoriales como operadores diferenciales de primer orden, $V=\sum_iV^i\partial_i$.
El generador es entonces
\[
  \mathcal L=V_0+\frac12\sum_{k=1}^mV_k^2,
\]
forma de ``suma de cuadrados'' de Hörmander, fundamental en el Capítulo~\ref{cap:hormander}.

\section{Propiedad de Markov, generador y Feynman--Kac}

\begin{teorema}
La solución es un proceso de Markov fuerte con semigrupo $P_tf(x)=\E f(X_t^x)$,
y para $f\in C_b^2$,
\[
  \frac{\dd}{\dd t}P_tf\Big|_{t=0}=\mathcal Lf=\sum_ib^i\partial_if+\frac12\sum_{i,j}(\sigma\sigma^\top)^{ij}\partial_{ij}f.
\]
\end{teorema}

\begin{teorema}[Feynman--Kac]
Sea $u\in C^{1,2}([0,T)\times\R^d)\cap C([0,T]\times\R^d)$ de crecimiento
polinomial solución de
\[
  \partial_tu+\mathcal Lu-ru=0,\qquad u(T,\cdot)=g,
\]
con $r\ge0$ continua y acotada. Entonces $u(t,x)=\E\big[e^{-\int_t^Tr(X_s)\,\dd s}g(X_T)\,\big|\,X_t=x\big]$.
\end{teorema}

\begin{proof}
Sea $R_s=e^{-\int_t^sr(X)}$. Por Itô, $\dd(R_su(s,X_s))=R_s(\partial_su+\mathcal Lu-ru)\,\dd s+
R_s\nabla u\,\sigma\,\dd W_s=R_s\nabla u\,\sigma\,\dd W_s$. Integrando de $t$ a $T-\varepsilon$, localizando, tomando esperanzas y
haciendo $\varepsilon\to0$ (convergencia dominada, usando el crecimiento polinomial y las cotas de
momentos) se concluye.
\end{proof}

\begin{intuicion}
Si $g$ no es regular, por ejemplo $g=\ind_{[K,\infty)}$ (una opción digital), ¿es $u$ regular para
$t<T$? Analíticamente, la respuesta depende de propiedades de regularización de
$\partial_t+\mathcal L$: elipticidad o, más generalmente, hipoelipticidad. Probabilísticamente, $u(t,x)=
\int g(y)p_{T-t}(x,y)\,\dd y$, y la pregunta es si $X_T$ tiene una \emph{densidad regular}. El cálculo de
Malliavin responde a la segunda pregunta directamente, sin pasar por las EDPs, y \emph{además}
permite derivar $u$ en $x$ sin derivar $g$: esto son las fórmulas de griegas de los Capítulos~\ref{cap:griegas}
y siguientes.
\end{intuicion}

\section{Ejemplos fundamentales}

\begin{ejemplo}[Ornstein--Uhlenbeck]
$\dd X=-\lambda X\,\dd t+\sigma\,\dd W$ tiene solución $X_t=e^{-\lambda t}x+
\sigma\int_0^te^{-\lambda(t-s)}\,\dd W_s$, gaussiana con varianza $\frac{\sigma^2}{2\lambda}(1-e^{-2\lambda t})$. Con $\lambda=1$, $\sigma=\sqrt2$, su ley invariante
es $N(0,1)$ y su semigrupo es la \emph{fórmula de Mehler}
\[
  P_tf(x)=\E f\big(e^{-t}x+\sqrt{1-e^{-2t}}\,\xi\big),\qquad\xi\sim N(0,1),
\]
que reaparecerá en dimensión infinita en el Capítulo~\ref{cap:ou}.
\end{ejemplo}

\begin{ejemplo}[Black--Scholes]
$\dd S_t=rS_t\,\dd t+\sigma S_t\,\dd W_t$ (bajo la medida de riesgo neutro) tiene
solución $S_t=S_0\exp((r-\frac{\sigma^2}{2})t+\sigma W_t)$, y $Y_t=\partial S_t/\partial S_0=S_t/S_0$.
\end{ejemplo}


\part{El núcleo del cálculo de Malliavin}
\chapter{Procesos gaussianos isonormales}\label{cap:isonormales}

\begin{panorama}
Abstraemos el marco: en lugar de trabajar en el espacio de trayectorias $C_0([0,T])$,
trabajamos con una isometría $W:\Hf\to L^2(\Omega)$ de un espacio de Hilbert real separable
$\Hf$ en un espacio gaussiano. Este punto de vista, debido esencialmente a Gross, Itô y
Segal y sistematizado por Nualart~\cite{Nualart06}, cubre de una vez el browniano $d$-dimensional, el
browniano fraccionario, el ruido blanco espacio-temporal y los ruidos coloreados de las
EDPEs.
\end{panorama}

\section{Definición y ejemplos}

\begin{definicion}
Sea $\Hf$ un espacio de Hilbert real separable. Un \emph{proceso gaussiano isonormal}
sobre $\Hf$ es una familia $W=\{W(h):h\in\Hf\}$ de variables gaussianas centradas definidas en
$(\Omega,\F,\P)$ tal que
\[
  \E[W(h)W(g)]=\ip{h}{g}_\Hf\qquad\text{para todo }h,g\in\Hf.
\]
Supondremos siempre que $\F$ es la $\sigma$-álgebra generada por $W$ (completada).
\end{definicion}

\begin{proposicion}
$h\mapsto W(h)$ es lineal, y por tanto una isometría lineal de $\Hf$ sobre un espacio
gaussiano $\Hc_1\subset L^2(\Omega)$, \textup{el primer caos}. Dado $\Hf$, siempre existe un proceso isonormal.
\end{proposicion}

\begin{proof}
$\E[(W(ah+bg)-aW(h)-bW(g))^2]=0$ desarrollando con la isometría. Para
la existencia, toma una base ortonormal $\{e_n\}$, variables i.i.d.\ $\xi_n\sim N(0,1)$ y define $W(h)=
\sum_n\xi_n\ip{h}{e_n}$, que converge en $L^2$ y es gaussiana por la Proposición~\ref{prop:gaussianos}.
\end{proof}

\begin{ejemplo}[Browniano]
$\Hf=L^2([0,T];\R^m)$ y $W(h)=\sum_k\int_0^Th^k(s)\,\dd W^k_s$, la \emph{integral de
Wiener}. Entonces $W^k_t=W(\ind_{[0,t]}e_k)$.
\end{ejemplo}

\begin{ejemplo}[Ruido blanco espacio-temporal]
$\Hf=L^2(\R_+\times\R^d)$. Para un boreliano $A$ de
medida finita, $W(A)=W(\ind_A)$ es $N(0,\abs{A})$, y $W(A),W(B)$ son independientes si $A\cap B=\emptyset$.
Es el ruido de la ecuación del calor estocástica de Walsh (Capítulo~\ref{cap:integracion-edpe}).
\end{ejemplo}

\begin{ejemplo}[Ruido coloreado en espacio]
Sea $\gamma$ una medida no negativa definida (la covarianza
espacial) con transformada de Fourier $\mu=\hat\gamma$ (la \emph{medida espectral}). Se toma $\Hf$ como la compleción
de $C_c^\infty(\R_+\times\R^d)$ respecto a
\[
  \ip{\varphi}{\psi}_\Hf=\int_0^\infty\int_{\R^{2d}}\varphi(s,x)\psi(s,y)\gamma(x-y)\,\dd x\,\dd y\,\dd s=\int_0^\infty\int_{\R^d}\hat\varphi(s,\xi)\overline{\hat\psi(s,\xi)}\mu(\dd\xi)\,\dd s.
\]
Es el marco de Dalang para EDPEs en dimensión $d\ge2$ (Capítulo~\ref{cap:integracion-edpe}). Notemos que $\Hf$ puede
contener distribuciones, no solo funciones.
\end{ejemplo}

\begin{ejemplo}[Browniano fraccionario]
Para $H\in(0,1)$, $\Hf$ es la compleción de las funciones
escalonadas en $[0,T]$ respecto a $\ip{\ind_{[0,t]}}{\ind_{[0,s]}}_\Hf=R_H(s,t)=\frac12(t^{2H}+s^{2H}-\abs{t-s}^{2H})$. Ver
Capítulo~\ref{cap:fraccionario}.
\end{ejemplo}

\section{Variables suaves y densidad}

\begin{definicion}
Denotamos por $\Sc$ el conjunto de \emph{variables aleatorias suaves} (o cilíndricas)
\[
  F=f\big(W(h_1),\dots,W(h_n)\big),\qquad h_i\in\Hf,\ f\in C_p^\infty(\R^n),
\]
donde $C_p^\infty$ son las funciones $C^\infty$ con todas sus derivadas de crecimiento polinomial. Análogamente
$\Sc_b$ (con $f\in C_b^\infty$) y $\mathcal P$ (con $f$ polinomio).
\end{definicion}

\begin{lema}\label{lem:densidad-suaves}
Para todo $p\in[1,\infty)$, los conjuntos $\mathcal P$, $\Sc_b$ y el espacio generado por $\{e^{W(h)}:h\in\Hf\}$
son densos en $L^p(\Omega,\F,\P)$.
\end{lema}

\begin{proof}
Lo vemos para las exponenciales en $L^2$; el caso $L^p$ es análogo usando dualidad
$L^p$--$L^q$. Sea $G\in L^2$ ortogonal a todas las $e^{W(h)}$ y sea $\{e_n\}$ una base de $\Hf$. Para cada $n$, la
función $\lambda\mapsto\E[G\exp(\sum_{i\le n}\lambda_iW(e_i))]$ se anula en $\R^n$ y es entera, luego se anula para $\lambda\in i\R^n$:
la medida $\E[G\mid W(e_1),\dots,W(e_n)]\,\dd\P$ tiene transformada de Fourier nula y $\cond{G}{\G_n}=0$,
con $\G_n=\sigma(W(e_1),\dots,W(e_n))$. Como $\G_n\uparrow\F$, por convergencia de martingalas, $G=0$. Los
polinomios aproximan las exponenciales en $L^2$ (serie de Taylor, con convergencia dominada por
$e^{\abs{W(h)}}$), luego también son densos.
\end{proof}

\section{El espacio de Wiener abstracto}
El marco de proceso isonormal es puramente $L^2$. A veces conviene una realización topológica.

\begin{definicion}
Un \emph{espacio de Wiener abstracto} es una terna $(B,H,\mu)$ donde $B$ es un espacio
de Banach separable, $H\hookrightarrow B$ un espacio de Hilbert inyectado continua y densamente, y $\mu$
una medida gaussiana centrada en $B$ con $\int_Be^{i\ell(x)}\mu(\dd x)=e^{-\frac12\norm{\ell}_H^2}$ para $\ell\in B^*\subset H^*\cong H$.
\end{definicion}

El ejemplo canónico es $B=C_0([0,T])$, $H$ el espacio de Cameron--Martin, $\mu$ la medida de
Wiener. Para $\ell\in B^*$, $W(\ell)=\ell(\omega)$ se extiende por isometría a $H$, dando un proceso isonormal.
El Teorema de Cameron--Martin vale en este marco: $\mu(\cdot-h)\ll\mu$ si y solo si $h\in H$.

\begin{intuicion}
¿Por qué no usar directamente $\Hf$ como espacio muestral? Porque una medida gaussiana
estándar no existe en un Hilbert de dimensión infinita: si $\xi=\sum\xi_ne_n$ con $\xi_n$ i.i.d., entonces
$\norm{\xi}^2=\sum\xi_n^2=\infty$ c.s. El ruido ``vive'' en un espacio más grande $B$, mientras que $H$ es la
parte de $B$ que describe las direcciones de traslación admisibles. Por eso la derivada de
Malliavin, aunque se calcula ``en el punto $\omega\in B$'', toma valores en $H$.
\end{intuicion}

\chapter{Caos de Wiener}\label{cap:caos}

\begin{panorama}
Descomponemos $L^2(\Omega)$ en suma ortogonal de \emph{caos}: $L^2(\Omega)=\bigoplus_{n\ge0}\Hc_n$. El caos $n$-ésimo
está formado por los polinomios de Hermite de grado $n$ en variables gaussianas, o
equivalentemente por las integrales múltiples de Itô $I_n(f)$. Es el análogo gaussiano del
desarrollo en serie de Fourier, y todos los operadores de Malliavin ($D$, $\delta$, $L$) actúan de
forma diagonal o casi diagonal en esta descomposición.
\end{panorama}

\section{Polinomios de Hermite}

\begin{definicion}
Los \emph{polinomios de Hermite} (probabilísticos) $H_n$ se definen por la función
generatriz
\[
  e^{tx-\frac{t^2}{2}}=\sum_{n=0}^\infty\frac{t^n}{n!}H_n(x).
\]
Así $H_0=1$, $H_1=x$, $H_2=x^2-1$, $H_3=x^3-3x$, $H_4=x^4-6x^2+3$.
\end{definicion}

\begin{proposicion}\label{prop:hermite}
\begin{enumerate}[label=(\roman*)]
  \item $H_n(x)=(-1)^ne^{x^2/2}\frac{\dd^n}{\dd x^n}e^{-x^2/2}$.
  \item $H_n'=nH_{n-1}$ y $H_{n+1}(x)=xH_n(x)-H_n'(x)$.
  \item Si $(X,Y)$ es gaussiano centrado con $\E X^2=\E Y^2=1$, entonces $\E[H_n(X)H_m(Y)]=
  \delta_{nm}\,n!\,(\E XY)^n$.
  \item $\{H_n/\sqrt{n!}\}$ es una base ortonormal de $L^2(\R,\gamma)$, con $\gamma=N(0,1)$.
\end{enumerate}
\end{proposicion}

\begin{proof}
(i) $e^{tx-t^2/2}=e^{x^2/2}e^{-(x-t)^2/2}$, y el coeficiente de Taylor en $t=0$ es
$e^{x^2/2}\frac{\dd^n}{\dd t^n}e^{-(x-t)^2/2}|_{t=0}=(-1)^ne^{x^2/2}\frac{\dd^n}{\dd x^n}e^{-x^2/2}$. (ii) Derivando la función generatriz en $x$ y
en $t$. (iii) $\E[e^{sX-s^2/2}e^{tY-t^2/2}]=e^{st\E XY}$; se identifican coeficientes de $s^nt^m/(n!m!)$. (iv) La
ortonormalidad es (iii) con $X=Y$; la completitud es el Lema~\ref{lem:densidad-suaves} en dimensión 1, ya que los
$H_n$ generan los polinomios.
\end{proof}

\section{Descomposición en caos}

\begin{definicion}
El \emph{$n$-ésimo caos de Wiener} $\Hc_n$ es el subespacio cerrado de $L^2(\Omega)$ generado
por $\{H_n(W(h)):\ h\in\Hf,\ \norm{h}=1\}$. $\Hc_0$ son las constantes.
\end{definicion}

\begin{teorema}[Wiener, Itô]\label{teo:caos}
$L^2(\Omega,\F,\P)=\bigoplus_{n=0}^\infty\Hc_n$ \textup{(}suma ortogonal\textup{)}.
\end{teorema}

\begin{proof}
La ortogonalidad $\Hc_n\perp\Hc_m$ para $n\ne m$ es la Proposición~\ref{prop:hermite}(iii).
Para la completitud: si $G\perp\Hc_n$ para todo $n$, entonces, para $\norm{h}=1$, $\E[Ge^{tW(h)}]=
e^{t^2/2}\sum_n\frac{t^n}{n!}\E[GH_n(W(h))]=0$ (el intercambio de suma y esperanza se justifica porque
$\sum_n\frac{\abs{t}^n}{n!}\norm{H_n(W(h))}_2=\sum_n\abs{t}^n/\sqrt{n!}<\infty$). Por el Lema~\ref{lem:densidad-suaves}, $G=0$.
\end{proof}

Denotamos por $J_n$ la proyección ortogonal sobre $\Hc_n$, de modo que $F=\sum_nJ_nF$ para
$F\in L^2$.

\section{Integrales múltiples}
Sea $\Hf^{\otimes n}$ el producto tensorial hilbertiano y $\Hf^{\odot n}$ su subespacio simétrico, con la simetrización
$\tilde f$. Cuando $\Hf=L^2(T,\mu)$ con $\mu$ sin átomos, $\Hf^{\odot n}=L^2_s(T^n)$ son las funciones simétricas.

\begin{proposicion}\label{prop:integrales-multiples}
Existe una única aplicación lineal $I_n:\Hf^{\odot n}\to\Hc_n$ tal que $I_n(h^{\otimes n})=
H_n(W(h))$ para $\norm{h}=1$, y satisface
\[
  \E[I_n(f)I_m(g)]=\delta_{nm}\,n!\,\ip{f}{g}_{\Hf^{\otimes n}}.
\]
Así $\sqrt{n!}^{-1}I_n$ es una isometría de $\Hf^{\odot n}$ sobre $\Hc_n$. Extendemos $I_n(f):=I_n(\tilde f)$ a $f\in\Hf^{\otimes n}$.
\end{proposicion}

\begin{proof}
Por la Proposición~\ref{prop:hermite}(iii), $\E[H_n(W(h))H_n(W(g))]=n!\ip{h}{g}^n=n!\ip{h^{\otimes n}}{g^{\otimes n}}$.
Los tensores $h^{\otimes n}$ generan un subespacio denso de $\Hf^{\odot n}$ (por polarización), y la aplicación es
isométrica salvo el factor $n!$ en ese generador, luego se extiende. Es sobreyectiva sobre $\Hc_n$ por
definición de este.
\end{proof}

\begin{teorema}[Representación como integrales iteradas]\label{teo:iteradas}
Si $\Hf=L^2([0,T])$ y $W$ es el
browniano, para $f\in L^2_s([0,T]^n)$,
\[
  I_n(f)=n!\int_0^T\int_0^{t_n}\cdots\int_0^{t_2}f(t_1,\dots,t_n)\,\dd W_{t_1}\cdots\dd W_{t_n}.
\]
\end{teorema}

\begin{proof}
Ambos lados son isometrías (salvo el factor $n!$) de $L^2_s$ en $L^2(\Omega)$, por aplicación
iterada de la isometría de Itô, y coinciden en $f=h^{\otimes n}$: denotando $M_t=\int_0^th\,\dd W$ y $\langle M\rangle_t=\int_0^th^2$,
el proceso $\langle M\rangle_t^{n/2}H_n(M_t/\langle M\rangle_t^{1/2})$ satisface, por Itô y la identidad $H_n'=nH_{n-1}$, la relación
$\dd(\,\cdot\,)_n=n(\,\cdot\,)_{n-1}h_t\,\dd W_t$, que es exactamente la recursión de las integrales iteradas.
\end{proof}

\begin{corolario}[Descomposición en caos de Itô]\label{cor:caos-ito}
Toda $F\in L^2(\Omega,\F,\P)$ se escribe de forma
única como $F=\sum_{n=0}^\infty I_n(f_n)$, con $f_n\in\Hf^{\odot n}$, $f_0=\E F$, y
\[
  \E F^2=\sum_{n=0}^\infty n!\norm{f_n}^2_{\Hf^{\otimes n}}.
\]
\end{corolario}

\begin{ejemplo}\label{ej:exponencial-caos}
$e^{W(h)-\frac12\norm{h}^2}=\sum_n\frac1{n!}I_n(h^{\otimes n})$. En efecto, para $\norm{h}=1$ es la función generatriz
de Hermite con $t=\norm{h}$.
\end{ejemplo}

\section{Contracciones y fórmula del producto}
Para $f\in\Hf^{\odot p}$, $g\in\Hf^{\odot q}$ y $r\le p\wedge q$, la \emph{contracción de orden $r$} $f\otimes_rg\in\Hf^{\otimes(p+q-2r)}$ se define, si
$\{e_i\}$ es una base ortonormal, por
\[
  f\otimes_rg=\sum_{i_1,\dots,i_r}\ip{f}{e_{i_1}\otimes\dots\otimes e_{i_r}}\otimes\ip{g}{e_{i_1}\otimes\dots\otimes e_{i_r}}.
\]
En $L^2(T)$: $(f\otimes_rg)(t,s)=\int_{T^r}f(t,u)g(s,u)\mu^r(\dd u)$, con $t\in T^{p-r}$, $s\in T^{q-r}$.

\begin{teorema}[Fórmula del producto]\label{teo:producto}
Para $f\in\Hf^{\odot p}$, $g\in\Hf^{\odot q}$,
\[
  I_p(f)I_q(g)=\sum_{r=0}^{p\wedge q}r!\binom pr\binom qrI_{p+q-2r}(f\otimes_rg).
\]
\end{teorema}

\begin{proof}[Esbozo]
Por bilinealidad y densidad basta tomar $f=h^{\otimes p}$, $g=k^{\otimes q}$. Se compara el coeficiente
de $s^pt^q$ en
\[
  e^{sW(h)-\frac{s^2}{2}\norm{h}^2}e^{tW(k)-\frac{t^2}{2}\norm{k}^2}=e^{W(sh+tk)-\frac12\norm{sh+tk}^2}e^{st\ip{h}{k}},
\]
expandiendo cada exponencial estocástica en caos (ejemplo anterior) y $e^{st\ip{h}{k}}=\sum_r\frac{(st)^r}{r!}\ip{h}{k}^r$.
La contracción $h^{\otimes p}\otimes_rk^{\otimes q}=\ip{h}{k}^rh^{\otimes(p-r)}\otimes k^{\otimes(q-r)}$, lo que produce los coeficientes combinatorios.
\end{proof}

\begin{ejemplo}
$I_1(h)^2=I_2(h\otimes h)+\norm{h}^2$, coherente con $H_2(x)=x^2-1$.
\end{ejemplo}

\begin{intuicion}
Las contracciones miden ``cuánta estructura comparten'' $f$ y $g$. Veremos en el Capítulo~\ref{cap:malliavin-stein}
que una sucesión $F_n=I_q(f_n)$ con varianza 1 converge a $N(0,1)$ si y solo si $\norm{f_n\otimes_rf_n}\to0$
para $r=1,\dots,q-1$: las contracciones controlan la gaussianidad. Es el \emph{teorema del
cuarto momento}.
\end{intuicion}

\section{Equivalencia de normas dentro de un caos}

\begin{teorema}\label{teo:normas-caos}
Para cada $n$ y $p\ge2$, en el caos $\Hc_n$ todas las normas $L^p$ son equivalentes:
\[
  \norm{F}_p\le(p-1)^{n/2}\norm{F}_2,\qquad F\in\Hc_n.
\]
\end{teorema}

Esto es consecuencia inmediata de la hipercontractividad de Nelson (Teorema~\ref{teo:hipercontractividad}), que
probaremos en el Capítulo~\ref{cap:ou}. Implica que la convergencia en $L^2$ dentro de un caos fijo arrastra
la convergencia de todos los momentos, hecho crucial para los teoremas límite.

\chapter{La derivada de Malliavin}\label{cap:derivada}

\begin{panorama}
Definimos la derivada $D$ en las variables suaves, probamos la fórmula de integración por
partes fundamental, la cerrabilidad de $D$ y construimos los espacios de Sobolev--Watanabe
$\D^{k,p}$. Establecemos la regla de la cadena, la expresión de $D$ en el caos de Wiener (incluida
la fórmula de Stroock) y calculamos la derivada de funcionales típicos: integrales de
Wiener, el supremo del browniano y esperanzas condicionales.
\end{panorama}

Trabajamos con un proceso isonormal $W$ sobre $\Hf$, y $\F=\sigma(W)$.

\section{Definición en las variables suaves}

\begin{definicion}
Para $F=f(W(h_1),\dots,W(h_n))\in\Sc$, la \emph{derivada de Malliavin} es la variable
con valores en $\Hf$
\[
  DF=\sum_{i=1}^n\partial_if\big(W(h_1),\dots,W(h_n)\big)\,h_i.
\]
\end{definicion}

Es fácil ver que $DF$ no depende de la representación de $F$ (si $F$ admite dos representaciones,
se reducen a una común con $h_i$ linealmente independientes). En el espacio de Wiener clásico,
$\Hf=L^2([0,T])$, $DF$ es un proceso $(D_tF)_{t\in[0,T]}$ y
\[
  \ip{DF}{\dot h}_{L^2}=\int_0^TD_tF\,\dot h(t)\,\dd t=\frac{\dd}{\dd\varepsilon}F(\omega+\varepsilon h)\Big|_{\varepsilon=0},\qquad h=\textstyle\int_0^\cdot\dot h,
\]
pues $W(g)(\omega+\varepsilon h)=W(g)(\omega)+\varepsilon\ip{g}{\dot h}$. Es la \emph{derivada de Fréchet en las direcciones de
Cameron--Martin}.

\begin{ejemplo}
\begin{enumerate}[label=(\alph*)]
  \item $DW(h)=h$. En particular $D_tW_s=\ind_{[0,s]}(t)$.
  \item $D_t(W_s^2)=2W_s\ind_{[0,s]}(t)$; $D_t\int_0^TW_s\,\dd s=T-t$.
  \item $D\,e^{W(h)}=e^{W(h)}h$.
\end{enumerate}
\end{ejemplo}

\begin{intuicion}
$D_tF$ mide cuánto cambia $F$ cuando se perturba el ruido \emph{en el instante $t$}: es la respuesta
de $F$ a un ``impulso'' $\delta_t$ en $\dot W$. Por ejemplo, $W_s$ responde con 1 a impulsos anteriores a $s$
y con 0 a impulsos posteriores. Si $F$ es la solución de una EDE en el instante $T$, $D_tX_T$
es la propagación de una perturbación desde $t$ hasta $T$ (Capítulo~\ref{cap:hormander}).
\end{intuicion}

\section{Integración por partes y cerrabilidad}

\begin{lema}[Integración por partes fundamental]\label{lem:ibp}
Para $F\in\Sc$ y $h\in\Hf$,
\[
  \E\ip{DF}{h}_\Hf=\E\big[F\,W(h)\big].
\]
\end{lema}

\begin{proof}
Por homogeneidad podemos suponer $\norm{h}=1$, y podemos escribir $F=
f(W(e_1),\dots,W(e_n))$ con $\{e_i\}$ ortonormales y $e_1=h$. Entonces $(W(e_1),\dots,W(e_n))\sim N(0,I_n)$
y, por el Lema~\ref{lem:stein},
\[
  \E\ip{DF}{h}=\E[\partial_1f(W(e))]=\int_{\R^n}\partial_1f(x)\phi_n(x)\,\dd x=\int_{\R^n}f(x)x_1\phi_n(x)\,\dd x=\E[FW(e_1)].
  \qedhere
\]
\end{proof}

\begin{lema}\label{lem:ibp2}
Para $F,G\in\Sc$ y $h\in\Hf$, $\E\big[G\ip{DF}{h}\big]=\E\big[-F\ip{DG}{h}+FGW(h)\big]$.
\end{lema}

\begin{proof}
Aplicar el Lema~\ref{lem:ibp} a $FG$ y usar la regla de Leibniz $D(FG)=F\,DG+G\,DF$,
obvia en $\Sc$.
\end{proof}

\begin{proposicion}[Cerrabilidad]\label{prop:cerrabilidad}
Para $p\ge1$, el operador $D:\Sc\subset L^p(\Omega)\to L^p(\Omega;\Hf)$ es
cerrable.
\end{proposicion}

\begin{proof}
Sea $F_n\in\Sc$ con $F_n\to0$ en $L^p$ y $DF_n\to\eta$ en $L^p(\Omega;\Hf)$. Hay que ver $\eta=0$. Sea
$G\in\Sc$ tal que $GW(h)$ esté acotada (por ejemplo $G=g\,e^{-\varepsilon W(h)^2}$ con $g\in\Sc_b$). Por el Lema~\ref{lem:ibp2},
\[
  \E\big[G\ip{\eta}{h}\big]=\lim_n\E\big[G\ip{DF_n}{h}\big]=\lim_n\E\big[-F_n\ip{DG}{h}+F_nGW(h)\big]=0,
\]
porque $\ip{DG}{h}$ y $GW(h)$ están acotadas. Tales $G$ son densas en $L^q$ (tomando $\varepsilon\to0$), luego
$\ip{\eta}{h}=0$ c.s.\ para todo $h$ de un subconjunto numerable denso, y $\eta=0$.
\end{proof}

\section{Espacios de Sobolev--Watanabe}

\begin{definicion}
Para $p\ge1$, $\D^{1,p}$ es el dominio de la clausura de $D$ en $L^p$, es decir, la
compleción de $\Sc$ respecto a
\[
  \norm{F}_{1,p}=\big(\E\abs{F}^p+\E\norm{DF}_\Hf^p\big)^{1/p}.
\]
Iterando, $D^kF$ es una variable con valores en $\Hf^{\otimes k}$, y $\D^{k,p}$ es la compleción con $\norm{F}_{k,p}^p=
\E\abs{F}^p+\sum_{j=1}^k\E\norm{D^jF}^p_{\Hf^{\otimes j}}$. Se definen análogamente $\D^{k,p}(V)$ para $V$-valuadas, con $V$ un
Hilbert, y
\[
  \D^\infty=\bigcap_{k,p}\D^{k,p}.
\]
\end{definicion}

$\D^{1,2}$ es un espacio de Hilbert. Se tiene $\D^{k,p}\subset\D^{k',p'}$ si $k\ge k'$, $p\ge p'$. Los elementos de $\D^\infty$
son los ``funcionales suaves de Wiener''. Veremos que, gracias a las desigualdades de Meyer, $\D^{k,p}$
coincide con el dominio de $(I-L)^{k/2}$ en $L^p$.

\section{Regla de la cadena}

\begin{proposicion}[Regla de la cadena]\label{prop:cadena}
Sea $\varphi\in C^1(\R^m)$ con derivadas parciales acotadas y
$F=(F^1,\dots,F^m)$ con $F^i\in\D^{1,p}$. Entonces $\varphi(F)\in\D^{1,p}$ y
\[
  D\varphi(F)=\sum_{i=1}^m\partial_i\varphi(F)\,DF^i.
\]
Si solo se supone $\varphi$ globalmente Lipschitz y $p>1$, entonces $\varphi(F)\in\D^{1,p}$ y existe un vector aleatorio
$G$ con $\abs{G}\le\norm{\varphi}_{\mathrm{Lip}}$ tal que $D\varphi(F)=\sum_iG_iDF^i$. Si además la ley de $F$ es absolutamente
continua, $G=\nabla\varphi(F)$.
\end{proposicion}

\begin{proof}
Para $F^i\in\Sc$ y $\varphi\in C_b^1$ es la regla de la cadena clásica. En general se aproxima
$F^i$ por $F^i_n\in\Sc$ en $\D^{1,p}$; entonces $\varphi(F_n)\to\varphi(F)$ en $L^p$ y $\sum\partial_i\varphi(F_n)DF^i_n\to\sum\partial_i\varphi(F)DF^i$ en
$L^p(\Hf)$ (convergencia dominada sobre una subsucesión con convergencia c.s.); por cerrabilidad,
se concluye. Para $\varphi$ Lipschitz, se regulariza $\varphi_\varepsilon=\varphi*\rho_\varepsilon$; la sucesión $D\varphi_\varepsilon(F)$ está acotada en
$L^p(\Hf)$, que es reflexivo para $p>1$, y un límite débil identifica la derivada.
\end{proof}

\begin{ejemplo}
Si $F\in\D^{1,2}$, entonces $\abs{F}\in\D^{1,2}$ y $D\abs{F}=\sgn(F)DF$; $F^+\in\D^{1,2}$ y $DF^+=
\ind_{\{F>0\}}DF$.
\end{ejemplo}

\begin{ejemplo}[Una variable no derivable]\label{ej:no-derivable}
$F=\ind_{\{W(h)>0\}}\notin\D^{1,2}$ si $h\ne0$. Si lo fuera, por
localidad (Proposición~\ref{prop:localidad}), $DF=0$ c.s.\ en $\{F=0\}\cup\{F=1\}=\Omega$, y por el Corolario~\ref{cor:constante},
$F$ sería constante. Los indicadores de sucesos no triviales nunca son Malliavin-diferenciables.
Esta es la razón de fondo de por qué las griegas de opciones digitales requieren integración por
partes (Capítulo~\ref{cap:griegas}).
\end{ejemplo}

\section{La derivada en el caos de Wiener}

\begin{teorema}\label{teo:derivada-caos}
Sea $F=\sum_nI_n(f_n)\in L^2(\Omega)$. Entonces $F\in\D^{1,2}$ si y solo si
\[
  \sum_{n=1}^\infty n\,n!\,\norm{f_n}^2_{\Hf^{\otimes n}}<\infty,
\]
y en tal caso, en el caso $\Hf=L^2(T,\mu)$,
\[
  D_tF=\sum_{n=1}^\infty nI_{n-1}\big(f_n(\cdot,t)\big),\qquad\E\norm{DF}^2_\Hf=\sum_{n\ge1}n\,n!\,\norm{f_n}^2.
\]
\end{teorema}

\begin{proof}
Para $F=I_n(h^{\otimes n})=H_n(W(h))$ con $\norm{h}=1$: $DF=H_n'(W(h))h=
nH_{n-1}(W(h))h=nI_{n-1}(h^{\otimes(n-1)})h$, que es la fórmula. Por linealidad y densidad vale para todo
$f_n\in\Hf^{\odot n}$, y $\E\norm{DI_n(f_n)}^2=n^2(n-1)!\norm{f_n}^2=n\,n!\norm{f_n}^2$. Además, $DI_n(f_n)$ y $DI_m(f_m)$ son
ortogonales en $L^2(\Omega;\Hf)$ si $n\ne m$. Si la serie converge, las sumas parciales $F_N=\sum_{n\le N}I_n(f_n)$
convergen en $\D^{1,2}$ y, por cerrabilidad, $F\in\D^{1,2}$. Recíprocamente, si $F\in\D^{1,2}$, se usa que $D$
conmuta con las proyecciones $J_n$ (salvo desplazamiento de índice), lo que se comprueba en $\Sc$, y
la serie es $\E\norm{DF}^2<\infty$.
\end{proof}

\begin{corolario}[Fórmula de Stroock]\label{cor:stroock}
Si $F\in\D^{n,2}$, el núcleo del $n$-ésimo caos es $f_n=
\frac1{n!}\E[D^nF]$.
\end{corolario}

\begin{proof}
Aplicando $n$ veces el Teorema~\ref{teo:derivada-caos}, $D^nF=\sum_{m\ge n}\frac{m!}{(m-n)!}I_{m-n}(f_m)$, cuya esperanza es el término $m=n$: $n!f_n$.
\end{proof}

\begin{ejemplo}
$F=W_1^3$. $\E F=0$, $\E DF=\E[3W_1^2]\ind_{[0,1]}=3\cdot\ind_{[0,1]}$, $\E D^2F=0$, $\E D^3F=6\cdot\ind_{[0,1]^3}$.
Así $W_1^3=I_1(3\cdot\ind)+I_3(\ind^{\otimes3})=3W_1+H_3(W_1)$. Correcto: $H_3(x)=x^3-3x$.
\end{ejemplo}

\section{Propiedades estructurales}

\begin{corolario}\label{cor:constante}
Si $F\in\D^{1,2}$ y $DF=0$, entonces $F=\E F$ es constante.
\end{corolario}

\begin{proof}
$0=\E\norm{DF}^2=\sum n\,n!\norm{f_n}^2$ obliga a $f_n=0$ para $n\ge1$.
\end{proof}

\begin{proposicion}[Localidad]\label{prop:localidad}
Si $F\in\D^{1,1}$, entonces $DF=0$ c.s.\ en el conjunto $\{F=0\}$.
Más en general, si $F=G$ en un suceso $A$, entonces $DF=DG$ c.s.\ en $A$.
\end{proposicion}

\begin{proof}[Esbozo]
Sea $\psi_\varepsilon(x)=\int_0^x\ind_{[-\varepsilon,\varepsilon]}(y)\,\dd y$, Lipschitz con $\psi_\varepsilon'=\ind_{[-\varepsilon,\varepsilon]}$ c.t.p.\ y $\abs{\psi_\varepsilon}\le\varepsilon$. Por la regla
de la cadena (versión Lipschitz), $D\psi_\varepsilon(F)=\ind_{\{\abs{F}\le\varepsilon\}}DF$. Para $u$ suave con valores en $\Hf$, la
dualidad del Capítulo~\ref{cap:divergencia} da
\[
  \Big|\E\big[\ind_{\{\abs{F}\le\varepsilon\}}\ip{DF}{u}\big]\Big|=\abs{\E[\psi_\varepsilon(F)\delta(u)]}\le\varepsilon\,\E\abs{\delta(u)}\xrightarrow[\varepsilon\to0]{}0.
\]
Por convergencia dominada, $\E[\ind_{\{F=0\}}\ip{DF}{u}]=0$ para todo $u$ de un conjunto denso, de donde
$\ind_{\{F=0\}}DF=0$.
\end{proof}

\begin{proposicion}[Derivada de una esperanza condicional]\label{prop:derivada-esp-cond}
Sea $\Hf=L^2(T,\mu)$, $A\subset T$ medible
y $\F_A=\sigma\{W(h):\ \operatorname{supp}h\subset A\}$. Si $F\in\D^{1,2}$, entonces $\cond{F}{\F_A}\in\D^{1,2}$ y
\[
  D_t\cond{F}{\F_A}=\cond{D_tF}{\F_A}\ind_A(t).
\]
En particular, con $A=[0,s]$ en el espacio de Wiener: si $F$ es $\F_s$-medible, $D_tF=0$ para $t>s$.
\end{proposicion}

\begin{proof}
Si $F=I_n(f_n)$, entonces $\cond{F}{\F_A}=I_n(f_n\ind_A^{\otimes n})$ (por Proposición~\ref{prop:hermite}(iii), basta
comprobarlo para $f_n=h^{\otimes n}$, donde se descompone $h=h\ind_A+h\ind_{A^c}$, cuyas integrales son
independientes). Aplicando el Teorema~\ref{teo:derivada-caos},
\[
  D_t\cond{F}{\F_A}=nI_{n-1}\big(f_n(\cdot,t)\ind_A^{\otimes(n-1)}\big)\ind_A(t)=\cond{D_tF}{\F_A}\ind_A(t).
\]
Se concluye sumando sobre $n$.
\end{proof}

\section{Un ejemplo no trivial: el supremo del browniano}

\begin{proposicion}\label{prop:supremo}
Sea $M=\sup_{t\in[0,1]}W_t$. Entonces $M\in\D^{1,2}$ y $D_sM=\ind_{[0,\tau]}(s)$, donde $\tau$ es
el único instante \textup{(c.s.)} en que se alcanza el máximo.
\end{proposicion}

\begin{proof}
Sea $\{t_i\}$ denso en $[0,1]$ y $M_n=\max_{i\le n}W_{t_i}$. La aplicación $\varphi_n(x)=\max_{i\le n}x_i$
es Lipschitz, y por la Proposición~\ref{prop:cadena}, $D_sM_n=\ind_{[0,\tau_n]}(s)$ con $\tau_n$ el $t_i$ donde se alcanza el máximo
(único c.s.\ por continuidad absoluta). Así $\norm{DM_n}^2_\Hf=\tau_n\le1$, y $M_n\to M$ en $L^2$. Como $\{M_n\}$
está acotada en $\D^{1,2}$, una subsucesión converge débilmente y, por cerrabilidad, $M\in\D^{1,2}$. La
unicidad del maximizante (consecuencia de la propiedad de Markov) da $\tau_n\to\tau$ c.s., e identifica
$DM=\ind_{[0,\tau]}$.
\end{proof}

\begin{intuicion}
Perturbar el ruido antes de $\tau$ desplaza toda la trayectoria posterior, incluido el máximo;
perturbarlo después de $\tau$ no cambia el máximo (a primer orden). El supremo no es
diferenciable en sentido Fréchet en $C_0$, pero sí en el sentido de Malliavin: la derivada solo
se pide en direcciones de Cameron--Martin y en $L^2$.
\end{intuicion}

\section{Caracterización por trayectorias}
En el espacio de Wiener clásico, $\D^{1,p}$ tiene una descripción intrínseca debida a Kusuoka y
Sugita: $F\in\D^{1,p}$, $p>1$, si y solo si (i) para cada $h\in H$ existe una versión de $F$ tal
que $\varepsilon\mapsto F(\omega+\varepsilon h)$ es absolutamente continua (\emph{ray absolute continuity}); y (ii) existe $DF\in
L^p(\Omega;H)$ tal que $\frac1\varepsilon(F(\omega+\varepsilon h)-F(\omega))\to\ip{DF}{h}$ en probabilidad para todo $h$ (\emph{diferenciabilidad
estocástica de Gâteaux}). En particular, todo funcional $F$ que sea $H$-Lipschitz, es decir, con
$\abs{F(\omega+h)-F(\omega)}\le C\norm{h}_H$, está en $\D^{1,p}$ con $\norm{DF}_H\le C$.

\chapter{La divergencia y la fórmula de Clark--Ocone}\label{cap:divergencia}

\begin{panorama}
El operador divergencia $\delta$ es el adjunto de $D$. En el espacio de Wiener, $\delta$ extiende la
integral de Itô a integrandos \emph{anticipativos}: es la \emph{integral de Skorokhod}. Probamos las
reglas de cálculo de $\delta$, su expresión en caos, la relación de conmutación de Heisenberg
$[D,\delta]$, y la fórmula de Clark--Ocone, que identifica explícitamente el integrando de la
representación de martingalas.
\end{panorama}

\section{Definición como adjunto}

\begin{definicion}
El \emph{operador divergencia} $\delta$ es el adjunto de $D:\D^{1,2}\subset L^2(\Omega)\to L^2(\Omega;\Hf)$.
Su dominio $\Dom\delta$ es el conjunto de $u\in L^2(\Omega;\Hf)$ tales que
\[
  \abs{\E\ip{DF}{u}_\Hf}\le c_u\norm{F}_{L^2}\qquad\text{para toda }F\in\Sc,
\]
y para tales $u$, $\delta(u)\in L^2(\Omega)$ es la única variable que satisface la \emph{fórmula de dualidad}
\begin{equation}\label{eq:dualidad}
  \E\big[F\,\delta(u)\big]=\E\ip{DF}{u}_\Hf\qquad\text{para toda }F\in\D^{1,2}.
\end{equation}
\end{definicion}

Tomando $F=1$: $\E\delta(u)=0$. Como $D$ es cerrado y densamente definido, $\delta$ es cerrado y
densamente definido.

\begin{proposicion}\label{prop:reglas-delta}
\begin{enumerate}[label=(\roman*)]
  \item Si $u=\sum_{j=1}^nF_jh_j$ con $F_j\in\Sc$ y $h_j\in\Hf$, entonces $u\in\Dom\delta$ y
  \[
    \delta(u)=\sum_jF_jW(h_j)-\sum_j\ip{DF_j}{h_j}.
  \]
  \item Si $F\in\D^{1,2}$, $u\in\Dom\delta$, $Fu\in L^2(\Omega;\Hf)$ y $F\delta(u)-\ip{DF}{u}\in L^2$, entonces $Fu\in\Dom\delta$ y
  \[
    \delta(Fu)=F\delta(u)-\ip{DF}{u}_\Hf.
  \]
\end{enumerate}
\end{proposicion}

\begin{proof}
(i) Por el Lema~\ref{lem:ibp2}, para $G\in\Sc$, $\E\ip{DG}{F_jh_j}=\E[G(F_jW(h_j)-\ip{DF_j}{h_j})]$.
(ii) Para $G\in\Sc$, $\E\ip{DG}{Fu}=\E\ip{D(FG)-G\,DF}{u}=\E[G(F\delta(u)-\ip{DF}{u})]$.
\end{proof}

\begin{intuicion}
La fórmula (ii) es un ``producto de Itô--Skorokhod'': sacar $F$ de la integral cuesta el
término de corrección $\ip{DF}{u}$. Si $F$ ``no ve'' el ruido que integra $u$ (por ejemplo, si $F$
es $\F_s$-medible y $u$ está soportado en $(s,T]$), entonces $\ip{DF}{u}=0$ y $F$ sale como una
constante: esa es exactamente la propiedad que hace funcionar la integral de Itô.
\end{intuicion}

\section{La divergencia en el caos}

\begin{proposicion}\label{prop:delta-caos}
Sea $\Hf=L^2(T,\mu)$ y $u\in L^2(\Omega;\Hf)$ con desarrollo $u_t=\sum_{n\ge0}I_n(f_n(\cdot,t))$,
$f_n\in\Hf^{\odot n}\otimes\Hf$. Entonces $u\in\Dom\delta$ si y solo si $\sum_n(n+1)!\,\norm{\tilde f_n}^2_{\Hf^{\otimes(n+1)}}<\infty$, y en tal caso
\[
  \delta(u)=\sum_{n=0}^\infty I_{n+1}(\tilde f_n),
\]
donde $\tilde f_n$ es la simetrización en las $n+1$ variables.
\end{proposicion}

\begin{proof}
Sea $G=I_{n+1}(g)$, $g$ simétrica. Por el Teorema~\ref{teo:derivada-caos} y la isometría,
\begin{align*}
  \E\ip{DG}{u}&=\int_T\E\big[(n+1)I_n(g(\cdot,t))I_n(f_n(\cdot,t))\big]\mu(\dd t)\\
  &=(n+1)!\ip{g}{f_n}_{\Hf^{\otimes(n+1)}}=(n+1)!\ip{g}{\tilde f_n}=\E[GI_{n+1}(\tilde f_n)].
\end{align*}
Para $G$ en otros caos ambos lados son cero. La condición de dominio es la acotación de esta
forma lineal.
\end{proof}

\section{La integral de Skorokhod extiende la de Itô}

\begin{teorema}\label{teo:skorokhod-ito}
En el espacio de Wiener ($\Hf=L^2([0,T])$), todo proceso $u\in L^2_{\mathrm{prog}}$ pertenece
a $\Dom\delta$ y
\[
  \delta(u)=\int_0^Tu_t\,\dd W_t.
\]
\end{teorema}

\begin{proof}
Sea $u=F\ind_{(s,t]}$ con $F\in\Sc$ y $\F_s$-medible. Por la Proposición~\ref{prop:reglas-delta}(i), $\delta(u)=
F(W_t-W_s)-\int_s^tD_rF\,\dd r$. Por la Proposición~\ref{prop:derivada-esp-cond}, $D_rF=0$ para $r>s$, luego $\delta(u)=F(W_t-W_s)$,
que es la integral de Itô. Por linealidad vale para elementales. Por densidad (Lema~\ref{lem:densidad-elementales}, con
$F\in\Sc$ $\F_s$-medibles densas en $L^2(\F_s)$), dado $u\in L^2_{\mathrm{prog}}$, existen elementales $u_n\to u$ con
$\delta(u_n)=\int u_n\,\dd W\to\int u\,\dd W$ en $L^2$. Como $\delta$ es cerrado, $u\in\Dom\delta$ y $\delta(u)=\int u\,\dd W$.
\end{proof}

\begin{ejemplo}[Un integrando anticipativo]
$\delta(W_T\ind_{[0,T]})=W_TW(\ind_{[0,T]})-\int_0^TD_tW_T\,\dd t=W_T^2-T$.
Comparar con $\int_0^TW_t\,\dd W_t=\frac12(W_T^2-T)$.
\end{ejemplo}

\section{Relación de conmutación y una isometría}

\begin{proposicion}[Relación de Heisenberg]\label{prop:heisenberg}
Si $u\in\D^{1,2}(\Hf)$ y $D_tu\in\Dom\delta$ para casi todo $t$,
con $(\delta(D_tu))_t\in L^2(\Omega;\Hf)$, entonces $\delta(u)\in\D^{1,2}$ y
\[
  D_t\delta(u)=u_t+\delta(D_tu).
\]
\end{proposicion}

\begin{proof}
Para $u=Fh$ con $F\in\Sc$: $\delta(u)=FW(h)-\ip{DF}{h}$, luego
\[
  D_t\delta(u)=D_tF\,W(h)+Fh_t-\ip{D(D_tF)}{h}=Fh_t+\big(D_tF\,W(h)-\ip{D(D_tF)}{h}\big)=u_t+\delta(D_tu).
\]
El caso general se sigue por linealidad y cerrabilidad.
\end{proof}

En términos de operadores, $[D,\delta]=I$ en un sentido apropiado; es la relación canónica entre
los operadores de aniquilación ($D$) y creación ($\delta$) del espacio de Fock bosónico $\Gamma(\Hf)\cong L^2(\Omega)$,
que es exactamente la descomposición en caos.

\begin{teorema}\label{teo:isometria-skorokhod}
$\D^{1,2}(\Hf)\subset\Dom\delta$ y, para $u\in\D^{1,2}(\Hf)$,
\[
  \E[\delta(u)^2]=\E\norm{u}_\Hf^2+\E\int_T\int_TD_su_t\,D_tu_s\,\mu(\dd s)\mu(\dd t)\le\E\norm{u}_\Hf^2+\E\norm{Du}^2_{\Hf\otimes\Hf}.
\]
\end{teorema}

\begin{proof}
Para $u$ suave elemental, usando la dualidad y la relación de Heisenberg,
\[
  \E[\delta(u)^2]=\E\ip{D\delta(u)}{u}=\E\int u_t(u_t+\delta(D_tu))\,\dd t=\E\norm{u}^2+\int\E\ip{D(u_t)}{D_tu}\,\dd t,
\]
y el último término es $\E\iint D_su_tD_tu_s\,\dd s\,\dd t$. La desigualdad es Cauchy--Schwarz. Se extiende
por densidad, lo que también prueba la inclusión.
\end{proof}

Si $u$ es adaptado, $D_su_t=0$ para $s>t$, luego $D_su_tD_tu_s=0$ salvo en la diagonal: se recupera
la isometría de Itô.

\begin{corolario}[Derivada de una integral de Itô]\label{cor:derivada-ito}
Si $u\in L^2_{\mathrm{prog}}\cap\D^{1,2}(\Hf)$, entonces
\[
  D_t\int_0^Tu_s\,\dd W_s=u_t+\int_t^TD_tu_s\,\dd W_s.
\]
\end{corolario}

\begin{proof}
Por la Proposición~\ref{prop:heisenberg}, y porque $s\mapsto D_tu_s$ es adaptado y nulo para $s<t$, de
modo que $\delta(D_tu)$ es una integral de Itô en $[t,T]$.
\end{proof}

\section{La fórmula de Clark--Ocone}

\begin{teorema}[Clark--Ocone]\label{teo:clark-ocone}
Sea $\F_t$ la filtración browniana. Si $F\in\D^{1,2}$, entonces
\[
  F=\E F+\int_0^T\cond{D_tF}{\F_t}\,\dd W_t.
\]
\end{teorema}

\begin{proof}[Primera demostración (por dualidad)]
Por el Teorema~\ref{teo:representacion}, $F=\E F+\int_0^Tu_t\,\dd W_t$ con $u\in L^2_{\mathrm{prog}}$.
Sea $v\in L^2_{\mathrm{prog}}$ arbitrario. Por la isometría de Itô, el Teorema~\ref{teo:skorokhod-ito} y la dualidad,
\[
  \E\int_0^Tu_tv_t\,\dd t=\E\Big[F\int_0^Tv\,\dd W\Big]=\E[F\delta(v)]=\E\int_0^TD_tF\,v_t\,\dd t=\E\int_0^T\cond{D_tF}{\F_t}v_t\,\dd t,
\]
donde en el último paso usamos que $v_t$ es $\F_t$-medible y Fubini. Aquí $(\cond{D_tF}{\F_t})_t$ es la
proyección predecible de $DF$, que es progresiva. Como $v$ es arbitrario en $L^2_{\mathrm{prog}}$, $u_t=\cond{D_tF}{\F_t}$
$\dd t\otimes\dd\P$-c.s.
\end{proof}

\begin{proof}[Segunda demostración (por caos)]
Si $F=I_n(f_n)$, por el Teorema~\ref{teo:iteradas},
\[
  F=n\int_0^TI_{n-1}\big(f_n(\cdot,t)\ind_{[0,t]}^{\otimes(n-1)}\big)\,\dd W_t.
\]
Por el Teorema~\ref{teo:derivada-caos} y la Proposición~\ref{prop:derivada-esp-cond},
\[
  \cond{D_tF}{\F_t}=nI_{n-1}\big(f_n(\cdot,t)\ind_{[0,t]}^{\otimes(n-1)}\big).
\]
Se suma en $n$.
\end{proof}

\begin{ejemplo}\label{ej:clark-ocone}
\begin{enumerate}[label=(\alph*)]
  \item $F=W_T^3$: $D_tF=3W_T^2$, $\cond{3W_T^2}{\F_t}=3(W_t^2+T-t)$. Así $W_T^3=
  \int_0^T3(W_t^2+T-t)\,\dd W_t$.
  \item $F=\int_0^TW_s\,\dd s$: $D_tF=T-t$, determinista. $F=\int_0^T(T-t)\,\dd W_t$.
  \item $F=M=\sup_{[0,T]}W$: $D_tM=\ind_{\{t\le\tau\}}$, luego $\cond{D_tM}{\F_t}=\P(\tau\ge t\mid\F_t)=
  \P\big(\sup_{s\in[t,T]}(W_s-W_t)\ge M_t-W_t\mid\F_t\big)=2\bar\Phi\big(\frac{M_t-W_t}{\sqrt{T-t}}\big)$, con $M_t=\sup_{s\le t}W_s$ y $\bar\Phi=1-\Phi$,
  usando el principio de reflexión.
\end{enumerate}
\end{ejemplo}

\begin{intuicion}
Clark--Ocone dice: \emph{la sensibilidad de $F$ al ruido en $t$, promediada con la información
disponible en $t$, es la exposición instantánea de la martingala $\cond{F}{\F_t}$ al incremento $\dd W_t$.}
En finanzas, si $F$ es un pago descontado, $\cond{D_tF}{\F_t}$ es (salvo un factor $\sigma S_t$) la cantidad
de activo de la cartera de cobertura en el instante $t$: es el \emph{delta}. Véase el Capítulo~\ref{cap:cobertura}.
\end{intuicion}

\section{Clark--Ocone bajo cambio de medida}
En finanzas, el pago se representa bajo una medida equivalente $\Q$, mientras que la derivada se
toma respecto al browniano original. El resultado siguiente, debido a Ocone y Karatzas~\cite{OconeKaratzas91},
resuelve el desajuste.

\begin{teorema}[Clark--Ocone generalizado]\label{teo:clark-ocone-generalizado}
Sea $\theta$ adaptado y acotado con $\theta_s\in\D^{1,2}$ y
condiciones de integrabilidad adecuadas, $\dd\Q=\mathcal E(-\theta)_T\,\dd\P$ y $\tilde W_t=W_t+\int_0^t\theta_s\,\dd s$, que es
$\Q$-browniano. Si $F\in\D^{1,2}$ y es suficientemente integrable bajo $\Q$,
\[
  F=\E_\Q F+\int_0^T\E_\Q\Big[D_tF-F\int_t^TD_t\theta_s\,\dd\tilde W_s\,\Big|\,\F_t\Big]\,\dd\tilde W_t.
\]
\end{teorema}

\begin{proof}[Esbozo]
Sea $Z_t=\mathcal E(-\theta)_t$. Por Bayes, $\E_\Q[F\mid\F_t]=Z_t^{-1}\E[FZ_T\mid\F_t]$. Se aplica Clark--Ocone
(bajo $\P$) a $FZ_T$, con $D_t(FZ_T)=Z_TD_tF+FD_tZ_T$, y la derivada de $\log Z_T=-\int\theta\,\dd W-\frac12\int\theta^2$
es $D_t\log Z_T=-\theta_t-\int_t^TD_t\theta_s\,\dd W_s-\int_t^T\theta_sD_t\theta_s\,\dd s=-\theta_t-\int_t^TD_t\theta_s\,\dd\tilde W_s$. Se combina con la
fórmula de Itô para el cociente $\E[FZ_T\mid\F_t]/Z_t$.
\end{proof}

\chapter{El operador de Ornstein--Uhlenbeck}\label{cap:ou}

\begin{panorama}
El tercer operador básico es $L=-\delta D$, generador del semigrupo de Ornstein--Uhlenbeck
en dimensión infinita. Es diagonal en el caos: $L=-\sum_nnJ_n$. Probamos la fórmula
de Mehler, la propiedad de difusión, la desigualdad de Poincaré y la representación de
covarianzas, la desigualdad logarítmica de Sobolev de Gross, la hipercontractividad de
Nelson y enunciamos (con esbozo de prueba) las desigualdades de Meyer, que hacen de
los espacios $\D^{k,p}$ una escala bien comportada.
\end{panorama}

\section{El semigrupo y su generador}

\begin{definicion}
El \emph{semigrupo de Ornstein--Uhlenbeck} es la familia de contracciones de $L^2(\Omega)$
\[
  P_tF=\sum_{n=0}^\infty e^{-nt}J_nF,\qquad t\ge0.
\]
Su generador es $L=-\sum_nnJ_n$, con dominio $\Dom L=\{F:\sum_nn^2\E(J_nF)^2<\infty\}$. La
\emph{pseudoinversa} es $L^{-1}F=-\sum_{n\ge1}\frac1nJ_nF$, que satisface $LL^{-1}F=F-\E F$.
\end{definicion}

\begin{teorema}[Fórmula de Mehler]\label{teo:mehler}
Sea $W'$ una copia independiente de $W$, definida en
$(\Omega',\F',\P')$. Para $F\in L^2(\Omega)$, viendo $F$ como funcional de $W$, $F=\Psi(W)$,
\[
  P_tF=\E'\Big[\Psi\big(e^{-t}W+\sqrt{1-e^{-2t}}\,W'\big)\Big].
\]
\end{teorema}

\begin{proof}
El proceso $W^{(t)}=e^{-t}W+\sqrt{1-e^{-2t}}W'$ es isonormal sobre $\Hf$, luego la fórmula
define una contracción de $L^p$ para todo $p$ (Jensen). Basta comprobarla en $F=\mathcal E(h)=
e^{W(h)-\frac12\norm{h}^2}$, que generan un subespacio denso:
\[
  \E'\big[e^{e^{-t}W(h)+\sqrt{1-e^{-2t}}W'(h)-\frac12\norm{h}^2}\big]=e^{e^{-t}W(h)-\frac12e^{-2t}\norm{h}^2}=\mathcal E(e^{-t}h).
\]
Por otra parte, $\mathcal E(h)=\sum_n\frac1{n!}I_n(h^{\otimes n})$, de modo que $\sum_ne^{-nt}\frac1{n!}I_n(h^{\otimes n})=\sum_n\frac1{n!}I_n((e^{-t}h)^{\otimes n})=
\mathcal E(e^{-t}h)$.
\end{proof}

\begin{corolario}
$(P_t)$ es un semigrupo markoviano: $P_t1=1$, $P_tF\ge0$ si $F\ge0$, y $\norm{P_tF}_p\le\norm{F}_p$
para todo $p\ge1$. Es simétrico respecto a $\P$ y ergódico: $P_tF\to\E F$ cuando $t\to\infty$.
\end{corolario}

\begin{proposicion}
$F\in\Dom L$ si y solo si $F\in\D^{1,2}$ y $DF\in\Dom\delta$, y en ese caso $LF=-\delta DF$.
\end{proposicion}

\begin{proof}
Para $F=I_n(f_n)$: $DF=nI_{n-1}(f_n(\cdot,t))$ y, por la Proposición~\ref{prop:delta-caos}, $\delta DF=
nI_n(f_n)=-LF$. Las condiciones de dominio coinciden con $\sum n^2n!\norm{f_n}^2<\infty$.
\end{proof}

\begin{proposicion}[Propiedades de cálculo]
\begin{enumerate}[label=(\roman*)]
  \item Para $F=f(W(h_1),\dots,W(h_n))\in\Sc$,
  \[
    LF=\sum_{i,j}\partial_{ij}f(\cdot)\ip{h_i}{h_j}-\sum_i\partial_if(\cdot)W(h_i).
  \]
  \item \textup{Propiedad de difusión}: si $F_1,\dots,F_m\in\Sc$ y $\varphi\in C^2$,
  \[
    L\varphi(F)=\sum_i\partial_i\varphi(F)LF_i+\sum_{i,j}\partial_{ij}\varphi(F)\ip{DF_i}{DF_j}.
  \]
  En particular $\Gamma(F,G):=\frac12\big(L(FG)-FLG-GLF\big)=\ip{DF}{DG}$.
  \item \textup{Conmutación}: $DP_tF=e^{-t}P_tDF$, donde $P_t$ actúa componente a componente sobre las
  variables $\Hf$-valuadas.
\end{enumerate}
\end{proposicion}

\begin{proof}
(i) Es $-\delta DF$ con la Proposición~\ref{prop:reglas-delta}(i). (ii) Se deduce de (i). (iii) En caos:
$DP_tI_n(f)=e^{-nt}nI_{n-1}(f)$ y $P_tDI_n(f)=e^{-(n-1)t}nI_{n-1}(f)$.
\end{proof}

\begin{intuicion}
En dimensión finita, con $\gamma_n=N(0,I_n)$, $L=\Delta-x\cdot\nabla$ es el generador de la difusión de
Ornstein--Uhlenbeck $\dd X=-X\,\dd t+\sqrt2\,\dd B$, cuya medida invariante es $\gamma_n$. El cálculo de
Malliavin es, en este sentido, el análisis armónico de la medida gaussiana en dimensión
infinita: $D$ es el gradiente, $\delta=D^*$ la divergencia, $L=-D^*D$ el laplaciano (de Witten) y
el caos de Wiener su descomposición espectral, con autovalores $0,-1,-2,\dots$
\end{intuicion}

\section{Desigualdad de Poincaré y representación de covarianzas}

\begin{teorema}[Poincaré gaussiana]\label{teo:poincare}
Para $F\in\D^{1,2}$, $\Var(F)\le\E\norm{DF}^2_\Hf$.
\end{teorema}

\begin{proof}
$\Var F=\sum_{n\ge1}n!\norm{f_n}^2\le\sum_{n\ge1}n\,n!\norm{f_n}^2=\E\norm{DF}^2$.
\end{proof}

La constante 1 es óptima (igualdad en el primer caos) y no depende de la dimensión: esta
es la fuente de los fenómenos de \emph{concentración de la medida} gaussiana.

\begin{proposicion}[Representación de covarianzas]\label{prop:covarianzas}
Para $F,G\in\D^{1,2}$,
\[
  \Cov(F,G)=\E\ip{DF}{-DL^{-1}G}_\Hf=\int_0^\infty e^{-t}\,\E\ip{DF}{P_tDG}_\Hf\,\dd t.
\]
\end{proposicion}

\begin{proof}
$G-\E G=LL^{-1}G=-\delta DL^{-1}G$, luego, por dualidad, $\E[F(G-\E G)]=
\E\ip{DF}{-DL^{-1}G}$. Para la segunda forma: $-L^{-1}=\int_0^\infty P_t\,\dd t$ sobre variables centradas, y
$DP_t=e^{-t}P_tD$.
\end{proof}

\begin{corolario}[Concentración gaussiana]
Si $F\in\D^{1,2}$ y $\norm{DF}_\Hf\le1$ c.s., entonces $\P(F-\E F\ge
r)\le e^{-r^2/2}$ para $r\ge0$.
\end{corolario}

\begin{proof}
Podemos suponer $\E F=0$. Sea $\phi(\lambda)=\E e^{\lambda F}$. Por la Proposición~\ref{prop:covarianzas} con
$G=e^{\lambda F}$, para la que $DG=\lambda e^{\lambda F}DF$,
\[
  \phi'(\lambda)=\E[Fe^{\lambda F}]=\int_0^\infty e^{-t}\E\ip{\lambda e^{\lambda F}DF}{P_tDF}\,\dd t\le\lambda\phi(\lambda),
\]
porque $\norm{P_tDF}\le P_t\norm{DF}\le1$ y $\int e^{-t}=1$. Así $\phi(\lambda)\le e^{\lambda^2/2}$ y se aplica Chebyshev exponencial
con $\lambda=r$.
\end{proof}

\begin{ejemplo}
$M=\sup_{[0,1]}W$ satisface $\norm{DM}^2=\tau\le1$, luego $\P(M-\E M\ge r)\le e^{-r^2/2}$. Más
en general, el supremo de cualquier proceso gaussiano acotado con $\sup_t\Var X_t\le\sigma^2$ satisface la
desigualdad de Borell--TIS, $\P(\sup X-\E\sup X\ge r)\le e^{-r^2/2\sigma^2}$.
\end{ejemplo}

\section{Desigualdad logarítmica de Sobolev}

\begin{teorema}[Gross~\cite{Gross75}]\label{teo:log-sobolev}
Para $F\in\D^{1,2}$,
\[
  \Ent(F^2):=\E[F^2\log F^2]-\E F^2\log\E F^2\le2\,\E\norm{DF}^2_\Hf.
\]
\end{teorema}

\begin{proof}
Probamos la forma equivalente $\Ent(G)\le\frac12\E[\norm{DG}^2/G]$ para $G=F^2$ suave
y $G\ge\varepsilon>0$ (luego se aproxima). Por ergodicidad, $\Ent(G)=-\int_0^\infty\frac{\dd}{\dd t}\E[P_tG\log P_tG]\,\dd t$.
Derivando y usando $\E[LH\cdot K]=-\E\ip{DH}{DK}$,
\[
  -\frac{\dd}{\dd t}\E[P_tG\log P_tG]=-\E[LP_tG\,(1+\log P_tG)]=\E\frac{\norm{DP_tG}^2}{P_tG}.
\]
Por la conmutación $DP_tG=e^{-t}P_tDG$ y Cauchy--Schwarz para la medida de Mehler, $\norm{P_tDG}^2\le
(P_t\norm{DG})^2\le P_tG\cdot P_t(\norm{DG}^2/G)$. Así
\[
  \E\frac{\norm{DP_tG}^2}{P_tG}\le e^{-2t}\E P_t\Big(\frac{\norm{DG}^2}{G}\Big)=e^{-2t}\E\frac{\norm{DG}^2}{G},
\]
e integrando, $\Ent(G)\le\frac12\E[\norm{DG}^2/G]$. Con $G=F^2$, $\norm{DG}^2/G=4\norm{DF}^2$.
\end{proof}

\section{Hipercontractividad}

\begin{teorema}[Nelson]\label{teo:hipercontractividad}
Sean $1<p<q<\infty$ y $t>0$ con $e^{2t}\ge\frac{q-1}{p-1}$. Entonces
$\norm{P_tF}_q\le\norm{F}_p$ para toda $F\in L^p$.
\end{teorema}

\begin{proof}[Demostración (Gross)]
Sea $F\ge0$ suave, acotada inferiormente por una constante positiva,
$q(t)=1+(p-1)e^{2t}$ y $N(t)=\norm{P_tF}_{q(t)}$. Basta ver que $N$ es decreciente. Con $G_t=P_tF$ y
$q=q(t)$, un cálculo directo da
\[
  \frac{\dd}{\dd t}\log N(t)=\frac{q'}{q^2}\frac{\Ent(G_t^q)}{\E G_t^q}+\frac{\E[G_t^{q-1}LG_t]}{\E G_t^q}.
\]
Ahora $\E[G^{q-1}LG]=-(q-1)\E[G^{q-2}\norm{DG}^2]=-\frac{4(q-1)}{q^2}\E\norm{D(G^{q/2})}^2$. Por la desigualdad
log-Sobolev aplicada a $G^{q/2}$, $\Ent(G^q)\le2\E\norm{D(G^{q/2})}^2$. Así
\[
  \frac{\dd}{\dd t}\log N\le\frac{2}{q^2\E G^q}\big(q'-2(q-1)\big)\E\norm{D(G^{q/2})}^2=0,
\]
porque $q'=2(q-1)$. Para $F$ de signo arbitrario se usa $\abs{P_tF}\le P_t\abs{F}$.
\end{proof}

\begin{proof}[Demostración del Teorema~\ref{teo:normas-caos}]
Si $F\in\Hc_n$, $P_tF=e^{-nt}F$. Con $p=2$ y $e^{2t}=q-1$:
$e^{-nt}\norm{F}_q\le\norm{F}_2$, es decir, $\norm{F}_q\le(q-1)^{n/2}\norm{F}_2$.
\end{proof}

\section{Desigualdades de Meyer}

\begin{teorema}[Meyer, 1984]\label{teo:meyer}
Para cada $p\in(1,\infty)$ existen $c_p,C_p>0$ tales que para toda
$F\in\Sc$,
\[
  c_p\norm{DF}_{L^p(\Omega;\Hf)}\le\norm{(-L)^{1/2}F}_{L^p}\le C_p\big(\norm{DF}_{L^p(\Omega;\Hf)}+\abs{\E F}\big).
\]
Más en general, $\norm{F}_{k,p}$ es equivalente a $\norm{(I-L)^{k/2}F}_p$, y $\D^{k,p}=\Dom_{L^p}(I-L)^{k/2}$.
\end{teorema}

\begin{proof}[Esbozo (método de Pisier)]
La idea es expresar la ``transformada de Riesz'' $D(-L)^{-1/2}$ como
un operador singular de tipo transformada de Hilbert. Se parte de la fórmula de subordinación
$(-L)^{-1/2}=\frac1{\sqrt\pi}\int_0^\infty t^{-1/2}P_t\,\dd t$ sobre variables centradas, de la fórmula de Mehler escrita con
$e^{-t}=\cos\theta$, $\sqrt{1-e^{-2t}}=\sin\theta$, y de la identidad, obtenida derivando en $\theta$ e integrando por
partes en $W'$,
\[
  \frac{\dd}{\dd\theta}\E'\big[F(W\cos\theta+W'\sin\theta)\big]\ \leftrightarrow\ \E'\big[\ip{DF(W\cos\theta+W'\sin\theta)}{\cdot}\,W'\big].
\]
El resultado (Pisier) es que $\ip{D(-L)^{-1/2}F}{h}$ se escribe como una integral en valor principal en
$\theta\in(-\pi/2,\pi/2)$ de $\E'[F(W\cos\theta+W'\sin\theta)W'(h)]$ contra un núcleo explícito que se comporta
como $1/\theta$ cerca de 0. Rotar $(W,W')\mapsto(W\cos\theta+W'\sin\theta,-W\sin\theta+W'\cos\theta)$ preserva la ley
conjunta, y el operador resultante es esencialmente la transformada de Hilbert en el círculo,
acotada en $L^p$ por el teorema de M.~Riesz (en su versión con valores en $L^p(\Omega')$, que es un
espacio UMD). Combinando con la desigualdad de Khintchine gaussiana para pasar de $W'(h)$ a
$\norm{\cdot}_\Hf$, se obtiene la primera desigualdad. La segunda se obtiene por dualidad. Véanse \cite[§1.5]{Nualart06}
y~\cite{Pisier88}.
\end{proof}

\begin{corolario}[Continuidad de $\delta$]\label{cor:continuidad-delta}
Para $p>1$, $\D^{1,p}(\Hf)\subset\Dom\delta$ y
\[
  \norm{\delta(u)}_p\le C_p\big(\norm{\E u}_\Hf+\norm{Du}_{L^p(\Omega;\Hf\otimes\Hf)}\big).
\]
Más en general $\delta$ es continuo de $\D^{k+1,p}(\Hf)$ en $\D^{k,p}$, y $\D^\infty(\Hf)\subset\Dom\delta$ con $\delta(\D^\infty(\Hf))\subset\D^\infty$.
\end{corolario}

\begin{proof}
Por dualidad: para $G\in\Sc$ centrada, $\E[G\delta(u)]=\E\ip{DG}{u}$, y se escribe $DG=
D(-L)^{-1/2}(-L)^{1/2}G$; Meyer controla $\norm{D(-L)^{-1/2}}_{L^q\to L^q(\Hf)}$, y se usa la relación de conmutación
para pasar una derivada a $u$.
\end{proof}

\begin{observacion}
Las desigualdades de Meyer son el ingrediente técnico que hace funcionar la
integración por partes iterada del Capítulo~\ref{cap:densidades}: allí se necesita que $\delta$ envíe $\D^\infty$ en $\D^\infty$ con cotas
$L^p$ controladas, para poder aplicar $\delta$ repetidamente.
\end{observacion}


\part{Densidades y el teorema de Hörmander}
\chapter{Existencia y regularidad de densidades}\label{cap:densidades}

\begin{panorama}
Esta es la aplicación fundacional del cálculo de Malliavin. Dada una variable $F=
(F^1,\dots,F^d)$ en el espacio de Wiener, queremos saber si su ley tiene densidad y si esta
es regular. La herramienta clave es la \emph{matriz de Malliavin} $\gamma_F=(\ip{DF^i}{DF^j}_\Hf)_{ij}$: si
es invertible, la ley de $F$ es absolutamente continua; si además $(\det\gamma_F)^{-1}$ tiene todos
los momentos y $F\in\D^\infty$, la densidad es $C^\infty$. Probamos ambos resultados, damos una
fórmula explícita para la densidad y, al final, la fórmula de Nourdin--Viens.
\end{panorama}

\section{La idea en dimensión uno}
Sea $F\in\D^{1,2}$ y $\varphi\in C_b^1$. Por la regla de la cadena, $\ip{D\varphi(F)}{DF}=\varphi'(F)\norm{DF}^2$. Si $\norm{DF}>0$,
podemos despejar
\[
  \varphi'(F)=\Big\langle D\varphi(F),\frac{DF}{\norm{DF}^2}\Big\rangle,
\]
y tomando esperanzas y usando la dualidad,
\begin{equation}\label{eq:ibp-1d}
  \E[\varphi'(F)]=\E\Big[\varphi(F)\,\delta\Big(\frac{DF}{\norm{DF}^2}\Big)\Big].
\end{equation}
Hemos \emph{trasladado la derivada de $\varphi$ a un peso}. Con $\varphi(y)=\int_{-\infty}^y\ind_{(a,b]}$, formalmente, $\P(a<F\le
b)=\E[\varphi'(F)]$ se controla por $\norm{\varphi}_\infty\le b-a$: la ley de $F$ no tiene átomos y, de hecho, tiene
densidad acotada.

\begin{teorema}[Fórmula de la densidad]\label{teo:densidad}
Sea $F\in\D^{1,2}$ tal que $\norm{DF}_\Hf>0$ c.s.\ y $u=
DF/\norm{DF}^2_\Hf\in\Dom\delta$. Entonces $F$ tiene una densidad continua y acotada
\[
  p_F(x)=\E\big[\ind_{\{F>x\}}\,\delta(u)\big].
\]
\end{teorema}

\begin{proof}
Sea $\psi\ge0$ continua con soporte compacto y $\varphi(y)=\int_{-\infty}^y\psi(z)\,\dd z$, que es $C^1$ y
acotada. Por~\eqref{eq:ibp-1d} (válida por aproximación),
\[
  \E[\psi(F)]=\E\Big[\delta(u)\int_{-\infty}^F\psi(x)\,\dd x\Big]=\int_\R\psi(x)\,\E\big[\ind_{\{F>x\}}\delta(u)\big]\,\dd x,
\]
por Fubini, ya que $\delta(u)\in L^1$. Así la ley de $F$ tiene densidad $x\mapsto\E[\ind_{\{F>x\}}\delta(u)]$, que es continua
(convergencia dominada, pues $\P(F=x)=0$) y acotada por $\E\abs{\delta(u)}$.
\end{proof}

\begin{observacion}
Como $\E\delta(u)=0$, también $p_F(x)=-\E[\ind_{\{F\le x\}}\delta(u)]$. Combinando ambas y
Hölder se obtienen \emph{cotas de colas}: $p_F(x)\le\norm{\delta(u)}_q\,\P(\abs{F}\ge\abs{x})^{1/p}$.
\end{observacion}

\section{El criterio de Bouleau--Hirsch}
La condición $u\in\Dom\delta$ es restrictiva. Sorprendentemente, la mera no degeneración basta para
la continuidad absoluta.

\begin{teorema}[Bouleau--Hirsch]\label{teo:bouleau-hirsch}
Sea $F=(F^1,\dots,F^d)$ con $F^i\in\D^{1,2}$ (o incluso $\D^{1,1}_{\mathrm{loc}}$). Si
la matriz de Malliavin $\gamma_F=\big(\ip{DF^i}{DF^j}_\Hf\big)_{i,j}$ es invertible c.s., entonces la ley de $F$ es
absolutamente continua respecto a la medida de Lebesgue en $\R^d$.
\end{teorema}

\begin{proof}[Demostración para $d=1$]
Hay que ver que si $B$ es boreliano con $\abs{B}=0$, entonces $\P(F\in B)=
0$. Podemos suponer $B$ acotado. Sea $\varphi(y)=\int_{-\infty}^y\ind_B(z)\,\dd z\equiv0$. Aproximamos $\ind_B$ por funciones
continuas $\psi_n\ge0$ con $\psi_n\to\ind_B$ en $L^1(\R,\dd x+\P\circ F^{-1})$ y uniformemente acotadas, y sea
$\varphi_n=\int_{-\infty}^\cdot\psi_n\to\varphi=0$ uniformemente. Por la regla de la cadena, $D\varphi_n(F)=\psi_n(F)DF$. Ahora
$\varphi_n(F)\to0$ en $L^2$ y $D\varphi_n(F)=\psi_n(F)DF\to\ind_B(F)DF$ en $L^2(\Omega;\Hf)$ (convergencia dominada,
con $\psi_n(F)\to\ind_B(F)$ en probabilidad). Por cerrabilidad de $D$, $\ind_B(F)DF=D0=0$. Como
$\norm{DF}>0$ c.s., $\ind_B(F)=0$ c.s.
\end{proof}

\begin{proof}[Idea para $d\ge2$]
El argumento anterior no se generaliza directamente, porque en $\R^d$ un conjunto
nulo no es la imagen de una primitiva. Bouleau y Hirsch usan una \emph{fórmula de la coárea} en
el espacio de Wiener, o alternativamente se prueba que $\P\circ F^{-1}\ll\mathrm{Leb}$ mediante el criterio:
para toda $\varphi\in C_b^1$, $\abs{\E[\partial_i\varphi(F)G]}\le C_G\norm{\varphi}_\infty$ para suficientes $G$; se toma $G$ que localiza en
$\{\abs{\gamma_F^{-1}}\le N\}$ y se usa el resultado de que una medida $\mu$ en $\R^d$ cuyas derivadas distribucionales
son medidas finitas es absolutamente continua (lema de Malliavin). Véase \cite[§2.1]{Nualart06} y~\cite{BouleauHirsch91}.
\end{proof}

\section{Integración por partes iterada y densidades regulares}

\begin{definicion}
Un vector $F\in(\D^\infty)^d$ es \emph{no degenerado} si $\gamma_F$ es invertible c.s.\ y
$(\det\gamma_F)^{-1}\in\bigcap_{p\ge1}L^p$.
\end{definicion}

\begin{lema}\label{lem:inversa-malliavin}
Si $F$ es no degenerado, entonces $(\gamma_F^{-1})_{ij}\in\D^\infty$ y
\[
  D(\gamma_F^{-1})_{ij}=-\sum_{k,l}(\gamma_F^{-1})_{ik}(\gamma_F^{-1})_{lj}\,D(\gamma_F)_{kl}.
\]
\end{lema}

\begin{proof}
$\gamma_F\in\D^\infty$ porque $D$ conmuta con productos (Leibniz) y Meyer. La inversa es
$\operatorname{adj}(\gamma_F)/\det\gamma_F$. Se aproxima $1/\det\gamma_F$ por $(\det\gamma_F+\varepsilon)^{-1}$, que es suave, y se pasa al límite
usando que $(\det\gamma_F)^{-1}$ tiene todos los momentos. La fórmula es la derivada de la inversa de
una matriz.
\end{proof}

\begin{proposicion}[Integración por partes]\label{prop:ibp-iterada}
Sea $F$ no degenerado y $G\in\D^\infty$. Para cada
multiíndice $\alpha=(\alpha_1,\dots,\alpha_k)\in\{1,\dots,d\}^k$ existe $H_\alpha(F,G)\in\D^\infty$ tal que, para toda $\varphi\in
C_p^\infty(\R^d)$,
\[
  \E\big[\partial_\alpha\varphi(F)\,G\big]=\E\big[\varphi(F)\,H_\alpha(F,G)\big].
\]
Los pesos se definen recursivamente por
\[
  H_{(i)}(F,G)=\sum_{j=1}^d\delta\big(G\,(\gamma_F^{-1})_{ij}\,DF^j\big),\qquad H_\alpha(F,G)=H_{(\alpha_k)}\big(F,H_{(\alpha_1,\dots,\alpha_{k-1})}(F,G)\big),
\]
y satisfacen $\norm{H_\alpha(F,G)}_p\le C\norm{\det\gamma_F^{-1}}_q^a\norm{DF}_{k,r}^b\norm{G}_{k,s}$ para exponentes adecuados.
\end{proposicion}

\begin{proof}
Por la regla de la cadena, $\ip{D\varphi(F)}{DF^j}=\sum_l\partial_l\varphi(F)(\gamma_F)_{lj}$. Invirtiendo,
$\partial_i\varphi(F)=\sum_j\ip{D\varphi(F)}{DF^j}(\gamma_F^{-1})_{ji}$. Multiplicando por $G$ y tomando esperanzas,
\[
  \E[\partial_i\varphi(F)G]=\E\Big\langle D\varphi(F),\sum_jG(\gamma_F^{-1})_{ij}DF^j\Big\rangle=\E\big[\varphi(F)H_{(i)}(F,G)\big],
\]
usando la simetría de $\gamma_F$ y la dualidad; el integrando está en $\D^\infty(\Hf)\subset\Dom\delta$ por el Lema~\ref{lem:inversa-malliavin},
y $\delta$ preserva $\D^\infty$ (Corolario~\ref{cor:continuidad-delta}). La recursión y las cotas siguen por inducción, usando Meyer
y Hölder.
\end{proof}

\begin{teorema}[Malliavin]\label{teo:malliavin-densidad}
Si $F$ es no degenerado, la ley de $F$ tiene una densidad $p_F\in\mathcal S(\R^d)$
(de Schwartz). Además
\[
  p_F(x)=\E\big[\ind_{\{F>x\}}\,H_{(1,2,\dots,d)}(F,1)\big],
\]
donde $\{F>x\}=\{F^i>x_i\ \forall i\}$, y para todo multiíndice $\beta$, $\partial_\beta p_F(x)=
(-1)^{\abs\beta}\E[\ind_{\{F>x\}}H_{(1,\dots,d,\beta)}(F,1)]$.
\end{teorema}

\begin{proof}
Sea $\psi\in C_c^\infty(\R^d)$ y $\varphi(y)=\int_{-\infty}^{y_1}\cdots\int_{-\infty}^{y_d}\psi(z)\,\dd z$, de modo que $\partial_{1\cdots d}\varphi=\psi$. Por
la Proposición~\ref{prop:ibp-iterada} y Fubini,
\[
  \E\psi(F)=\E[\varphi(F)H_{(1,\dots,d)}(F,1)]=\int_{\R^d}\psi(x)\E\big[\ind_{\{F>x\}}H_{(1,\dots,d)}(F,1)\big]\,\dd x.
\]
Para las derivadas, la transformada de Fourier: con $\varphi=e^{i\ip{\xi}{\cdot}}$, $\abs{\xi^\alpha}\abs{\hat\mu_F(\xi)}=\abs{\E[\partial_\alpha\varphi(F)]}\le
\E\abs{H_\alpha(F,1)}<\infty$, luego $\hat\mu_F$ decrece más rápido que cualquier polinomio y $p_F\in C^\infty$. El decrecimiento
de $p_F$ y sus derivadas en el infinito sale de la fórmula con $\ind_{\{F>x\}}$ (o su complemento) y
Hölder, usando $\P(\abs{F}>\abs{x})\le C_N\abs{x}^{-N}$.
\end{proof}

\begin{intuicion}
El esquema de la prueba es el de las ``integraciones por partes iteradas'': cada derivada de
$\varphi$ se transfiere a un peso, pagando con una potencia de $\gamma_F^{-1}$ y una derivada de Malliavin
más. La no degeneración es exactamente lo que hace que todos esos pesos sean integrables.
En el Capítulo~\ref{cap:hormander} veremos que, para soluciones de EDEs, la condición de corchetes de
Hörmander implica la no degeneración. En el Capítulo~\ref{cap:griegas}, que los \emph{mismos pesos $H_\alpha$} dan
las griegas en finanzas.
\end{intuicion}

\section{Densidades locales y positividad}
A menudo $\gamma_F$ solo es no degenerada en un abierto. Si $F\in(\D^\infty)^d$ y $\gamma_F$ es no degenerada \emph{en
un abierto} $U\subset\R^d$ (es decir, $\E[\ind_{\{F\in U\}}(\det\gamma_F)^{-p}]<\infty$ para todo $p$, con localización suave), la
restricción de la ley a $U$ tiene densidad $C^\infty$. Esto se prueba multiplicando por $G=\chi(F)$ con
$\chi\in C_c^\infty(U)$. Véase también~\cite{BallyCaramellino11}, que obtiene estos criterios con transformadas de Riesz.

\begin{teorema}[Fang; Hirsch--Song]
Si $F\in\D^{1,p}$, $p>1$, el soporte de la ley de $F$ es conexo.
En particular, si $F\in\D^{1,2}$ es real, el soporte es un intervalo cerrado.
\end{teorema}

La \emph{positividad estricta} de la densidad es un problema más delicado; un criterio general es el
de Bally--Pardoux, basado en la fórmula de Cameron--Martin: si para algún $h\in H$ el funcional
trasladado $F(\omega+h)$ cae en un entorno de $y$ con matriz de Malliavin no degenerada, entonces
$p_F(y)>0$. Lo usaremos en el Capítulo~\ref{cap:densidades-edpe} para la ecuación del calor.

\section{La fórmula de Nourdin--Viens}
Un enfoque distinto, más reciente, evita invertir la matriz de Malliavin usando el operador $L^{-1}$.

\begin{teorema}[Nourdin--Viens, 2009]\label{teo:nourdin-viens}
Sea $F\in\D^{1,2}$ con $\E F=0$, y sea
\[
  g_F(x)=\E\big[\ip{DF}{-DL^{-1}F}_\Hf\mid F=x\big].
\]
Entonces $g_F\ge0$ c.s.\ sobre el soporte de la ley. Además, $F$ tiene densidad si y solo si
$g_F(F)>0$ c.s., y en ese caso el soporte es un intervalo $[\alpha,\beta]\ni0$ y, para $x\in(\alpha,\beta)$,
\[
  p_F(x)=\frac{\E\abs{F}}{2g_F(x)}\exp\Big(-\int_0^x\frac{y\,\dd y}{g_F(y)}\Big).
\]
\end{teorema}

\begin{proof}[Demostración de la fórmula]
Sea $\varphi\in C_b^1$ con primitiva $\Phi$. Por la Proposición~\ref{prop:covarianzas} con $G=F$,
$\E[F\varphi(F)]=\E[\varphi'(F)\ip{DF}{-DL^{-1}F}]=\E[\varphi'(F)g_F(F)]$. Si $F$ tiene densidad $\rho$, esto dice que,
distribucionalmente,
\[
  \int x\varphi(x)\rho(x)\,\dd x=\int\varphi'(x)g_F(x)\rho(x)\,\dd x\quad\Longrightarrow\quad(g_F\rho)'=-x\rho.
\]
Así $g_F(x)\rho(x)=\int_x^\infty y\rho(y)\,\dd y$, que es $>0$ en el interior del soporte; en particular
$g_F\ge0$. Llamando $\Psi(x)=\int_x^\infty y\rho(y)\,\dd y$, la relación $\Psi'=-x\rho=-x\Psi/g_F$ da
$\Psi(x)=\Psi(0)\exp(-\int_0^xy/g_F(y)\,\dd y)$ con $\Psi(0)=\E[F^+]=\frac12\E\abs{F}$. Dividiendo por $g_F$ se
obtiene la fórmula. La equivalencia con la existencia de densidad requiere un argumento
adicional; véase~\cite{NourdinViens09}.
\end{proof}

\begin{corolario}[Cotas gaussianas]
Si $0<\sigma^2_{\min}\le g_F(F)\le\sigma^2_{\max}$ c.s., entonces
\[
  \frac{\E\abs{F}}{2\sigma^2_{\max}}e^{-x^2/2\sigma^2_{\min}}\le p_F(x)\le\frac{\E\abs{F}}{2\sigma^2_{\min}}e^{-x^2/2\sigma^2_{\max}}.
\]
\end{corolario}

\begin{ejemplo}
En el primer caos, $F=W(h)$, $-DL^{-1}F=h$ y $g_F\equiv\norm{h}^2$: se recupera la
densidad gaussiana. Para el supremo $M$ del browniano, centrado, se tiene $g_{M-\E M}\le1$ (porque
$\norm{DM}\le1$ y $\norm{P_tDM}\le1$ en la representación integral de $-DL^{-1}$); combinado con una cota
inferior de $g$ en una región, se obtienen cotas gaussianas explícitas para la densidad allí. El
método se ha aplicado con éxito a soluciones de EDPEs y a polímeros dirigidos.
\end{ejemplo}

\chapter{EDEs y el teorema de Hörmander}\label{cap:hormander}

\begin{panorama}
Calculamos la derivada de Malliavin de la solución de una EDE y probamos que $X_t\in\D^\infty$.
Expresamos la matriz de Malliavin con el proceso tangente, tratamos el caso elíptico
y demostramos el teorema que motivó la teoría: bajo la condición de corchetes de
Hörmander, $X_t$ tiene densidad $C^\infty$. La prueba que presentamos sigue a Norris~\cite{Norris86} y
Hairer~\cite{Hairer11}.
\end{panorama}

Consideramos la EDE de Stratonovich en $\R^d$
\begin{equation}\label{eq:ede-strat}
  \dd X_t=V_0(X_t)\,\dd t+\sum_{k=1}^mV_k(X_t)\circ\dd W^k_t,\qquad X_0=x,
\end{equation}
con campos $V_k\in C_b^\infty$ (derivadas de orden $\ge1$ acotadas). Sean $Y_t=\nabla_xX_t$ y $Z_t=Y_t^{-1}$
(Teorema~\ref{teo:kunita}). En notación de Stratonovich,
\[
  \dd Y_t=\nabla V_0(X_t)Y_t\,\dd t+\sum_k\nabla V_k(X_t)Y_t\circ\dd W^k_t,\qquad\dd Z_t=-Z_t\nabla V_0(X_t)\,\dd t-\sum_kZ_t\nabla V_k(X_t)\circ\dd W^k_t.
\]

\section{La derivada de Malliavin de la solución}

\begin{teorema}\label{teo:derivada-ede}
Para todo $t$, $X_t\in(\D^\infty)^d$, y para $r\le t$,
\[
  D^k_rX_t=Y_tZ_rV_k(X_r),\qquad k=1,\dots,m,
\]
mientras que $D_rX_t=0$ si $r>t$. Además, $\sup_{r\le t}\E\norm{D_rX_t}^p<\infty$ para todo $p$, y lo mismo
para las derivadas de orden superior.
\end{teorema}

\begin{proof}
Escribimos la ecuación en forma de Itô, $\dd X=b(X)\,\dd t+\sigma_k(X)\,\dd W^k$. Las
iteraciones de Picard $X^{(n)}$ están en $\D^{1,2}$ (por el Corolario~\ref{cor:derivada-ito} y la regla de la cadena), y
derivando formalmente:
\[
  D^k_rX_t=\sigma_k(X_r)+\int_r^t\nabla b(X_s)D^k_rX_s\,\dd s+\sum_l\int_r^t\nabla\sigma_l(X_s)D^k_rX_s\,\dd W^l_s,\qquad r\le t.
\]
Las estimaciones BDG--Gronwall del Teorema~\ref{teo:existencia-ede} muestran que $\sup_n\E\sup_t\norm{DX^{(n)}_t}^p_\Hf<\infty$;
como $X^{(n)}\to X$ en $L^p$, por cerrabilidad (con un argumento de compacidad débil), $X_t\in\D^{1,p}$
y $DX$ satisface la ecuación lineal anterior. Pero esta ecuación, como función de $t\ge r$, es la
ecuación de primera variación~\eqref{eq:primera-variacion} con dato inicial $\sigma_k(X_r)$ en el instante $r$. Por unicidad y la
propiedad de flujo, su solución es $Y_tY_r^{-1}\sigma_k(X_r)=Y_tZ_rV_k(X_r)$ (recordemos que $\sigma_k=V_k$). Las
derivadas superiores se tratan por inducción: satisfacen ecuaciones lineales con términos fuente
polinomiales en derivadas de orden inferior.
\end{proof}

\begin{intuicion}
Un impulso en la $k$-ésima coordenada del ruido en el instante $r$ desplaza el estado en la
dirección $V_k(X_r)$; ese desplazamiento infinitesimal se transporta hasta $t$ por el linealizado
del flujo, $Y_tZ_r=\nabla_{X_r}X_{r\to t}$. Es exactamente lo que uno esperaría.
\end{intuicion}

\begin{corolario}[Matriz de Malliavin]\label{cor:matriz-malliavin}
La matriz de Malliavin de $X_t$ es
\[
  \gamma_{X_t}=\sum_{k=1}^m\int_0^tD^k_rX_t\,(D^k_rX_t)^\top\,\dd r=Y_t\,C_t\,Y_t^\top,\qquad C_t=\sum_{k=1}^m\int_0^tZ_rV_k(X_r)V_k(X_r)^\top Z_r^\top\,\dd r.
\]
$C_t$ es la \textup{matriz de Malliavin reducida}. Como $Y_t$ es invertible con inversa en todo $L^p$, $X_t$ es no
degenerado si y solo si $(\det C_t)^{-1}\in\bigcap_pL^p$.
\end{corolario}

\section{El caso elíptico}

\begin{teorema}
Si $\sigma\sigma^\top(x)\ge\lambda I$ con $\lambda>0$ para todo $x$, entonces para todo $t>0$, $X_t$ tiene
densidad $C^\infty$.
\end{teorema}

\begin{proof}
Para $\xi\in S^{d-1}$, $\xi^\top C_t\xi=\int_0^t\xi^\top Z_r\sigma\sigma^\top(X_r)Z_r^\top\xi\,\dd r\ge\lambda\int_0^t\abs{Z_r^\top\xi}^2\,\dd r\ge
\lambda\int_0^t\abs{Y_r}^{-2}\,\dd r$, ya que $\abs{\xi}\le\abs{Y_r^\top}\abs{Z_r^\top\xi}$. Así $\det C_t\ge(\lambda\int_0^t\abs{Y_r}^{-2}\,\dd r)^d$, y por Jensen,
$\big(\frac1t\int_0^t\abs{Y_r}^{-2}\,\dd r\big)^{-1}\le\frac1t\int_0^t\abs{Y_r}^2\,\dd r$, que tiene todos los momentos. Se concluye con el
Teorema~\ref{teo:malliavin-densidad}.
\end{proof}

\section{La condición de Hörmander}
En el caso degenerado, el ruido actúa solo en algunas direcciones, pero la deriva y la no
conmutatividad de los campos pueden ``propagarlo'' a todas. Recordemos el corchete de Lie
$[V,U]=\nabla U\,V-\nabla V\,U$, es decir, $[V,U]f=V(Uf)-U(Vf)$.

\begin{definicion}
Sea $\mathcal V_0=\{V_1,\dots,V_m\}$ y, recursivamente, $\mathcal V_{n+1}=\mathcal V_n\cup\{[U,V_j]:\ U\in
\mathcal V_n,\ j=0,\dots,m\}$. Se dice que se satisface la \emph{condición de Hörmander (parabólica)} en $x$ si
\[
  \operatorname{span}\{U(x):\ U\in\textstyle\bigcup_n\mathcal V_n\}=\R^d.
\]
\end{definicion}

Notemos que $V_0$ \emph{no} está en $\mathcal V_0$: la deriva sola no genera aleatoriedad, solo la transporta.

\begin{teorema}[Hörmander, 1967; Malliavin, 1976~\cite{Malliavin78}]\label{teo:hormander}
Si los campos $V_0,\dots,V_m\in C_b^\infty$
satisfacen la condición de Hörmander en $x$, entonces para todo $t>0$, $X_t^x$ es no degenerado
y su ley tiene densidad $p_t(x,\cdot)\in\mathcal S(\R^d)$.
\end{teorema}

\begin{ejemplo}[Kolmogorov]
$\dd X=V\,\dd t$, $\dd V=\dd W$ en $\R^2$. $V_0=(v,0)=v\partial_x$, $V_1=\partial_v$.
$[V_1,V_0]=\partial_x$. Los campos $V_1,[V_1,V_0]$ generan $\R^2$. Así $(X_t,V_t)=(\int_0^tW,W_t)$ tiene densidad
regular (de hecho gaussiana no degenerada, con covarianza $\big(\begin{smallmatrix}t^3/3&t^2/2\\t^2/2&t\end{smallmatrix}\big)$). Notemos la escala: la
dirección $\partial_x$, obtenida con un corchete, se ``llena'' en tiempo $t^{3/2}$ frente a $t^{1/2}$.
\end{ejemplo}

\begin{ejemplo}[Opción asiática]
$\dd S=\sigma S\,\dd W$ (Stratonovich: $V_1=\sigma s\partial_s$, $V_0=-\frac{\sigma^2}{2}s\partial_s$) y
$\dd A=S\,\dd t$. En $\R^2$ con $V_0=-\frac{\sigma^2}{2}s\partial_s+s\partial_a$ se tiene $[V_1,V_0]=\sigma s\partial_a$, y la condición se cumple en
$\{s>0\}$. Por tanto $(S_t,\int_0^tS_u\,\dd u)$ tiene densidad regular, lo que justifica las fórmulas de griegas
para opciones asiáticas del Capítulo~\ref{cap:griegas}.
\end{ejemplo}

\begin{ejemplo}[Área de Lévy]
$\dd X^1=\dd W^1$, $\dd X^2=\dd W^2$, $\dd X^3=\frac12(X^1\,\dd W^2-X^2\,\dd W^1)$.
$V_1=\partial_1-\frac{x_2}{2}\partial_3$, $V_2=\partial_2+\frac{x_1}{2}\partial_3$, $[V_1,V_2]=\partial_3$. Es el browniano en el grupo de Heisenberg, el
prototipo de la geometría sub-riemanniana.
\end{ejemplo}

\section{Demostración del teorema de Hörmander}
Por el Corolario~\ref{cor:matriz-malliavin} y el Teorema~\ref{teo:malliavin-densidad}, basta probar:
\begin{equation}\label{eq:hormander-objetivo}
  \text{para todo }p\ge1\text{ existe }C_p:\qquad\sup_{\abs\xi=1}\P\big(\xi^\top C_t\xi<\varepsilon\big)\le C_p\varepsilon^p\quad\text{para }\varepsilon\in(0,1].
\end{equation}
En efecto, un argumento de recubrimiento estándar de la esfera (usando que $\xi\mapsto\xi^\top C_t\xi$ es Lipschitz
con constante $2\abs{C_t}$, que tiene todos los momentos) promueve esta cota a $\P(\inf_{\abs\xi=1}\xi^\top C_t\xi<
\varepsilon)\le C_p'\varepsilon^p$, y esto da $\E(\det C_t)^{-p}\le\E(\lambda_{\min}(C_t))^{-dp}<\infty$.

\paragraph{Paso 1: la dinámica de $Z_tU(X_t)$.} Para un campo $U$ suave, por la regla de Stratonovich,
\begin{equation}\label{eq:dinamica-ZU}
  \dd\big(Z_tU(X_t)\big)=Z_t[V_0,U](X_t)\,\dd t+\sum_{k=1}^mZ_t[V_k,U](X_t)\circ\dd W^k_t.
\end{equation}
En efecto, $\dd(U(X_t))=\nabla U(X_t)\circ\dd X_t=\nabla U\,V_0\,\dd t+\sum_k\nabla U\,V_k\circ\dd W^k$, y $\dd Z_t\,U=-Z\nabla V_0\,U\,\dd t-
\sum Z\nabla V_k\,U\circ\dd W^k$; sumando, aparecen $\nabla U\,V_j-\nabla V_j\,U=[V_j,U]$.

\paragraph{Paso 2: el lema de Norris.} Es la herramienta probabilística clave: una semimartingala no
puede ser pequeña sin que lo sean su parte de variación finita y su parte martingala.

\begin{lema}[Norris]\label{lem:norris}
Sean $a,y$ procesos de Itô reales, $u=(u^k)$ adaptado, con
\[
  Y_t=y+\int_0^ta_s\,\dd s+\sum_k\int_0^tu^k_s\,\dd W^k_s,\qquad a_t=\alpha+\int_0^t\beta_s\,\dd s+\sum_k\int_0^t\gamma^k_s\,\dd W^k_s.
\]
Supongamos que $\E\sup_{t\le T}(\abs{\beta_t}+\abs{\gamma_t}+\abs{a_t}+\abs{u_t})^p<\infty$ para todo $p$. Entonces, para todo $q>8$
y $r<\frac{q-8}{18}$, existe $C$ tal que, para $\varepsilon\in(0,1]$,
\[
  \P\Big(\int_0^TY_t^2\,\dd t<\varepsilon^q,\ \int_0^T\big(a_t^2+\abs{u_t}^2\big)\,\dd t\ge\varepsilon\Big)\le C\varepsilon^r\quad\text{(con }C=C_{q,r}).
\]
\end{lema}

\begin{proof}[Idea de la prueba]
Se usa repetidamente la \emph{desigualdad exponencial de martingalas}: si $M$ es
martingala continua, $\P(\sup_{t\le T}\abs{M_t}\ge\delta,\ \langle M\rangle_T<\rho)\le2e^{-\delta^2/2\rho}$. (1) Si $\int Y^2$ es pequeño,
entonces $\int Y\,\dd Y$, que es el término $\frac12(Y_T^2-Y_0^2)$ salvo la variación cuadrática, controla $\int\abs{u}^2$: por
Itô, $Y_T^2=y^2+2\int Ya\,\dd t+2\int Yu\,\dd W+\int\abs{u}^2\,\dd t$; la integral estocástica $\int Yu\,\dd W$ tiene variación
cuadrática $\int Y^2\abs{u}^2$, pequeña, luego es pequeña salvo un suceso de probabilidad $O(e^{-1/\varepsilon^c})$; y
$\int Ya$ es pequeño por Cauchy--Schwarz. Así $\int\abs{u}^2$ es pequeño. (2) Si $Y$ es pequeño en $L^2$ y $u$
también, entonces la parte de variación acotada $\int a$ es pequeña en sup (interpolación entre $L^2$ y
la regularidad Hölder que da la cota de $\beta,\gamma$), y aplicando el mismo argumento a $a$ en lugar de
$Y$, $\int a^2$ es pequeño. Ambos pasos pierden una potencia de $\varepsilon$, lo que explica las restricciones
$q>8$, $r<(q-8)/18$. Véase~\cite{Norris86} o~\cite[Lemma~2.3.2]{Nualart06}.
\end{proof}

\paragraph{Paso 3: iteración por corchetes.} Fijamos $\xi\in S^{d-1}$ y, para un campo $U$, escribimos
$Y^U_t=\xi^\top Z_tU(X_t)$. Entonces
\[
  \xi^\top C_t\xi=\sum_{k=1}^m\int_0^t\big(Y^{V_k}_s\big)^2\,\dd s.
\]
Por~\eqref{eq:dinamica-ZU}, en forma de Itô, $Y^U$ es un proceso de Itô con deriva $\xi^\top Z([V_0,U]+
\frac12\sum_k[V_k,[V_k,U]])(X)$ y coeficientes de martingala $Y^{[V_k,U]}$. Aplicando el Lema~\ref{lem:norris} con
$Y=Y^U$: si $\int_0^t(Y^U)^2<\varepsilon^q$, entonces, salvo un suceso de probabilidad $O(\varepsilon^p)$,
\[
  \int_0^t\big(Y^{[V_k,U]}_s\big)^2\,\dd s<\varepsilon,\quad k=1,\dots,m,\qquad\int_0^t\Big(Y_s^{[V_0,U]+\frac12\sum_k[V_k,[V_k,U]]}\Big)^2\,\dd s<\varepsilon.
\]
Combinando estas dos (el término $[V_k,[V_k,U]]$ ya es pequeño por la aplicación previa del lema
a $[V_k,U]$), $\int(Y^{[V_0,U]})^2$ es también pequeño, con una pérdida de exponente.

Partiendo de $U\in\mathcal V_0$ y repitiendo $N$ veces, obtenemos: existen $q_N$ y, para todo $p$, $C$ tales
que
\[
  \P\Big(\xi^\top C_t\xi<\varepsilon^{q_N},\ \sum_{U\in\mathcal V_N}\int_0^t\big(Y^U_s\big)^2\,\dd s\ge\varepsilon\Big)\le C\varepsilon^p.
\]

\paragraph{Paso 4: conclusión.} Por la condición de Hörmander en $x$ y la compacidad de la esfera,
existe $N$ y $c>0$ tales que $\sum_{U\in\mathcal V_N}(\xi^\top U(x))^2\ge c$ para todo $\abs\xi=1$. Como $Z_0=I$, $X_0=x$ y
las trayectorias son continuas, el suceso
\[
  \Big\{\sum_{U\in\mathcal V_N}\int_0^t\big(Y^U_s\big)^2\,\dd s<\varepsilon\Big\}
\]
implica que $s\mapsto\sum_U(Y^U_s)^2$, que vale al menos $c$ en $s=0$, debe caer por debajo de $c/2$ antes
del tiempo $2\varepsilon/c$. La probabilidad de esto es $\le\P(\sup_{s\le2\varepsilon/c}\abs{Z_sU(X_s)-U(x)}\ge c')\le C_p\varepsilon^p$
por las cotas de momentos de los incrementos (Kolmogorov). Combinando con el Paso~3,
$\P(\xi^\top C_t\xi<\varepsilon^{q_N})\le C\varepsilon^p$, que es~\eqref{eq:hormander-objetivo} tras reparametrizar $\varepsilon$.
\qed

\begin{corolario}[Hipoelipticidad]
Bajo la condición de Hörmander en todo punto, el operador
$\partial_t-\mathcal L$, con $\mathcal L=V_0+\frac12\sum V_k^2$, es hipoelíptico en $(0,\infty)\times\R^d$: su núcleo de calor $p_t(x,y)$ es $C^\infty$
en $(t,x,y)$.
\end{corolario}

\begin{proof}[Esbozo]
La regularidad en $y$ es el Teorema~\ref{teo:hormander}, con cotas uniformes en $x$ en compactos. La
regularidad en $x$ sale de diferenciar $\E[\varphi(X_t^x)]$ en $x$, usando el flujo $Y_t$ y las mismas integraciones
por partes, y en $t$, de la ecuación de Kolmogorov. El argumento de Hörmander para $\partial_t-\mathcal L$
como operador en $(t,x)$ requiere además aplicar el resultado al adjunto.
\end{proof}

\section{Comentarios y extensiones}
\begin{itemize}
  \item \textbf{Cotas cuantitativas.} Kusuoka y Stroock~\cite{KusuokaStroock} obtuvieron cotas gaussianas para $p_t(x,y)$
  y sus derivadas, con escala $t^{-\text{grado}/2}$ en cada dirección según el número de corchetes
  necesarios para generarla, en consonancia con el ejemplo de Kolmogorov.
  \item \textbf{Enfoque de Bismut.} Bismut~\cite{Bismut81} dio una prueba alternativa basada en Girsanov:
  perturbar el ruido en direcciones de Cameron--Martin y diferenciar la fórmula de Girsanov
  produce directamente las fórmulas de integración por partes. Es la base de la fórmula de
  Bismut--Elworthy--Li, que usaremos en finanzas.
  \item \textbf{Dimensión infinita.} Hairer y Mattingly~\cite{HairerMattingly11} extendieron el programa a EDPEs semilineales
  (Navier--Stokes 2D con ruido degenerado), probando ergodicidad única. Ahí la
  matriz de Malliavin es un operador en dimensión infinita y nunca es invertible; se prueba
  una no degeneración en un subespacio de dimensión finita.
  \item \textbf{Rough paths.} Cass y Friz~\cite{CassFriz10} probaron que la condición de Hörmander también da
  densidades para EDRs dirigidas por procesos gaussianos rugosos (browniano fraccionario
  con $H>1/4$); véase el Capítulo~\ref{cap:rough-paths}.
\end{itemize}


\part{Aplicaciones a las finanzas matemáticas}
\chapter{Griegas por integración por partes}\label{cap:griegas}

\begin{panorama}
Las \emph{griegas} son las derivadas del precio de un derivado respecto a sus parámetros:
$\Delta=\partial_{S_0}V$, $\Gamma=\partial^2_{S_0}V$, $\mathcal V=\partial_\sigma V$ (vega), etc. Cuando el pago es discontinuo (digitales,
barreras) o depende de la trayectoria (asiáticas, lookback), las técnicas Monte Carlo
ingenuas fallan o son muy ruidosas. El trabajo de Fournié, Lasry, Lebuchoux, Lions y
Touzi~\cite{Fournie99} (1999) mostró que el cálculo de Malliavin produce estimadores de la forma
$\E[\Phi(\cdot)\,\pi]$, con un \emph{peso} $\pi$ que no depende del pago. Desarrollamos la teoría, calculamos
los pesos clásicos, discutimos su optimalidad en varianza y la localización, e ilustramos
numéricamente.
\end{panorama}

\section{Planteamiento: tres métodos para una derivada}
Sea $X^\theta$ un vector aleatorio que depende de un parámetro $\theta$ y $V(\theta)=\E[\Phi(X^\theta)]$. Hay tres formas
clásicas de estimar $V'(\theta)$ por Monte Carlo:
\begin{enumerate}
  \item \textbf{Diferencias finitas:} $\frac1{2h}\E[\Phi(X^{\theta+h})-\Phi(X^{\theta-h})]$. Sesgo $O(h^2)$ (si $V$ es regular) y varianza
  $O(h^{-1})$ para $\Phi$ discontinua aun con números aleatorios comunes. Hay que equilibrar
  ambas.
  \item \textbf{Derivada trayectorial:} $\E[\nabla\Phi(X^\theta)\,\partial_\theta X^\theta]$. Insesgado y de baja varianza, pero exige $\Phi$
  (Lipschitz) diferenciable: no sirve para digitales.
  \item \textbf{Razón de verosimilitud:} si $X^\theta$ tiene densidad $p_\theta$ conocida, $V'(\theta)=\E[\Phi(X^\theta)\,\partial_\theta\log p_\theta(X^\theta)]$.
  No requiere regularidad de $\Phi$, pero sí la densidad explícita, que rara vez se conoce.
\end{enumerate}
El método de Malliavin combina las ventajas de 2 y 3: parte de la derivada trayectorial y
\emph{transfiere} la derivada de $\Phi$ a un peso mediante integración por partes en el espacio de Wiener,
\emph{sin conocer la densidad}.

\section{El principio general}

\begin{proposicion}[Pesos de Malliavin]\label{prop:pesos}
Sea $F\in\D^{1,2}$ real, $G\in\D^{1,2}$, y $a\in L^2(\Omega;\Hf)$ un
proceso tal que $\ip{DF}{a}_\Hf\ne0$ c.s.\ y $u=Ga/\ip{DF}{a}\in\Dom\delta$. Entonces, para toda $\Phi$ medible
con crecimiento polinomial,
\[
  \E\big[\Phi'(F)\,G\big]=\E\big[\Phi(F)\,\delta(u)\big],
\]
donde el lado izquierdo se interpreta, para $\Phi$ no derivable, como $\lim_n\E[\Phi_n'(F)G]$ con $\Phi_n$ suaves
que aproximan $\Phi$ (por ejemplo, en $L^2$ de la ley de $F$, que tiene densidad bajo estas hipótesis).
\end{proposicion}

\begin{proof}
Para $\Phi\in C_b^1$: $\ip{D\Phi(F)}{a}=\Phi'(F)\ip{DF}{a}$, luego $\Phi'(F)G=\ip{D\Phi(F)}{u}$ y basta
aplicar la dualidad. Para $\Phi$ general, $\E[\Phi_n(F)\delta(u)]\to\E[\Phi(F)\delta(u)]$ si $\Phi_n\to\Phi$ en $L^2(\P\circ F^{-1})$ y
$\delta(u)\in L^2$.
\end{proof}

\begin{intuicion}
El proceso $a$ es una elección libre: es la \emph{dirección} en el espacio de Cameron--Martin en
la que perturbamos el ruido para reproducir el efecto de mover $\theta$. Diferentes elecciones
dan diferentes pesos, todos insesgados pero con distinta varianza (§\ref{sec:pesos-optimos}). La elección
canónica en el espacio de Wiener es $a\equiv1$ en $[0,T]$, pero en problemas de trayectoria
(asiáticas) conviene $a_t=S_t$.
\end{intuicion}

\section{Black--Scholes: pesos explícitos}
Bajo la medida de riesgo neutro, $S_T=S_0\exp\big((r-\frac{\sigma^2}{2})T+\sigma W_T\big)$ y $V=e^{-rT}\E[\Phi(S_T)]$. Se tiene
$D_tS_T=\sigma S_T\ind_{[0,T]}(t)$.

\begin{proposicion}\label{prop:pesos-bs}
Para toda $\Phi$ con crecimiento polinomial,
\begin{align*}
  \Delta&=\frac{e^{-rT}}{S_0\sigma T}\,\E\big[\Phi(S_T)\,W_T\big],\\
  \Gamma&=\frac{e^{-rT}}{S_0^2\sigma T}\,\E\Big[\Phi(S_T)\Big(\frac{W_T^2}{\sigma T}-W_T-\frac1\sigma\Big)\Big],\\
  \mathcal V&=e^{-rT}\,\E\Big[\Phi(S_T)\Big(\frac{W_T^2}{\sigma T}-W_T-\frac1\sigma\Big)\Big].
\end{align*}
\end{proposicion}

\begin{proof}
\emph{Delta.} $\partial_{S_0}\E\Phi(S_T)=\E[\Phi'(S_T)S_T/S_0]$. Con $F=S_T$, $G=S_T/S_0$, $a\equiv1$:
$\ip{DF}{a}=\sigma TS_T$ y $u=\frac1{S_0\sigma T}$, constante. Así $\delta(u)=\frac{W_T}{S_0\sigma T}$.

\emph{Vega.} $\partial_\sigma S_T=S_T(W_T-\sigma T)$, luego $\partial_\sigma\E\Phi(S_T)=\E[\Phi'(S_T)S_T(W_T-\sigma T)]$. Con $a\equiv1$,
$u=\frac{W_T-\sigma T}{\sigma T}$ y, por la Proposición~\ref{prop:reglas-delta}(ii) con $D_tW_T=1$,
\[
  \delta(u)=\frac{(W_T-\sigma T)W_T}{\sigma T}-\frac{T}{\sigma T}=\frac{W_T^2}{\sigma T}-W_T-\frac1\sigma.
\]

\emph{Gamma.} $\partial^2_{S_0}\E\Phi(S_T)=\frac1{S_0^2}\E[\Phi''(S_T)S_T^2]$. Integramos por partes una vez con $G=S_T^2$: $u=\frac{S_T}{\sigma T}$,
$\delta(u)=\frac{S_TW_T}{\sigma T}-\frac{\sigma TS_T}{\sigma T}=S_T(\frac{W_T}{\sigma T}-1)$. Así $\E[\Phi''(S_T)S_T^2]=\E[\Phi'(S_T)S_T(\frac{W_T}{\sigma T}-1)]$. Una segunda
integración por partes con $G=S_T(\frac{W_T}{\sigma T}-1)$, $u=\frac1{\sigma T}(\frac{W_T}{\sigma T}-1)$, da $\delta(u)=\frac1{\sigma T}\big(\frac{W_T^2}{\sigma T}-W_T-\frac1\sigma\big)$.
\end{proof}

Observemos la relación clásica $\mathcal V=\sigma TS_0^2\Gamma$, que se lee directamente en los pesos.

\section{Difusiones generales: la fórmula de Bismut--Elworthy--Li}
Sea $X^x$ la solución de $\dd X=b(X)\,\dd t+\sigma(X)\,\dd W$ en $\R^d$ con $\sigma$ invertible y $\sigma^{-1}$ acotada,
$Y_t=\nabla_xX_t$.

\begin{teorema}[Bismut--Elworthy--Li]\label{teo:bel}
Para toda $\Phi$ medible acotada, $x\mapsto\E[\Phi(X_T^x)]$ es
diferenciable y
\[
  \nabla_x\E\big[\Phi(X_T^x)\big]=\E\Big[\Phi(X_T^x)\,\frac1T\int_0^T\big(\sigma^{-1}(X_t)Y_t\big)^\top\dd W_t\Big].
\]
Más en general, puede sustituirse $\frac1T\ind_{[0,T]}(t)$ por cualquier $\alpha(t)$ determinista con $\int_0^T\alpha=1$.
\end{teorema}

\begin{proof}
Para $\Phi\in C_b^1$, $\nabla_x\E\Phi(X_T)=\E[\nabla\Phi(X_T)Y_T]$. Por el Teorema~\ref{teo:derivada-ede}, $D_tX_T=
Y_TZ_t\sigma(X_t)$, de modo que $Y_T=D_tX_T\,\sigma^{-1}(X_t)Y_t$ para todo $t\le T$. Promediando en $t$ con peso
$\alpha$:
\[
  Y_T=\int_0^TD_tX_T\,\sigma^{-1}(X_t)Y_t\,\alpha(t)\,\dd t.
\]
Por tanto, para cada vector $v$,
\[
  \E[\nabla\Phi(X_T)Y_Tv]=\E\int_0^TD_t\big(\Phi(X_T)\big)\,\sigma^{-1}(X_t)Y_tv\,\alpha(t)\,\dd t=\E\Big[\Phi(X_T)\,\delta\big(\sigma^{-1}(X_\cdot)Y_\cdot v\,\alpha\big)\Big].
\]
El integrando es adaptado, luego $\delta$ es la integral de Itô. Por densidad, la fórmula vale para $\Phi$
acotada medible.
\end{proof}

\begin{observacion}
La fórmula da, como subproducto, una estimación de regularización:
$\abs{\nabla P_T\Phi(x)}\le\frac{C}{\sqrt T}\norm{\Phi}_\infty$. Es una herramienta estándar en el estudio de EDPs parabólicas y de
ergodicidad.
\end{observacion}

\section{Opciones dependientes de la trayectoria: la asiática}
Sea $A=\int_0^TS_t\,\dd t$ y $V=e^{-rT}\E[\Phi(A/T)]$ en Black--Scholes.

\begin{proposicion}\label{prop:asiatica}
\[
  \Delta_{\text{asiática}}=\frac{e^{-rT}}{S_0}\,\E\Big[\Phi\Big(\frac AT\Big)\Big(\frac{2(S_T-S_0)}{\sigma^2A}-\frac{2r}{\sigma^2}+1\Big)\Big].
\]
\end{proposicion}

\begin{proof}
$\partial_{S_0}A=A/S_0$. Elegimos $a_t=S_t$. Como $D_tA=\sigma\int_t^TS_u\,\dd u$,
\[
  \ip{DA}{a}=\sigma\int_0^TS_t\int_t^TS_u\,\dd u\,\dd t=\frac\sigma2A^2.
\]
Así $u_t=\frac{A}{S_0}\frac{S_t}{\sigma A^2/2}=\frac{2}{\sigma S_0A}S_t$. Por la Proposición~\ref{prop:reglas-delta}(ii), con $F=\frac{2}{\sigma S_0A}$ y $\ip{DF}{S}=
-\frac{2}{\sigma S_0A^2}\ip{DA}{S}=-\frac1{S_0}$,
\[
  \delta(u)=\frac{2}{\sigma S_0A}\int_0^TS_t\,\dd W_t+\frac1{S_0}.
\]
Finalmente, $\int_0^T\sigma S_t\,\dd W_t=S_T-S_0-rA$, de donde se obtiene la fórmula.
\end{proof}

La elección $a=S$ en lugar de $a=1$ es la que hace que $\ip{DA}{a}$ sea un cuadrado perfecto y
el peso sea explícito.

\section{Pesos óptimos y localización}\label{sec:pesos-optimos}

\begin{proposicion}[Peso de varianza mínima]
Sea $F$ real, y sea $\mathcal W$ el conjunto de pesos $\pi\in L^2$
tales que $\E[\Phi'(F)G]=\E[\Phi(F)\pi]$ para toda $\Phi\in C_b^1$. Entonces $\mathcal W=\pi_0+\{\pi:\ \cond{\pi}{F}=0\}$
para cualquier $\pi_0\in\mathcal W$, y el peso de varianza mínima, simultáneamente para todas las $\Phi$, es
$\pi^*=\cond{\pi_0}{F}$.
\end{proposicion}

\begin{proof}
$\E[\Phi(F)(\pi-\pi_0)]=0$ para toda $\Phi$ equivale a $\cond{\pi-\pi_0}{F}=0$. Para
$\pi\in\mathcal W$, $\E[\Phi(F)^2\pi^2]=\E[\Phi(F)^2\pi^{*2}]+\E[\Phi(F)^2(\pi-\pi^*)^2]$, porque el término cruzado es
$\E[\Phi(F)^2\pi^*\cond{\pi-\pi^*}{F}]=0$.
\end{proof}

En Black--Scholes, el peso $W_T/(S_0\sigma T)$ es función de $S_T$, luego ya es óptimo: coincide con el
peso de razón de verosimilitud $\partial_{S_0}\log p(S_T)$. En problemas multidimensionales o dependientes
de la trayectoria, $\cond{\pi_0}{F}$ no es calculable, y la elección de $a$ importa.

\paragraph{Localización.} El peso tiene varianza grande cuando $\Phi(F)$ es grande lejos de las discontinuidades.
La idea de~\cite{Fournie99,Fournie01} es separar: $\Phi=\Phi_{\mathrm{reg}}+\Phi_{\mathrm{sing}}$, donde $\Phi_{\mathrm{sing}}$ tiene soporte compacto
cerca de la singularidad, y usar la derivada trayectorial para la parte regular y el peso solo para
la singular:
\[
  \Delta=\E\big[\Phi_{\mathrm{reg}}'(F)G\big]+\E\big[\Phi_{\mathrm{sing}}(F)\pi\big].
\]

\section{Ilustración numérica}
Calculamos la delta de una opción digital $\Phi(S)=\ind_{\{S>K\}}$ en Black--Scholes con $S_0=K=100$,
$r=0{,}05$, $\sigma=0{,}2$, $T=1$. El valor exacto es $\Delta=e^{-rT}\varphi(d_2)/(S_0\sigma\sqrt T)=0{,}018762$. Cada
estimador usa $N=10^5$ trayectorias; repetimos 50 veces para medir la desviación típica. Las
diferencias finitas usan números aleatorios comunes. La versión localizada usa $\Phi_{\mathrm{reg}}(S)=
(1+e^{-(S-K)/5})^{-1}$.

\begin{center}
\begin{tabular}{lcc}
\toprule
Estimador & Media & Desv.\ típica\\
\midrule
Diferencias finitas, $h=1$ & 0,018803 & $2{,}8\times10^{-4}$\\
Diferencias finitas, $h=0{,}1$ & 0,018798 & $9{,}6\times10^{-4}$\\
Malliavin, peso $W_T/(S_0\sigma T)$ & 0,018758 & $9{,}0\times10^{-5}$\\
Malliavin localizado & 0,018771 & $5{,}4\times10^{-5}$\\
\bottomrule
\end{tabular}
\end{center}

El peso de Malliavin reduce la desviación típica en un factor 3 frente a la mejor diferencia
finita, y la localización la reduce otro 40\%; las diferencias entre las medias de los cuatro
estimadores están dentro del error de Monte Carlo. Para griegas de segundo orden (gamma de
digitales), la ventaja es mucho mayor: las diferencias finitas centradas tienen varianza $O(h^{-3})$.

\section{Extensiones}
\begin{itemize}
  \item \textbf{Volatilidad estocástica.} En Heston o SABR la matriz de Malliavin del par $(S_T,v_T)$
  no es explícita, pero el principio de la Proposición~\ref{prop:pesos} con $a$ elegido en la dirección del
  browniano del subyacente da pesos calculables por simulación de $Y_t$ (\cite{Fournie99,MalliavinThalmaier06}).
  \item \textbf{Saltos.} En modelos de Lévy se usa el cálculo de Malliavin en el espacio de Poisson o se
  integra por partes solo respecto a la componente browniana (\cite{DiNunnoOksendalProske09}).
  \item \textbf{Barreras y lookback.} Las discontinuidades en la trayectoria exigen localizar alrededor
  del tiempo de toque; la Proposición~\ref{prop:supremo} da la derivada del supremo, y se combinan pesos
  con puentes brownianos.
  \item \textbf{Diferenciación automática adjunta (AAD).} En la práctica bancaria actual, las
  griegas de pagos regulares se calculan con AAD; los pesos de Malliavin siguen siendo
  la herramienta de referencia para pagos discontinuos y se combinan con AAD mediante
  suavizado o localización.
\end{itemize}

\chapter{Cobertura y Clark--Ocone en finanzas}\label{cap:cobertura}

\begin{panorama}
En un mercado completo, todo derivado se replica con una cartera autofinanciada. El
teorema de representación de martingalas garantiza su existencia; la fórmula de Clark--Ocone
la \emph{calcula}. Desarrollamos el marco, obtenemos las carteras de cobertura de opciones
europeas, digitales, asiáticas y lookback, y resolvemos con la fórmula generalizada de
Ocone--Karatzas el problema de inversión óptima de Merton con coeficientes aleatorios,
donde aparece la \emph{demanda de cobertura intertemporal} como derivada de Malliavin del
precio de mercado del riesgo.
\end{panorama}

\section{El marco de mercado}
Sea $(\Omega,\F,(\F_t),\P)$ con $\F_t$ la filtración de un browniano $W$. El mercado tiene un activo sin
riesgo $B_t=e^{\int_0^tr_s\,\dd s}$ y un activo con riesgo
\[
  \dd S_t=S_t(\mu_t\,\dd t+\sigma_t\,\dd W_t),
\]
con $r,\mu,\sigma$ progresivos y acotados, $\sigma_t\ge\sigma_0>0$. El \emph{precio de mercado del riesgo} es $\theta_t=
(\mu_t-r_t)/\sigma_t$. Por Girsanov (Teorema~\ref{teo:girsanov}), bajo $\dd\Q=\mathcal E(-\theta)_T\,\dd\P$ el proceso $\tilde W_t=W_t+\int_0^t\theta_s\,\dd s$
es browniano y los precios descontados $\tilde S=S/B$ son martingalas: $\dd\tilde S_t=\sigma_t\tilde S_t\,\dd\tilde W_t$.

Una \emph{cartera autofinanciada} es un proceso $\phi_t$ (número de acciones) tal que el valor $V_t=
\phi_tS_t+\psi_tB_t$ satisface $\dd V_t=\phi_t\,\dd S_t+\psi_t\,\dd B_t$. Equivalentemente, el valor descontado $\tilde V=V/B$
satisface
\[
  \dd\tilde V_t=\phi_t\,\dd\tilde S_t=\phi_t\sigma_t\tilde S_t\,\dd\tilde W_t.
\]

\begin{teorema}[Cobertura de Clark--Ocone]\label{teo:cobertura}
Sea $H\ge0$ un pago $\F_T$-medible con $\tilde H=
H/B_T\in L^2(\Q)$.
\begin{enumerate}[label=(\roman*)]
  \item \textup{(Coeficientes deterministas.)} Si $r,\sigma$ son deterministas, la filtración de $\tilde W$ coincide con
  $(\F_t)$ y, si $\tilde H\in\D^{1,2}$ respecto a $\tilde W$, la cartera
  \[
    \phi_t=\frac{\E_\Q\big[\tilde D_t\tilde H\,\big|\,\F_t\big]}{\sigma_t\tilde S_t}
  \]
  replica $H$, donde $\tilde D$ es la derivada de Malliavin respecto a $\tilde W$. El precio es $V_0=\E_\Q\tilde H$.
  \item \textup{(Coeficientes aleatorios.)} En general, con $D$ la derivada respecto a $W$,
  \[
    \phi_t=\frac1{\sigma_t\tilde S_t}\,\E_\Q\Big[D_t\tilde H-\tilde H\int_t^TD_t\theta_s\,\dd\tilde W_s\,\Big|\,\F_t\Big].
  \]
\end{enumerate}
\end{teorema}

\begin{proof}
(i) Por Clark--Ocone (Teorema~\ref{teo:clark-ocone}) bajo $\Q$ para el browniano $\tilde W$: $\tilde H=\E_\Q\tilde H+
\int_0^T\E_\Q[\tilde D_t\tilde H\mid\F_t]\,\dd\tilde W_t$. Comparando con $\tilde V_T=V_0+\int_0^T\phi_t\sigma_t\tilde S_t\,\dd\tilde W_t$ se obtiene $\phi$. (ii) Es el
Teorema~\ref{teo:clark-ocone-generalizado}.
\end{proof}

\begin{intuicion}
La cobertura es \emph{la sensibilidad esperada del pago a un impulso del ruido}, dividida por la
sensibilidad del activo al mismo impulso ($\sigma_t\tilde S_t$). Es una ``regla de tres'' en el espacio de
trayectorias. Cuando el pago es $\Phi(S_T)$ esto se reduce al delta de la EDP de Black--Scholes,
pero la fórmula funciona sin EDP para pagos que dependen de toda la trayectoria.
\end{intuicion}

\section{Ejemplos en Black--Scholes}
Tomamos $r,\sigma$ constantes. Entonces $\tilde D_tS_T=\sigma S_T$ para $t\le T$, y escribimos $\tau=T-t$.

\begin{ejemplo}[Call europea]
$H=(S_T-K)^+$. Por la regla de la cadena (Lipschitz),
$\tilde D_tH=\ind_{\{S_T>K\}}\sigma S_T$. Así
\[
  \phi_t=\frac{e^{-rT}\sigma\E_\Q[S_T\ind_{\{S_T>K\}}\mid\F_t]}{\sigma e^{-rt}S_t}=\frac{e^{-r\tau}}{S_t}\E_\Q[S_T\ind_{\{S_T>K\}}\mid\F_t]=N(d_1),
\]
con $d_1=\frac{\log(S_t/K)+(r+\sigma^2/2)\tau}{\sigma\sqrt\tau}$, usando que bajo la medida $\dd\Q^S/\dd\Q=S_Te^{-rT}/S_0$ el suceso tiene
probabilidad $N(d_1)$. Es la delta clásica.
\end{ejemplo}

\begin{ejemplo}[Digital]
$H=\ind_{\{S_T>K\}}\notin\D^{1,2}$. Sin embargo, aproximando por $H_n=\Phi_n(S_T)$
suaves, las carteras $\phi^n_t=e^{-r\tau}\E_\Q[\Phi_n'(S_T)S_T\mid\F_t]/S_t$ convergen para $t<T$, porque la ley
condicional de $S_T$ dado $\F_t$ tiene densidad regular. El límite es
\[
  \phi_t=\frac{e^{-r\tau}\,\varphi(d_2)}{\sigma S_t\sqrt\tau},\qquad t<T.
\]
Formalmente, $\tilde D_t\ind_{\{S_T>K\}}=\delta_K(S_T)\sigma S_T$ es una \emph{distribución de Watanabe}: el cálculo de Malliavin
se extiende a $\D^{-\infty}$, el dual de $\D^\infty$, y Clark--Ocone vale ahí para $t<T$.
\end{ejemplo}

\begin{ejemplo}[Asiática]
$H=(\bar A-K)^+$, con $\bar A=\frac1T\int_0^TS_u\,\dd u$. Como $\tilde D_t\bar A=\frac\sigma T\int_t^TS_u\,\dd u$,
\[
  \phi_t=\frac{e^{-r\tau}}{S_t}\,\E_\Q\Big[\ind_{\{\bar A>K\}}\frac1T\int_t^TS_u\,\dd u\,\Big|\,\F_t\Big].
\]
Esta expresión se evalúa por Monte Carlo anidado o, mejor, se reduce a una función de
$(t,S_t,\int_0^tS)$ por la propiedad de Markov del par.
\end{ejemplo}

\begin{ejemplo}[Lookback flotante]
$H=M_T-S_T$, con $M_T=\max_{u\le T}S_u$. Sea $\tau^*$ el instante
del máximo. Por la Proposición~\ref{prop:supremo}, $\tilde D_tM_T=\sigma M_T\ind_{\{t\le\tau^*\}}$, luego
\[
  \phi_t=\frac{e^{-r\tau}}{S_t}\Big(\E_\Q\big[M_T\ind_{\{\tau^*\ge t\}}\mid\F_t\big]-\E_\Q[S_T\mid\F_t]\Big).
\]
El primer término es la parte de la cartera que cubre la posibilidad de que el máximo aún no se
haya alcanzado; se calcula con el principio de reflexión para el browniano con deriva.
\end{ejemplo}

\section{Inversión óptima: Merton con coeficientes aleatorios}
Un inversor con riqueza inicial $x$ elige la fracción $\pi_t$ de su riqueza invertida en $S$ para maximizar
$\E[U(X_T)]$, con $U(x)=\frac{x^{1-\gamma}}{1-\gamma}$ (aversión relativa al riesgo $\gamma>0$; $\gamma=1$ es log). Para simplificar,
$r\equiv0$. La riqueza satisface $\dd X_t=\pi_tX_t\sigma_t(\dd W_t+\theta_t\,\dd t)=\pi_tX_t\sigma_t\,\dd\tilde W_t$.

\paragraph{Método de martingala.} En mercado completo, el problema dinámico equivale al estático
\[
  \max\ \E[U(X_T)]\quad\text{s.a.}\quad\E[\xi_TX_T]=x,\qquad\xi_T=\mathcal E(-\theta)_T=\frac{\dd\Q}{\dd\P}.
\]
La condición de primer orden $U'(X_T)=y\xi_T$ da $X_T^*=(y\xi_T)^{-1/\gamma}$, con $y$ fijado por la restricción
presupuestaria. Queda por \emph{encontrar la estrategia} que replica $X_T^*$: es un problema de cobertura.

\begin{teorema}[Ocone--Karatzas]\label{teo:ocone-karatzas}
Bajo condiciones de integrabilidad y $\theta_s\in\D^{1,2}$, la cartera
óptima es
\[
  \pi_t^*=\underbrace{\frac{\theta_t}{\gamma\sigma_t}}_{\text{miope}}\ +\ \underbrace{\Big(\frac1\gamma-1\Big)\frac1{\sigma_t}\frac{\E_\Q\big[X_T^*\int_t^TD_t\theta_s\,\dd\tilde W_s\mid\F_t\big]}{\E_\Q[X_T^*\mid\F_t]}}_{\text{demanda de cobertura intertemporal}}.
\]
\end{teorema}

\begin{proof}
La riqueza óptima es $X_t=\E_\Q[X_T^*\mid\F_t]$ y $\dd X_t=\pi_tX_t\sigma_t\,\dd\tilde W_t$. Por el Teorema~\ref{teo:clark-ocone-generalizado}
aplicado a $F=X_T^*$,
\[
  \pi_tX_t\sigma_t=\E_\Q\Big[D_tX_T^*-X_T^*\int_t^TD_t\theta_s\,\dd\tilde W_s\,\Big|\,\F_t\Big].
\]
Ahora $\log X_T^*=-\frac1\gamma\log(y\xi_T)$ y $\log\xi_T=-\int_0^T\theta\,\dd W-\frac12\int_0^T\theta^2\,\dd s$. Por el Corolario~\ref{cor:derivada-ito},
\[
  D_t\log\xi_T=-\theta_t-\int_t^TD_t\theta_s\,\dd W_s-\int_t^T\theta_sD_t\theta_s\,\dd s=-\theta_t-\int_t^TD_t\theta_s\,\dd\tilde W_s.
\]
Así $D_tX_T^*=\frac1\gamma X_T^*\big(\theta_t+\int_t^TD_t\theta_s\,\dd\tilde W_s\big)$. Sustituyendo,
\[
  \pi_tX_t\sigma_t=\frac{\theta_t}{\gamma}X_t+\Big(\frac1\gamma-1\Big)\E_\Q\Big[X_T^*\int_t^TD_t\theta_s\,\dd\tilde W_s\,\Big|\,\F_t\Big],
\]
y se divide por $X_t\sigma_t$.
\end{proof}

\begin{observacion}
\begin{itemize}
  \item Con utilidad logarítmica ($\gamma=1$) la demanda de cobertura desaparece:
  el inversor log es miope.
  \item Si $\theta$ es determinista, $D\theta=0$ y la cartera es miope para toda $\gamma$: la solución clásica de
  Merton $\pi=\frac{\mu-r}{\gamma\sigma^2}$.
  \item Si $\theta$ depende de un factor estocástico correlacionado con $W$ (predecibilidad de rentabilidades,
  volatilidad estocástica), $D_t\theta_s$ mide cómo un choque hoy cambia las oportunidades
  de inversión futuras, y el inversor con $\gamma>1$ se cubre contra su deterioro. La fórmula es
  la base de los métodos de simulación de Detemple, García y Rindisbacher para carteras
  óptimas en alta dimensión.
\end{itemize}
\end{observacion}

\section{Información privilegiada}
¿Cuánto vale saber algo sobre el futuro? Supongamos que un inversor conoce desde $t=0$ una
variable $G$ (por ejemplo, $G=W_T$ o $G=\ind_{\{S_T>K\}}$). Su filtración es $\G_t=\F_t\vee\sigma(G)$. Bajo la
\emph{hipótesis de Jacod} (la ley condicional de $G$ dado $\F_t$ tiene densidad $p_t(\cdot)$ respecto a la ley de $G$),
$W$ es una $\G$-semimartingala:
\[
  W_t=\hat W_t+\int_0^t\frac{\dd\langle p_\cdot(x),W\rangle_s}{p_s(x)}\Big|_{x=G}\qquad\text{con }\hat W\text{ un }\G\text{-browniano}.
\]
Imkeller~\cite{Imkeller03} observó que el término de deriva, la ``ventaja informativa'' $\mu^G_s$, se expresa con
la derivada de Malliavin de la densidad condicional: $\dd\langle p_\cdot(x),W\rangle_s=\cond{D_s\,\delta_x(G)}{\F_s}\,\dd s$ por
Clark--Ocone. La utilidad logarítmica adicional del insider es
\[
  \E\big[\log X^\G_T\big]-\E\big[\log X^\F_T\big]=\frac12\E\int_0^T(\mu^G_s)^2\,\dd s,
\]
que coincide con la información mutua entre $G$ y la trayectoria $W$. Si $G=W_T$, $\mu^G_s=\frac{W_T-W_s}{T-s}$ y
la utilidad adicional es infinita: hay arbitraje. El enfoque anticipativo alternativo usa integrales
de Skorokhod o de tipo ``forward'' para definir la riqueza del insider sin pasar por $\G$.

\chapter{Esperanzas condicionales, opciones americanas y volatilidad estocástica}\label{cap:esperanzas-condicionales}

\begin{panorama}
Tres aplicaciones financieras más avanzadas. (1) Una representación de las esperanzas
condicionales $\E[\Phi(G)\mid F=x]$ como cociente de esperanzas \emph{no condicionales}, que permite
valorar opciones americanas por Monte Carlo (Fournié et al., Lions--Regnier, Bally--Caramellino--Zanette).
(2) La fórmula de descomposición de Alòs, que generaliza la
fórmula de Hull--White a volatilidad correlacionada mediante el cálculo anticipativo. (3)
Su consecuencia más célebre: el comportamiento en tiempo corto del \emph{skew} de volatilidad
implícita, que conecta con la volatilidad rugosa.
\end{panorama}

\section{Esperanzas condicionales por integración por partes}
Formalmente, $\E[\Phi(G)\mid F=x]=\E[\Phi(G)\delta_x(F)]/\E[\delta_x(F)]$. El cálculo de Malliavin permite
reescribir $\E[\Phi(G)\delta_x(F)]$ sustituyendo la delta de Dirac por un indicador y un peso.

\begin{teorema}\label{teo:esp-cond-malliavin}
Sea $F\in\D^{1,2}$ real, $G\in\D^{1,2}$ (posiblemente vectorial), y $u\in L^2(\Omega;\Hf)$ tal que
\[
  \ip{DF}{u}_\Hf=1\quad\text{y}\quad\ip{DG}{u}_\Hf=0\qquad\text{c.s.},
\]
con $u\in\Dom\delta$. Entonces, para $\Phi$ con crecimiento polinomial y $x$ en el interior del soporte
de $F$ donde $p_F(x)>0$,
\[
  \E[\Phi(G)\mid F=x]=\frac{\E\big[\Phi(G)\,\ind_{\{F>x\}}\,\delta(u)\big]}{\E\big[\ind_{\{F>x\}}\,\delta(u)\big]}.
\]
\end{teorema}

\begin{proof}
Sea $\psi_\varepsilon$ una aproximación de $\delta_x$ y $\Psi_\varepsilon(y)=\int_{-\infty}^y\psi_\varepsilon$. Para $\Phi$ suave, la condición
$\ip{DG}{u}=0$ da $\ip{D(\Phi(G)\Psi_\varepsilon(F))}{u}=\Phi(G)\psi_\varepsilon(F)\ip{DF}{u}=\Phi(G)\psi_\varepsilon(F)$. Por dualidad,
\[
  \E[\Phi(G)\psi_\varepsilon(F)]=\E[\Phi(G)\Psi_\varepsilon(F)\delta(u)].
\]
Cuando $\varepsilon\to0$, el lado izquierdo tiende a $\E[\Phi(G)\mid F=x]p_F(x)$ y el derecho a
$\E[\Phi(G)\ind_{\{F>x\}}\delta(u)]$. Con $\Phi\equiv1$ se obtiene el denominador: $p_F(x)=\E[\ind_{\{F>x\}}\delta(u)]$
(Teorema~\ref{teo:densidad}). Se extiende a $\Phi$ general por densidad.
\end{proof}

\begin{ejemplo}[Black--Scholes]\label{ej:esp-cond-bs}
Sean $0<s<t$, $F=S_s$, $G=S_t$. Buscamos $u$ con $\ip{DS_t}{u}=0$
y $\ip{DS_s}{u}\ne0$. Como $D_rS_s=\sigma S_s\ind_{[0,s]}(r)$ y $D_rS_t=\sigma S_t\ind_{[0,t]}(r)$, la elección
\[
  a_r=\frac1s\ind_{[0,s]}(r)-\frac1{t-s}\ind_{(s,t]}(r)
\]
da $\ip{DS_t}{a}=0$ y $\ip{DS_s}{a}=\sigma S_s$. Con $u=a/(\sigma S_s)$ y la Proposición~\ref{prop:reglas-delta}(ii),
\[
  \delta(u)=\frac1{\sigma S_s}\Big(\frac{W_s}{s}-\frac{W_t-W_s}{t-s}\Big)+\frac1{S_s}=\frac1{\sigma S_s}\Big(\frac{tW_s-sW_t}{s(t-s)}+\sigma\Big).
\]
Por tanto
\[
  \E[\Phi(S_t)\mid S_s=x]=\frac{\E\Big[\Phi(S_t)\ind_{\{S_s>x\}}\frac1{S_s}\Big(\frac{tW_s-sW_t}{s(t-s)}+\sigma\Big)\Big]}{\E\Big[\ind_{\{S_s>x\}}\frac1{S_s}\Big(\frac{tW_s-sW_t}{s(t-s)}+\sigma\Big)\Big]}.
\]
\end{ejemplo}

\begin{intuicion}
La dirección $a$ perturba el ruido en $[0,s]$ hacia arriba y en $(s,t]$ hacia abajo de modo
que $S_t$ no se mueve: es una ``variación con extremo final fijo''. El peso resultante es un
estimador de $\delta_x(S_s)$ que no requiere que las trayectorias pasen exactamente por $x$, lo que
haría el estimador ingenuo inviable.
\end{intuicion}

\paragraph{Localización.} La varianza de estos estimadores es alta cuando $x$ está lejos de la masa de $F$.
Para cualquier densidad $\psi$ en $\R$ con primitiva $\Psi$, se tiene la identidad, válida para todo $\psi$,
\[
  \E[\Phi(G)\delta_x(F)]=\E\big[\Phi(G)\psi(F-x)\big]+\E\big[\Phi(G)\big(\ind_{\{F>x\}}-\Psi(F-x)\big)\delta(u)\big],
\]
que se obtiene aplicando el Teorema~\ref{teo:esp-cond-malliavin} a la diferencia $\delta_x-\psi(\cdot-x)$. Bouchard, Ekeland
y Touzi probaron que, en la familia $\psi(z)=\frac\eta2e^{-\eta\abs z}$, el parámetro que minimiza la varianza
integrada es explícito: $\eta^2=\E[\Phi(G)^2\delta(u)^2]/\E[\Phi(G)^2]$.

\section{Opciones americanas por Monte Carlo}
Una put americana (o bermuda, en una rejilla $0=t_0<\dots<t_n=T$) tiene precio dado por el
principio de programación dinámica:
\[
  V_{t_n}=\Phi(S_{t_n}),\qquad V_{t_k}=\max\Big\{\Phi(S_{t_k}),\ e^{-r\Delta t}\,\E\big[V_{t_{k+1}}\mid S_{t_k}\big]\Big\}.
\]
El problema es la esperanza condicional, que debe evaluarse en cada punto de cada trayectoria
simulada. Hay tres familias de métodos:
\begin{enumerate}
  \item \textbf{Regresión (Longstaff--Schwartz, Tsitsiklis--Van Roy):} se proyecta $V_{t_{k+1}}$ sobre una
  base de funciones de $S_{t_k}$. Es el estándar industrial, pero requiere elegir la base.
  \item \textbf{Malliavin (Lions--Regnier, Bouchard--Touzi, Bally--Caramellino--Zanette~\cite{BallyCaramellinoZanette05}):}
  se usa el Teorema~\ref{teo:esp-cond-malliavin}, con $G=S_{t_{k+1}}$ y $F=S_{t_k}$, y los pesos del Ejemplo~\ref{ej:esp-cond-bs}. Es no
  paramétrico: no requiere base.
  \item \textbf{Núcleos:} se aproxima $\delta_x$ por un núcleo; la versión localizada de Malliavin puede verse
  como un núcleo corregido.
\end{enumerate}

\paragraph{Delta de una americana.} En la región de continuación, la delta es la derivada en $x$ del
cociente $\E[\Phi(G)\ind_{\{F>x\}}\delta(u)]/\E[\ind_{\{F>x\}}\delta(u)]$. Derivar un indicador vuelve a producir una delta
de Dirac, que se elimina con una segunda integración por partes; el resultado es otro cociente
de esperanzas no condicionales con pesos de segundo orden (iterados de $\delta$), explícitos en
Black--Scholes~\cite{BallyCaramellinoZanette05}. Así, el método de Malliavin produce precio y cobertura con las mismas
trayectorias.

La convergencia del método es del orden $N^{-1/2}$ por fecha, con constantes que crecen con el
número de fechas de ejercicio; en dimensión alta (cestas) la ventaja de no necesitar una base es
significativa, aunque la varianza del peso crece con la dimensión.

\section{La fórmula de descomposición de Alòs}
Consideramos un modelo de volatilidad estocástica para el log-precio $X_t=\log S_t$ bajo la medida
de riesgo neutro:
\[
  \dd X_t=\Big(r-\frac{\sigma_t^2}{2}\Big)\dd t+\sigma_t\Big(\rho\,\dd W_t+\sqrt{1-\rho^2}\,\dd B_t\Big),
\]
con $W,B$ brownianos independientes, $\rho\in(-1,1)$, y $\sigma_t$ un proceso positivo, adaptado a la
filtración de $W$, con $\sigma\in\D^{1,2}_W$ (derivada de Malliavin respecto a $W$). Sea $BS(t,x,\sigma)$ el precio
de Black--Scholes de una call con log-subyacente $x$ y volatilidad $\sigma$, y $V_t$ el precio en el modelo.

\begin{definicion}
La \emph{volatilidad media futura} es $v_t=\Big(\frac1{T-t}\int_t^T\sigma_s^2\,\dd s\Big)^{1/2}$. No es adaptada:
anticipa el futuro.
\end{definicion}

\begin{teorema}[Hull--White, 1987]
Si $\rho=0$, $V_t=\E_t\big[BS(t,X_t,v_t)\big]$.
\end{teorema}

\begin{proof}
Condicionando a toda la trayectoria de $\sigma$ (es decir, a $\F^W_T$), $X_T$ es gaussiana
con varianza $\int_t^T\sigma^2=(T-t)v_t^2$, y el precio condicional es exactamente $BS(t,X_t,v_t)$.
\end{proof}

Cuando $\rho\ne0$ el argumento falla, porque la trayectoria de $\sigma$ y la del subyacente están
correlacionadas. Alòs~\cite{Alos06} encontró la corrección exacta.

\begin{teorema}[Alòs, 2006]\label{teo:alos}
Bajo condiciones de integrabilidad,
\[
  V_t=\E_t\big[BS(t,X_t,v_t)\big]+\frac\rho2\,\E_t\Big[\int_t^Te^{-r(s-t)}H(s,X_s,v_s)\,\Lambda_s\,\dd s\Big],
\]
donde $H=(\partial_x^3-\partial_x^2)BS$ y
\[
  \Lambda_s=\sigma_s\int_s^TD^W_s\big(\sigma_r^2\big)\,\dd r.
\]
\end{teorema}

\begin{proof}[Esbozo]
Se aplica una \emph{fórmula de Itô anticipativa} al proceso $e^{-rs}BS(s,X_s,v_s)$ en $[t,T]$. Como
$v_s$ no es adaptado, $\int\cdot\,\dd W$ se entiende como integral de Skorokhod, y la fórmula de Itô para
integrales de Skorokhod (Nualart--Pardoux) contiene un término adicional: si $Y_s=\int_s^T\sigma_r^2\,\dd r$,
que es $\F_T$-medible, el término de corrección es $\frac12\partial^2_{xY}(\cdot)\,\dd\langle X,Y\rangle$ ``anticipativo'', que se escribe
$\rho\,\sigma_s\,(D^W_sY_s)\,\partial_x\partial_Y(\cdot)\,\dd s$. Usando que $BS$ satisface la EDP de Black--Scholes y la relación
$\partial_\sigma BS=\sigma(T-t)(\partial_x^2-\partial_x)BS$ (propia de la gaussiana), los términos sin $\rho$ se cancelan, la integral
de Skorokhod tiene esperanza cero, y queda exactamente $\frac\rho2\int H\Lambda\,\dd s$. Como $BS(T,X_T,v_T)=
(e^{X_T}-K)^+$, se obtiene la fórmula. Véase~\cite{Alos06,AlosGarciaLorite21}.
\end{proof}

\begin{intuicion}
El primer término es el precio de Hull--White: ``Black--Scholes con la volatilidad media
futura''. El segundo corrige por la correlación: $\Lambda_s$ mide cuánto cambia la volatilidad
futura ante un choque en el browniano del subyacente en el instante $s$. Si $\rho<0$ (efecto
apalancamiento: cuando el precio baja, la volatilidad sube), la corrección hace las puts
fuera del dinero más caras, lo que se traduce en un \emph{skew} negativo de volatilidad implícita.
\end{intuicion}

\section{El skew de volatilidad implícita en tiempo corto}
Sea $I(t,k)$ la volatilidad implícita para log-strike $k$ y $k_t^*=X_t+r(T-t)$ el strike forward en el
dinero. La fórmula de Alòs permite desarrollar $I$ en torno a $k^*$ cuando $T\to t$.

\begin{teorema}[Alòs--León--Vives~\cite{AlosLeonVives07}, forma heurística]\label{teo:skew}
Bajo condiciones de regularidad,
cuando $T-t\to0$, la pendiente en el dinero satisface, a primer orden,
\[
  \partial_kI(t,k_t^*)\approx\frac{\rho}{\sigma_t\,(T-t)^2}\int_t^T\int_s^T\E_t\big[D^W_s\sigma_r\big]\,\dd r\,\dd s.
\]
En particular:
\begin{enumerate}[label=(\alph*)]
  \item si $\E_t[D^W_s\sigma_r]\to c_t$ cuando $r,s\downarrow t$ (volatilidad dirigida por una difusión, como Heston o
  SABR), el skew converge a una constante, $\frac{\rho c_t}{2\sigma_t}$;
  \item si $\E_t[D^W_s\sigma_r]\sim c\,(r-s)^{H-1/2}$ con $H\in(0,1/2)$, entonces $\partial_kI(t,k^*)\sim C\rho\,(T-t)^{H-1/2}$,
  que explota cuando $T\to t$.
\end{enumerate}
\end{teorema}

Empíricamente, el skew en el dinero de los índices bursátiles se comporta como $(T-t)^{-\alpha}$
con $\alpha\approx0{,}4$ para vencimientos cortos, algo que ningún modelo de difusión clásico reproduce
(por (a)). El caso (b) es exactamente el que producen los modelos de \textbf{volatilidad rugosa}.

\begin{ejemplo}[Bergomi rugoso~\cite{BayerFrizGatheral16}]
Sea $\hat W^H_t=\int_0^tK(t,s)\,\dd W_s$ con $K(t,s)=\sqrt{2H}(t-s)^{H-1/2}$
(proceso de Riemann--Liouville), y
\[
  \sigma_t^2=\xi_0(t)\exp\Big(\eta\hat W^H_t-\frac{\eta^2}{2}t^{2H}\Big).
\]
Por la regla de la cadena, $D^W_s\sigma_r^2=\eta\,\sigma_r^2K(r,s)=\eta\sqrt{2H}\,\sigma_r^2\,(r-s)^{H-1/2}$ para $s<r$, de modo
que $D^W_s\sigma_r=\frac\eta2\sqrt{2H}\,\sigma_r(r-s)^{H-1/2}$. Sustituyendo en el Teorema~\ref{teo:skew},
\[
  \int_t^T\int_s^T(r-s)^{H-1/2}\,\dd r\,\dd s=\frac{(T-t)^{H+3/2}}{(H+\frac12)(H+\frac32)},
\]
y el skew se comporta como $(T-t)^{H-1/2}$. Con $H\approx0{,}1$, como estiman Gatheral, Jaisson y
Rosenbaum~\cite{GatheralJaissonRosenbaum18} a partir de series de volatilidad realizada, se obtiene $\alpha\approx0{,}4$.
\end{ejemplo}

\begin{observacion}
La fórmula de Alòs y sus desarrollos se han convertido en una herramienta
estándar para obtener aproximaciones de precios y volatilidades implícitas en modelos rugosos,
donde no hay EDP de valoración (el modelo no es markoviano) y la simulación es costosa. La
monografía de Alòs y García-Lorite~\cite{AlosGarciaLorite21} recoge el estado de la cuestión, incluidas fórmulas para la
curvatura de la sonrisa y para la volatilidad implícita del VIX. El modelo de Heston rugoso de
El Euch y Rosenbaum~\cite{ElEuchRosenbaum19} tiene función característica semiexplícita vía ecuaciones de Riccati
fraccionarias, lo que permite comparar con los desarrollos de Malliavin.
\end{observacion}


\part{Ecuaciones en derivadas parciales estocásticas}
\chapter{Integración estocástica para EDPEs}\label{cap:integracion-edpe}

\begin{panorama}
Las ecuaciones en derivadas parciales estocásticas (EDPEs) son ecuaciones de evolución
dirigidas por un ruido que depende del espacio y del tiempo. Presentamos el enfoque de
\emph{campos aleatorios} de Walsh y su extensión por Dalang, en el que la solución es una función
$u(t,x)$ y el ruido es un proceso isonormal sobre un espacio $\Hf$ adaptado a su covarianza
espacial. Probamos la \emph{condición de Dalang}, criterio necesario y suficiente para que la
ecuación lineal tenga solución en forma de función, y el teorema de existencia y unicidad
para la ecuación no lineal. Este marco es el escenario natural para el cálculo de Malliavin
de los Capítulos~\ref{cap:densidades-edpe} y~\ref{cap:anderson}.
\end{panorama}

\section{Motivación: la ecuación del calor estocástica}
La ecuación modelo es
\begin{equation}\label{eq:calor-estocastica}
  \partial_tu(t,x)=\tfrac12\Delta u(t,x)+b(u(t,x))+\sigma(u(t,x))\,\dot W(t,x),\qquad t>0,\ x\in\R^d,
\end{equation}
con dato inicial $u(0,\cdot)=u_0$. Aquí $\dot W$ es un ruido gaussiano centrado, blanco en tiempo, con
covarianza espacial $\gamma$:
\[
  \E\big[\dot W(t,x)\dot W(s,y)\big]=\delta(t-s)\,\gamma(x-y).
\]
El caso $\gamma=\delta_0$ es el \emph{ruido blanco espacio-temporal}. Las aplicaciones incluyen la propagación de
calor con fuentes aleatorias, la dinámica de poblaciones (Fisher--KPP estocástica), los polímeros
dirigidos en medio aleatorio y la superficie de crecimiento KPZ (vía la transformación de
Cole--Hopf, Capítulo~\ref{cap:anderson}).

\paragraph{Una primera dificultad.} Con $\sigma\equiv1$, $b\equiv0$, $u_0\equiv0$ y ruido blanco, la solución sería
\[
  u(t,x)=\int_0^t\int_{\R^d}G(t-s,x-y)\,W(\dd s,\dd y),\qquad G(t,x)=(2\pi t)^{-d/2}e^{-\abs x^2/2t},
\]
una variable gaussiana de varianza $\int_0^t\int G(s,y)^2\,\dd y\,\dd s=\int_0^t(4\pi s)^{-d/2}\,\dd s$, que es finita \emph{si y solo
si} $d=1$. En $d\ge2$ el ruido blanco es demasiado irregular y la solución solo existe como
distribución. Esto motiva tanto los ruidos coloreados (Dalang) como la teoría de estructuras de
regularidad (Capítulo~\ref{cap:estructuras}).

\section{El ruido como proceso isonormal}
Sea $\gamma$ una medida no negativa, simétrica y \emph{definida no negativa} en $\R^d$, de modo que, por el
teorema de Bochner--Schwartz, $\gamma=\hat\mu$ para una medida temperada $\mu\ge0$, la \emph{medida espectral}.
Definimos $\Hc_0$ como la compleción de $C_c^\infty(\R^d)$ con
\[
  \ip{\varphi}{\psi}_{\Hc_0}=\int_{\R^{2d}}\varphi(x)\psi(y)\gamma(x-y)\,\dd x\,\dd y=\int_{\R^d}\hat\varphi(\xi)\overline{\hat\psi(\xi)}\,\mu(\dd\xi),
\]
y $\Hf=L^2([0,T];\Hc_0)$. El ruido es un proceso isonormal $\{W(h):h\in\Hf\}$. Para $h=\ind_{[0,t]}\varphi$,
$W_t(\varphi)=W(\ind_{[0,t]}\varphi)$ define, para cada $\varphi$, un browniano con varianza $t\norm{\varphi}^2_{\Hc_0}$: en el lenguaje de
Da Prato--Zabczyk, $W$ es un \emph{browniano cilíndrico} en $\Hc_0$.

\begin{ejemplo}
\begin{enumerate}[label=(\alph*)]
  \item Ruido blanco: $\gamma=\delta_0$, $\mu(\dd\xi)=(2\pi)^{-d}\,\dd\xi$, $\Hc_0=L^2(\R^d)$.
  \item Núcleo de Riesz: $\gamma(x)=\abs x^{-\beta}$, $0<\beta<d$, $\mu(\dd\xi)=c_{d,\beta}\abs\xi^{\beta-d}\,\dd\xi$.
  \item Ruido fraccionario en espacio, en $d=1$: $\gamma(x)=H(2H-1)\abs x^{2H-2}$ con $H>1/2$~\cite{HuNualartSong11}.
\end{enumerate}
\end{ejemplo}

\section{La integral estocástica de Walsh--Dalang}
Sea $\F_t$ la filtración generada por $\{W_s(\varphi):\ s\le t,\ \varphi\in\Hc_0\}$. Para un campo aleatorio $X(s,y)$
\emph{predecible} (medible respecto a la $\sigma$-álgebra generada por los elementales $Z\ind_{(a,b]}(s)\ind_A(y)$ con
$Z\in\F_a$) definimos
\[
  \int_0^t\int_{\R^d}X(s,y)\,W(\dd s,\dd y)
\]
primero para elementales, como $Z(W_b(\ind_A)-W_a(\ind_A))$, y después por la isometría.

\begin{teorema}[Isometría de Walsh--Dalang]\label{teo:walsh-dalang}
Para $X$ predecible con $\E\int_0^T\norm{X(s,\cdot)}^2_{\Hc_0}\,\dd s<
\infty$, la integral está bien definida, es una martingala continua en $t$ y
\[
  \E\Big[\Big(\int_0^t\!\!\int X\,\dd W\Big)^2\Big]=\E\int_0^t\int_{\R^{2d}}X(s,y)X(s,z)\gamma(y-z)\,\dd y\,\dd z\,\dd s=\E\int_0^t\norm{X(s,\cdot)}^2_{\Hc_0}\,\dd s.
\]
Además vale la desigualdad de BDG: $\E\big|\int_0^t\int X\,\dd W\big|^p\le C_p\E\big(\int_0^t\norm{X(s)}^2_{\Hc_0}\,\dd s\big)^{p/2}$.
\end{teorema}

\begin{proof}
Para elementales es el cálculo de la isometría de Itô (Lema~\ref{lem:isometria-elemental}), usando que
$\E[W_b(\ind_A)W_b(\ind_B)]=b\ip{\ind_A}{\ind_B}_{\Hc_0}$. La extensión es por densidad, como en el Teorema~\ref{teo:integral-ito}. La
BDG es el Teorema~\ref{teo:bdg} aplicado a la martingala $t\mapsto\int_0^t\int X\,\dd W$, cuya variación cuadrática es
$\int_0^t\norm{X(s)}^2_{\Hc_0}\,\dd s$.
\end{proof}

Dalang~\cite{Dalang99} extendió la integral a integrandos de la forma $G(t-s,x-\cdot)Z(s,\cdot)$ en los que
$G(t-s,x-\cdot)$ es una \emph{distribución} (como el núcleo de ondas en $d=3$), lo que requiere que $\Hc_0$
contenga distribuciones.

\section{La condición de Dalang}
Consideramos un operador diferencial $\mathcal L$ con solución fundamental $G$ y la ecuación lineal $\mathcal Lu=\dot W$
con dato nulo: $u(t,x)=\int_0^t\int G(t-s,x-y)W(\dd s,\dd y)$.

\begin{teorema}[Dalang, 1999]\label{teo:dalang}
Sea $\mathcal L=\partial_t-\frac12\Delta$ (calor) o $\mathcal L=\partial^2_{tt}-\Delta$ (ondas, con $d\le3$).
La solución de la ecuación lineal es un campo aleatorio, con $\E u(t,x)^2<\infty$, si y solo si
\begin{equation}\label{eq:dalang}
  \Upsilon:=\int_{\R^d}\frac{\mu(\dd\xi)}{1+\abs\xi^2}<\infty.
\end{equation}
\end{teorema}

\begin{proof}
Por la isometría y Plancherel,
\[
  \E u(t,x)^2=\int_0^t\norm{G(s,\cdot)}^2_{\Hc_0}\,\dd s=\int_0^t\int_{\R^d}\abs{\mathcal FG(s,\cdot)(\xi)}^2\,\mu(\dd\xi)\,\dd s.
\]
\emph{Calor:} $\mathcal FG(s)(\xi)=e^{-s\abs\xi^2/2}$ y $\int_0^te^{-s\abs\xi^2}\,\dd s=\frac{1-e^{-t\abs\xi^2}}{\abs\xi^2}$. Esta función es comparable, con
constantes dependientes de $t$, a $(1+\abs\xi^2)^{-1}$.

\emph{Ondas:} $\mathcal FG(s)(\xi)=\frac{\sin(s\abs\xi)}{\abs\xi}$ y $\int_0^t\frac{\sin^2(s\abs\xi)}{\abs\xi^2}\,\dd s=\frac{t}{2\abs\xi^2}-\frac{\sin(2t\abs\xi)}{4\abs\xi^3}$, de nuevo comparable a
$(1+\abs\xi^2)^{-1}$ (cerca de $\xi=0$ vale $\approx t^3/3$).

En ambos casos $\E u(t,x)^2<\infty\iff\Upsilon<\infty$.
\end{proof}

\begin{ejemplo}
Ruido blanco: $\Upsilon=\int\frac{\dd\xi}{1+\abs\xi^2}<\infty\iff d=1$. Núcleo de Riesz $\abs x^{-\beta}$:
$\int\frac{\abs\xi^{\beta-d}}{1+\abs\xi^2}\,\dd\xi<\infty\iff\beta<2$, en cualquier dimensión.
\end{ejemplo}

\section{La ecuación no lineal}

\begin{definicion}
Una \emph{solución mild} de~\eqref{eq:calor-estocastica} es un campo aleatorio predecible $u$ tal que,
para todo $(t,x)$,
\begin{align*}
  u(t,x)={}&\int G(t,x-y)u_0(y)\,\dd y+\int_0^t\!\!\int G(t-s,x-y)b(u(s,y))\,\dd y\,\dd s\\
  &+\int_0^t\!\!\int G(t-s,x-y)\sigma(u(s,y))W(\dd s,\dd y).
\end{align*}
\end{definicion}

\begin{teorema}[Walsh; Dalang]\label{teo:existencia-edpe}
Supongamos~\eqref{eq:dalang}, $b,\sigma$ globalmente Lipschitz y $u_0$ acotado
y medible. Entonces existe una única solución mild con $\sup_{t\le T,x\in\R^d}\E\abs{u(t,x)}^p<\infty$ para
todo $p$. Si además $\int\frac{\mu(\dd\xi)}{(1+\abs\xi^2)^{1-\eta}}<\infty$ para algún $\eta\in(0,1)$, entonces $u$ tiene una versión
localmente Hölder de orden $<\eta/2$ en $t$ y $<\eta$ en $x$.
\end{teorema}

\begin{proof}[Esbozo]
Iteración de Picard: $u^{n+1}=\Phi(u^n)$ con $\Phi$ el lado derecho. Por BDG (Teorema~\ref{teo:walsh-dalang}) y
Minkowski,
\[
  \sup_x\norm{u^{n+1}(t,x)-u^n(t,x)}_p^2\le C\int_0^t\sup_y\norm{u^n(s,y)-u^{n-1}(s,y)}_p^2\,J(t-s)\,\dd s,
\]
con $J(s)=\norm{G(s)}^2_{\Hc_0}+\norm{G(s)}_{L^1}$ (para la parte de $\sigma$ se usa que $\gamma\ge0$: desarrollando
$\norm{G(t-s,x-\cdot)Z(s,\cdot)}^2_{\Hc_0}$ como integral doble contra $\gamma$ y aplicando Minkowski y Cauchy--Schwarz,
se obtiene la cota $\sup_y\norm{Z(s,y)}_p^2\norm{G(t-s)}^2_{\Hc_0}$). La condición de Dalang dice exactamente
$\int_0^TJ<\infty$. Un \emph{lema de Gronwall con núcleo} (Dalang, Lemma~15) muestra que si
$f_n(t)\le\int_0^tf_{n-1}(s)J(t-s)\,\dd s$, entonces $\sum_nf_n(t)^{1/2}<\infty$ uniformemente en $[0,T]$. La regularidad
se prueba estimando incrementos y aplicando Kolmogorov (Teorema~\ref{teo:kolmogorov}).
\end{proof}

\begin{ejemplo}
Para ruido blanco en $d=1$, $\mu=\dd\xi/2\pi$ y la condición con $\eta$ vale para $\eta<1/2$:
$u$ es Hölder de orden $<1/4$ en tiempo y $<1/2$ en espacio, como el browniano en espacio y
como el browniano fraccionario de Hurst $1/4$ en tiempo.
\end{ejemplo}

\section{El marco de Malliavin para EDPEs}
El ruido es un proceso isonormal sobre $\Hf=L^2([0,T];\Hc_0)$, y por tanto toda la Parte~II se
aplica. La derivada de Malliavin de un funcional $F$ es un elemento de $\Hf$, que denotamos $D_{r,\cdot}F$
o, cuando es función, $D_{r,z}F$. En el caso de ruido blanco, $\Hf=L^2([0,T]\times\R^d)$ y
\[
  D_{r,z}W(\ind_{[0,t]\times A})=\ind_{[0,t]}(r)\ind_A(z),\qquad\norm{DF}^2_\Hf=\int_0^T\int_{\R^d}\abs{D_{r,z}F}^2\,\dd z\,\dd r.
\]
La divergencia $\delta$ restringida a campos predecibles es la integral de Walsh--Dalang, como en el
Teorema~\ref{teo:skorokhod-ito}, y la fórmula de Clark--Ocone tiene la forma $F=\E F+\int_0^T\int\E[D_{r,z}F\mid\F_r]W(\dd r,\dd z)$.

\chapter{Densidades para EDPEs}\label{cap:densidades-edpe}

\begin{panorama}
Probamos que la solución de la ecuación del calor estocástica es Malliavin-diferenciable a
todos los órdenes, escribimos la ecuación lineal que satisface su derivada y deducimos la
existencia de densidad (Pardoux--Zhang) y su regularidad $C^\infty$ (Bally--Pardoux, Mueller--Nualart).
La pieza clave es una estimación de pequeñas probabilidades para $\norm{Du(t,x)}^2_\Hf$
que explota la ventana temporal $[t-\delta,t]$. Discutimos la positividad estricta, las densidades
conjuntas y las probabilidades de alcance, y la ecuación de ondas.
\end{panorama}

Trabajamos con la ecuación~\eqref{eq:calor-estocastica} en $d=1$ con ruido blanco, salvo indicación contraria; la
mayoría de los resultados se extienden a ruidos que satisfacen la condición de Dalang.

\section{Diferenciabilidad de Malliavin}

\begin{teorema}\label{teo:derivada-edpe}
Si $b,\sigma\in C^\infty$ con derivadas acotadas, entonces $u(t,x)\in\D^\infty$ para todo $(t,x)$,
y la derivada satisface, para $r<t$,
\begin{equation}\label{eq:derivada-edpe}
\begin{aligned}
  D_{r,z}u(t,x)={}&G(t-r,x-z)\sigma(u(r,z))+\int_r^t\!\!\int G(t-s,x-y)\,b'(u(s,y))D_{r,z}u(s,y)\,\dd y\,\dd s\\
  &+\int_r^t\!\!\int G(t-s,x-y)\,\sigma'(u(s,y))D_{r,z}u(s,y)\,W(\dd s,\dd y),
\end{aligned}
\end{equation}
y $D_{r,z}u(t,x)=0$ si $r>t$. Además, $\sup_{(t,x)}\E\norm{D^ku(t,x)}^p_{\Hf^{\otimes k}}<\infty$ para todo $k,p$.
\end{teorema}

\begin{proof}[Esbozo]
Las iteraciones de Picard $u^n$ son Malliavin-diferenciables. Derivando la ecuación mild
con la regla de la cadena y el Corolario~\ref{cor:derivada-ito} (versión Walsh), $Du^{n+1}$ satisface la versión lineal
de~\eqref{eq:derivada-edpe} con $u^n$ en los coeficientes. El término $G(t-r,x-z)\sigma(u^n(r,z))$ es la ``derivada
del integrando respecto al ruido en el propio punto $(r,z)$'', el análogo del término $u_t$ en el
Corolario~\ref{cor:derivada-ito}. Las estimaciones del Teorema~\ref{teo:existencia-edpe}, aplicadas ahora en el espacio $L^p(\Omega;\Hf)$, dan
cotas uniformes en $n$, y la cerrabilidad de $D$ permite pasar al límite. Las derivadas superiores
satisfacen ecuaciones lineales análogas con términos fuente que involucran derivadas de orden
inferior. Véase~\cite{BallyPardoux98,Nualart06,SanzSole05}.
\end{proof}

\begin{intuicion}
La ecuación~\eqref{eq:derivada-edpe} es la versión de dimensión infinita del Teorema~\ref{teo:derivada-ede}: un impulso del
ruido en $(r,z)$ produce un desplazamiento $\sigma(u(r,z))$ en el punto $z$, que después se propaga
por el calor y por la linealización de la ecuación. Si se desprecia la parte linealizada,
$D_{r,z}u(t,x)\approx G(t-r,x-z)\sigma(u(r,z))$: la sensibilidad está concentrada en un ``cono
parabólico'' detrás de $(t,x)$.
\end{intuicion}

\section{Existencia de densidad}

\begin{teorema}[Pardoux--Zhang, 1993]
Supongamos $b,\sigma$ Lipschitz, $\sigma\in C^1$, $u_0$ continuo y
acotado en $\R$, y que existe $y_0$ con $\sigma(u_0(y_0))\ne0$. Entonces para todo $t>0$ y $x\in\R$, la ley de
$u(t,x)$ es absolutamente continua.
\end{teorema}

\begin{proof}[Esbozo]
Por el criterio de Bouleau--Hirsch (Teorema~\ref{teo:bouleau-hirsch}) basta ver que $\norm{Du(t,x)}_\Hf>0$ c.s.
La dificultad es que $\sigma(u(r,z))$ puede anularse. Se procede en dos pasos. (1) Por continuidad,
$\sigma(u(r,z))\ne0$ para $(r,z)$ en un entorno de $(0,y_0)$ con probabilidad positiva, y un argumento de
\emph{propagación} usando la positividad del núcleo del calor muestra que el conjunto de $(r,z)$ donde
$\sigma(u(r,z))\ne0$ alcanza, c.s., un entorno de cualquier $z$ para $r$ próximo a $t$. (2) En una ventana
$[t-\delta,t]\times\{\abs{z-x}\le\delta^{1/2}\}$ donde $\abs{\sigma(u)}\ge c>0$, el término principal de~\eqref{eq:derivada-edpe} domina al resto
(siguiente sección), lo que da $\norm{Du(t,x)}>0$.
\end{proof}

\section{Regularidad de la densidad}

\begin{teorema}[Bally--Pardoux, 1998; Mueller--Nualart, 2008]\label{teo:regularidad-edpe}
Supongamos $b,\sigma\in C^\infty$ con
derivadas acotadas. Si $\abs{\sigma(v)}\ge c>0$ para todo $v$ \textup{(elipticidad)}, entonces para todo $t>0$
y $x\in\R$, $u(t,x)$ tiene densidad $C^\infty$. Mueller y Nualart~\cite{MuellerNualart08} prueban la misma conclusión,
para la ecuación en $[0,1]$ con condiciones de Dirichlet, bajo la mera hipótesis $\sigma(u_0(y_0))\ne0$
para algún $y_0$.
\end{teorema}

\begin{proof}[Demostración en el caso elíptico]
Por los Teoremas~\ref{teo:malliavin-densidad} y~\ref{teo:derivada-edpe}, basta probar que
$\norm{Du(t,x)}^{-2}_\Hf\in\bigcap_pL^p$, y para ello que, para todo $p$, $\P(\norm{Du(t,x)}^2_\Hf<\varepsilon)\le C_p\varepsilon^p$ para $\varepsilon$
pequeño.

Sea $\delta\in(0,t)$ que fijaremos en función de $\varepsilon$. Restringiendo la norma a la ventana $[t-\delta,t]$ y
usando $(a+b)^2\ge\frac12a^2-b^2$,
\[
  \norm{Du(t,x)}^2_\Hf\ge\int_{t-\delta}^t\int_\R\abs{D_{r,z}u(t,x)}^2\,\dd z\,\dd r\ge\frac12I_1(\delta)-I_2(\delta),
\]
donde
\[
  I_1(\delta)=\int_{t-\delta}^t\int G(t-r,x-z)^2\sigma(u(r,z))^2\,\dd z\,\dd r,\qquad I_2(\delta)=\int_{t-\delta}^t\int\abs{R_{r,z}(t,x)}^2\,\dd z\,\dd r,
\]
y $R_{r,z}(t,x)$ es la integral (deriva más ruido) en~\eqref{eq:derivada-edpe}.

\emph{Cota inferior de $I_1$.} Como $\int_\R G(s,y)^2\,\dd y=(4\pi s)^{-1/2}$,
\[
  I_1(\delta)\ge c^2\int_0^\delta(4\pi s)^{-1/2}\,\dd s=c_1\,\delta^{1/2}.
\]

\emph{Cota superior de $I_2$.} Usando BDG para la parte de ruido, Hölder para la deriva, y las cotas
uniformes de $\E\abs{D_{r,z}u(s,y)}^p$ integradas en $(r,z)$ sobre la ventana, se prueba
\[
  \E\abs{I_2(\delta)}^p\le C_p\,\delta^p.
\]
La razón heurística: $R_{r,z}$ es una integral sobre $s\in[r,t]$, intervalo de longitud $\le\delta$, de una
cantidad del tamaño de $D_{r,z}u(s,y)$; integrando su cuadrado en la ventana se gana un factor
$\delta^{1/2}$ adicional respecto a $I_1$.

\emph{Conclusión.} Elegimos $\delta$ de modo que $\frac{c_1}{2}\delta^{1/2}=2\varepsilon$, es decir, $\delta=(4\varepsilon/c_1)^2$. Entonces
\[
  \P\big(\norm{Du(t,x)}^2_\Hf<\varepsilon\big)\le\P\big(I_2(\delta)\ge\tfrac{c_1}{2}\delta^{1/2}-\varepsilon\big)=\P(I_2(\delta)\ge\varepsilon)\le\frac{\E\abs{I_2(\delta)}^p}{\varepsilon^p}\le C\frac{\delta^p}{\varepsilon^p}=C'\varepsilon^p.
  \qedhere
\]
\end{proof}

\begin{observacion}
El exponente $1/2$ de la cota de $I_1$ refleja la regularidad temporal $1/4$ de $u$
(pues $\norm{Du}$ mide una ``varianza local'' de orden $\delta^{1/2}=(\delta^{1/4})^2$). En ruidos coloreados con medida
espectral $\mu$, se obtiene $I_1(\delta)\ge c\int_0^\delta\int e^{-s\abs\xi^2}\mu(\dd\xi)\,\dd s$, y la prueba funciona si esta cantidad es al
menos $\delta^\alpha$ y $\E I_2^p\le C\delta^{\beta p}$ con $\beta>\alpha$. Esto da la teoría de Nualart y Quer-Sardanyons~\cite{NualartQuerSardanyons07} para
ecuaciones espacialmente homogéneas.
\end{observacion}

\section{Positividad estricta}

\begin{teorema}[Mueller--Nualart; Chen--Hu--Nualart~\cite{ChenHuNualart21}]
Bajo las hipótesis del Teorema~\ref{teo:regularidad-edpe} en el caso elíptico, la densidad de $u(t,x)$ es estrictamente positiva en todo $\R$.
Para la ecuación con $\sigma(0)=0$ y $u_0\ge0$ (por ejemplo, el modelo parabólico de Anderson,
$\sigma(u)=u$), la solución es estrictamente positiva y la densidad es positiva en $(0,\infty)$.
\end{teorema}

La prueba usa el criterio de Bally--Pardoux, que es esencialmente una versión cuantitativa
del teorema de Cameron--Martin: si $F$ es no degenerado en un entorno de $y$ y existe una
dirección $h\in\Hf$ tal que $F(\omega+h)$ ``apunta'' a $y$ con probabilidad positiva, entonces $p_F(y)>0$.
Las direcciones $h$ se construyen resolviendo un problema de control: elegir $h$ como control de la
EDP determinista $\partial_tv=\frac12\Delta v+b(v)+\sigma(v)h$ que lleva $v(t,x)$ a $y$.

\section{Densidades conjuntas y probabilidades de alcance}
Para el vector $(u(t,x_1),\dots,u(t,x_n))$ o para sistemas $u=(u^1,\dots,u^d)$ de ecuaciones del calor
con valores en $\R^d$, la matriz de Malliavin es $d\times d$ y la no degeneración requiere controlar su
autovalor mínimo. Dalang, Khoshnevisan y Nualart~\cite{DalangKhoshnevisanNualart07} obtuvieron \emph{cotas gaussianas} para la
densidad conjunta de $(u(t,x),u(s,y))$:
\[
  p_{t,x;s,y}(z_1,z_2)\le C\,\Delta^{-d/2}\exp\Big(-\frac{\abs{z_1-z_2}^2}{C\Delta}\Big),\qquad\Delta=\abs{t-s}^{1/2}+\abs{x-y},
\]
y, mediante las técnicas de Khoshnevisan de capacidad, las \emph{probabilidades de alcance}: para un
conjunto $A\subset\R^d$,
\[
  c\operatorname{Cap}_{d-6}(A)\le\P\big(u(I\times J)\cap A\ne\emptyset\big)\le C\,\mathcal H_{d-6}(A),
\]
donde $\mathcal H_\beta$ es la medida de Hausdorff y $\operatorname{Cap}_\beta$ la capacidad de Bessel--Riesz. El número 6 es
la ``dimensión'' de $(t,x)\in\R_+\times\R$ con la métrica parabólica $\abs{t-s}^{1/4}+\abs{x-y}^{1/2}$: $4+2$. En
particular, la superficie aleatoria $u$ alcanza puntos fijos con probabilidad positiva si $d<6$ y no
los alcanza si $d>6$; el caso crítico $d=6$ requirió trabajo posterior.

\section{La ecuación de ondas}
Para $\partial^2_{tt}u=\Delta u+b(u)+\sigma(u)\dot W$ en $d\le3$ con ruido que satisface la condición de Dalang, el
análogo de la ecuación~\eqref{eq:derivada-edpe} vale con $G$ el núcleo de ondas. En $d=3$, $G(t)$ es la medida uniforme
en la esfera de radio $t$ (normalizada), una distribución no funcional, y hay que usar la integral de
Dalang. La estrategia de la ventana temporal funciona con $I_1(\delta)\ge c\int_0^\delta\int\frac{\sin^2(s\abs\xi)}{\abs\xi^2}\mu(\dd\xi)\,\dd s\ge c\delta^3$
cerca de 0. Quer-Sardanyons y Sanz-Solé probaron la existencia de densidades regulares en
$d=3$; la falta de positividad del núcleo de ondas en $d\ge4$, y de efecto regularizante, hace la
teoría mucho más delicada que la del calor. Véase~\cite{SanzSole05}.

\chapter{El modelo parabólico de Anderson y teoremas límite espaciales}\label{cap:anderson}

\begin{panorama}
El modelo parabólico de Anderson (PAM) es la ecuación del calor con ruido multiplicativo
lineal, $\sigma(u)=u$. Es exactamente resoluble por caos de Wiener y por una fórmula
de Feynman--Kac, y exhibe \emph{intermitencia}: sus momentos crecen a tasas exponenciales
distintas, lo que indica picos raros y altísimos. Su logaritmo es la solución de Cole--Hopf
de la ecuación KPZ. En la segunda mitad del capítulo presentamos uno de los avances
recientes más elegantes del cálculo de Malliavin aplicado: los \emph{teoremas centrales del
límite cuantitativos para promedios espaciales} de soluciones de EDPEs (Huang, Nualart,
Viitasaari y Zheng, 2020).
\end{panorama}

\section{El modelo y su desarrollo en caos}
Consideramos, en $\R_+\times\R$ con ruido blanco,
\begin{equation}\label{eq:pam}
  \partial_tu=\tfrac12\partial_{xx}u+u\,\dot W,\qquad u(0,\cdot)\equiv1.
\end{equation}

\begin{teorema}
La solución mild de~\eqref{eq:pam} admite el desarrollo en caos
\[
  u(t,x)=1+\sum_{n=1}^\infty I_n\big(f_n(\cdot;t,x)\big),
\]
con
\[
  f_n(s_1,y_1,\dots,s_n,y_n;t,x)=\frac1{n!}\prod_{j=1}^nG\big(s_{\pi(j+1)}-s_{\pi(j)},y_{\pi(j+1)}-y_{\pi(j)}\big),
\]
donde $\pi$ es la permutación que ordena $s_{\pi(1)}<\dots<s_{\pi(n)}$ y $(s_{\pi(n+1)},y_{\pi(n+1)})=(t,x)$.
Equivalentemente, $n!f_n$ es la simetrización de $\prod_jG(s_{j+1}-s_j,y_{j+1}-y_j)\ind_{\{0<s_1<\dots<s_n<t\}}$.
\end{teorema}

\begin{proof}
Iterando la ecuación mild, $u(t,x)=1+\int_0^t\int G(t-s,x-y)u(s,y)W(\dd s,\dd y)$, $n$
veces, se obtiene una suma de integrales iteradas de Walsh de núcleos deterministas más un
resto. Por el Teorema~\ref{teo:iteradas}, la integral iterada $n$-ésima es $I_n$ del núcleo simetrizado. El resto
tiende a cero en $L^2$ porque, con $\norm{f_n}^2$ calculado abajo, $\sum_nn!\norm{f_n}^2<\infty$. Por unicidad de la
descomposición en caos, se concluye.
\end{proof}

\begin{proposicion}
$\E[u(t,x)^2]=\sum_{n\ge0}n!\norm{f_n}^2=\E\exp\Big(\int_0^t\delta_0(B_s-\tilde B_s)\,\dd s\Big)$, donde $B,\tilde B$ son
brownianos independientes y la integral es el tiempo local en 0 de $B-\tilde B$ (respecto a la medida
de Lebesgue). En particular, $\E[u(t,x)^2]=2e^{t/4}\Phi\big(\sqrt{t/2}\big)$.
\end{proposicion}

\begin{proof}
Usando la representación de $G$ como densidad de transición browniana y el
teorema de Fubini, $n!\norm{f_n}^2$ es el $n$-ésimo término $\frac1{n!}\E[(\int_0^t\delta_0(B_s-\tilde B_s)\,\dd s)^n]$ del desarrollo de la
exponencial. Ahora $B-\tilde B=\sqrt2\beta$ con $\beta$ browniano estándar, y $\int_0^t\delta_0(\sqrt2\beta_s)\,\dd s=\frac1{\sqrt2}L^0_t(\beta)$. Por
el teorema de Lévy, $L^0_t(\beta)\law\abs{\beta_t}\law\sqrt t\abs N$. Así, con $a=\sqrt{t/2}$, $\E[u^2]=\E e^{a\abs N}=2e^{a^2/2}\Phi(a)=
2e^{t/4}\Phi(\sqrt{t/2})$.
\end{proof}

\section{Intermitencia}

\begin{teorema}[Bertini--Cancrini, 1995]
Para todo entero $p\ge2$, la \textup{función de Lyapunov}
$\lambda(p)=\lim_{t\to\infty}\frac1t\log\E[u(t,x)^p]$ existe y vale
\[
  \lambda(p)=\frac{p(p^2-1)}{24}.
\]
\end{teorema}

\begin{proof}[Idea]
La fórmula de Feynman--Kac para momentos da $\E[u(t,x)^p]=\E\exp\big(\sum_{1\le i<j\le p}L^0_t(B^i-
B^j)\big)$, con $B^1,\dots,B^p$ brownianos independientes desde $x$. El límite $\frac1t\log$ es el autovalor principal
del operador $\frac12\Delta+\sum_{i<j}\delta(x_i-x_j)$ en $\R^p$, que es el hamiltoniano de Lieb--Liniger de $p$ bosones
con interacción atractiva delta, y se calcula exactamente por el ansatz de Bethe.
\end{proof}

\begin{intuicion}
La función $p\mapsto\lambda(p)/p$ es estrictamente creciente: es la definición de \emph{intermitencia}.
Implica que $u(t,x)$ está típicamente muy por debajo de su media $\E u=1$ (de hecho
$\frac1t\log u(t,x)\to-\frac1{24}$ c.s.), mientras que los momentos están dominados por sucesos raros
en que $u$ es exponencialmente grande. Espacialmente, la solución se organiza en ``islas''
aisladas de masa muy concentrada.
\end{intuicion}

\paragraph{Conexión con KPZ.} Formalmente, $h=\log u$ resuelve la ecuación de Kardar--Parisi--Zhang
$\partial_th=\frac12\partial_{xx}h+\frac12(\partial_xh)^2+\dot W-\infty$, donde ``$-\infty$'' es una constante de renormalización (Itô). La
\emph{solución de Cole--Hopf} se \emph{define} como $\log u$; es la que seleccionan las aproximaciones discretas y
la que reconstruye la teoría de estructuras de regularidad de Hairer (Capítulo~\ref{cap:estructuras}). Que $u>0$
(Mueller) es imprescindible para esta definición.

\section{Promedios espaciales y estacionariedad}
Si $u_0\equiv1$, por la invariancia del ruido blanco bajo traslaciones espaciales, el proceso $x\mapsto u(t,x)$
es \emph{estacionario}. Más aún, es \emph{ergódico} (Chen, Khoshnevisan, Nualart, Pu~\cite{ChenKhoshnevisanNualartPu21}, vía una desigualdad
tipo Poincaré), de modo que
\[
  \frac1{2R}\int_{-R}^Ru(t,x)\,\dd x\xrightarrow[R\to\infty]{\text{c.s.}}\E u(t,0)=1.
\]
Las fluctuaciones son la pregunta natural: ¿cómo se comporta
\[
  F_R(t)=\int_{-R}^R\big(u(t,x)-1\big)\,\dd x
\]
cuando $R\to\infty$? Para el PAM, cada $u(t,x)$ tiene infinitos caos y las correlaciones espaciales no
son finitas en rango, así que los TCLs clásicos para campos mixing no se aplican directamente.
El método de Malliavin--Stein (Capítulo~\ref{cap:malliavin-stein}) sí.

\section{El TCL cuantitativo de Huang--Nualart--Viitasaari}
El punto de partida es una cota de Stein para variables que \emph{son divergencias}.

\begin{proposicion}\label{prop:stein-divergencia}
Sea $F=\delta(v)$ con $v\in\Dom\delta$, $F\in\D^{1,2}$, y $\sigma^2=\E F^2>0$. Entonces
\[
  d_{TV}\Big(\frac F\sigma,N(0,1)\Big)\le\frac2{\sigma^2}\sqrt{\Var\big(\ip{DF}{v}_\Hf\big)}.
\]
\end{proposicion}

\begin{proof}
Por el lema de Stein (Lema~\ref{lem:stein-ecuacion} del Capítulo~\ref{cap:malliavin-stein}), $d_{TV}(X,N)\le
\sup_f\abs{\E[f'(X)-Xf(X)]}$ sobre funciones $f\in C^1$ con $\norm f_\infty\le\sqrt{\pi/2}$, $\norm{f'}_\infty\le2$. Con
$X=F/\sigma$ y la dualidad,
\[
  \E[Xf(X)]=\frac1\sigma\E[\delta(v)f(X)]=\frac1\sigma\E\ip{Df(X)}{v}=\frac1{\sigma^2}\E\big[f'(X)\ip{DF}{v}\big].
\]
Así $\E[f'(X)-Xf(X)]=\frac1{\sigma^2}\E\big[f'(X)(\sigma^2-\ip{DF}{v})\big]$. Como $\E\ip{DF}{v}=\E[F\delta(v)]=\sigma^2$, Cauchy--Schwarz
da la cota.
\end{proof}

\begin{teorema}[Huang--Nualart--Viitasaari~\cite{HuangNualartViitasaari20}, 2020]\label{teo:hnv}
Sea $u$ la solución de~\eqref{eq:calor-estocastica} con $b=0$,
$u_0\equiv1$, $d=1$, ruido blanco, $\sigma$ Lipschitz con $\sigma(1)\ne0$. Para $t>0$ fijo, $\sigma_R^2=\Var F_R(t)$
satisface $\sigma_R^2\sim\kappa_tR$ con
\[
  \kappa_t=2\int_0^t\E\big[\sigma(u(s,0))^2\big]\,\dd s>0\quad\text{(en el caso }\E\sigma(u)^2\text{ estacionario)},
\]
y existe $C=C_t$ tal que
\[
  d_{TV}\Big(\frac{F_R(t)}{\sigma_R},N(0,1)\Big)\le\frac C{\sqrt R}.
\]
Además, el proceso $\big(R^{-1/2}F_R(t)\big)_{t\in[0,T]}$ converge en $C([0,T])$ a un proceso gaussiano centrado.
\end{teorema}

\begin{proof}[Estructura de la prueba]
\emph{Paso 1: $F_R$ es una divergencia.} Integrando la ecuación mild en
$x\in[-R,R]$ y usando Fubini estocástico,
\[
  F_R(t)=\int_0^t\int_\R\varphi_R(s,y)\sigma(u(s,y))W(\dd s,\dd y)=\delta(v_R),\qquad\varphi_R(s,y)=\int_{-R}^RG(t-s,x-y)\,\dd x,
\]
con $v_R(s,y)=\varphi_R(s,y)\sigma(u(s,y))$, adaptado.

\emph{Paso 2: la varianza.} Por la isometría, $\sigma_R^2=\int_0^t\int\varphi_R(s,y)^2\E[\sigma(u(s,y))^2]\,\dd y\,\dd s$, y como
$\varphi_R\approx\ind_{[-R,R]}$ salvo en una capa de anchura $O(\sqrt t)$, $\sigma_R^2\sim R\cdot\text{const}$.

\emph{Paso 3: la derivada.} Por la Proposición~\ref{prop:heisenberg} (versión Walsh) y la ecuación~\eqref{eq:derivada-edpe},
\[
  D_{r,z}F_R=\varphi_R(r,z)\sigma(u(r,z))+\int_r^t\!\!\int\varphi_R(s,y)\sigma'(u(s,y))D_{r,z}u(s,y)W(\dd s,\dd y).
\]
Por tanto $\ip{DF_R}{v_R}_\Hf=\int_0^t\int\varphi_R^2\sigma(u)^2\,\dd z\,\dd r+\text{(término de ruido)}$.

\emph{Paso 4: la varianza de $\ip{DF_R}{v_R}$.} Es la parte técnica. Se acota $\Var$ de cada término
mediante la desigualdad de Poincaré (Teorema~\ref{teo:poincare}) y la clave es la \emph{cota de la derivada}
\[
  \norm{D_{r,z}u(t,x)}_p\le C\,G(t-r,x-z),
\]
que dice que la sensibilidad de $u(t,x)$ al ruido en $(r,z)$ decae como el núcleo del calor
(se prueba por iteración de~\eqref{eq:derivada-edpe} y un lema de Gronwall con núcleo). Con ella, las covarianzas
$\Cov(\sigma(u(s,y))^2,\sigma(u(s,y'))^2)$ decaen lo bastante rápido en $\abs{y-y'}$ como para que
$\Var\ip{DF_R}{v_R}=O(R)$. Por la Proposición~\ref{prop:stein-divergencia}, $d_{TV}\le\frac CR\sqrt R=\frac C{\sqrt R}$.

\emph{Paso 5: TCL funcional.} Convergencia de las distribuciones finito-dimensionales por la
versión multivariante, y tensión por Kolmogorov usando momentos de los incrementos en $t$.
\end{proof}

\begin{observacion}
La tasa $R^{-1/2}$ es la del teorema de Berry--Esseen para sumas de variables
independientes: en la escala espacial, la solución se comporta como si tuviera dependencia
de corto alcance. El resultado se ha extendido en múltiples direcciones: ruido coloreado de
tipo Riesz, con tasas $R^{-\beta/2}$ y escalas $R^{d-\beta/2}$ (\cite{HuangNualartViitasaariZheng20}); la ecuación de ondas; ruido fraccionario
en tiempo; dato inicial de Dirac (el PAM con $u_0=\delta_0$, relevante para polímeros); y versiones
cuasi-casi seguras. La herramienta de \emph{promedio de funcionales gaussianos} de Nualart y Zheng
\cite{NualartZheng20} sistematiza el método.
\end{observacion}


\part{Extensiones}
\chapter{Browniano fraccionario y procesos gaussianos generales}\label{cap:fraccionario}

\begin{panorama}
El browniano fraccionario (fBm) $B^H$ es el proceso gaussiano autosimilar con incrementos
estacionarios más sencillo después del browniano. Para $H\ne1/2$ no es semimartingala,
de modo que el cálculo de Itô no se aplica. El cálculo de Malliavin, que solo necesita un
proceso isonormal, sí: la divergencia $\delta$ proporciona una integral estocástica natural, y su
relación con las integrales trayectoriales (Young, Stratonovich) se expresa con un término
de traza que involucra la derivada de Malliavin. Es el lenguaje de la volatilidad rugosa
del Capítulo~\ref{cap:esperanzas-condicionales}.
\end{panorama}

\section{Definición y propiedades}

\begin{definicion}
Para $H\in(0,1)$, el \emph{browniano fraccionario} de Hurst $H$ es el proceso
gaussiano centrado continuo con
\[
  R_H(s,t)=\E[B^H_sB^H_t]=\tfrac12\big(t^{2H}+s^{2H}-\abs{t-s}^{2H}\big).
\]
\end{definicion}

\begin{proposicion}
\begin{enumerate}[label=(\roman*)]
  \item \textup{Autosimilitud:} $(B^H_{ct})_t\law(c^HB^H_t)_t$. \textup{Incrementos estacionarios:}
  $\E(B^H_t-B^H_s)^2=\abs{t-s}^{2H}$.
  \item Las trayectorias son $\gamma$-Hölder para todo $\gamma<H$ y para ningún $\gamma>H$.
  \item Si $H>1/2$, los incrementos están positivamente correlacionados y $\sum_n\Cov(B^H_1,B^H_{n+1}-
  B^H_n)=\infty$ (memoria larga). Si $H<1/2$, están negativamente correlacionados (trayectorias
  rugosas).
  \item Para $H\ne\frac12$, $B^H$ no es semimartingala: su $p$-variación es 0 si $p>1/H$ e infinita si
  $p<1/H$.
\end{enumerate}
\end{proposicion}

\begin{proof}
(i) Directo de la covarianza. (ii) Kolmogorov con $\E\abs{B^H_t-B^H_s}^{2m}=
c_m\abs{t-s}^{2Hm}$, y la cota inferior por la ley del logaritmo iterado local. (iii) $\Cov(B^H_1,B^H_{n+1}-B^H_n)=
\frac12((n+1)^{2H}-2n^{2H}+(n-1)^{2H})\sim H(2H-1)n^{2H-2}$. (iv) Por autosimilitud y ergodicidad,
$\sum_i\abs{B^H_{(i+1)/n}-B^H_{i/n}}^p\approx n^{1-pH}\E\abs{B^H_1}^p$. Si $H<1/2$, la variación cuadrática es infinita; si
$H>1/2$, es nula, y una semimartingala continua con variación cuadrática nula sería de
variación finita, lo que contradice (iv) con $p=1$.
\end{proof}

\section{El espacio de Hilbert asociado}
Sea $\mathcal E$ el espacio de funciones escalonadas en $[0,T]$ y $\Hf$ la compleción respecto a $\ip{\ind_{[0,t]}}{\ind_{[0,s]}}_\Hf=
R_H(s,t)$. Entonces $B^H(\varphi)$, extendido por linealidad e isometría, es un proceso isonormal sobre
$\Hf$, y la Parte~II se aplica íntegramente.

\paragraph{El caso $H>1/2$.} Como $\partial_s\partial_tR_H(s,t)=\alpha_H\abs{t-s}^{2H-2}$, con $\alpha_H=H(2H-1)$,
\[
  \ip{\varphi}{\psi}_\Hf=\alpha_H\int_0^T\int_0^T\varphi(s)\psi(t)\abs{t-s}^{2H-2}\,\dd s\,\dd t,
\]
y $L^{1/H}([0,T])\subset\Hf$. Los elementos de $\Hf$ pueden ser distribuciones.

\paragraph{El caso $H<1/2$.} $\Hf$ es un espacio de funciones más pequeño, contenido en $L^2$, e igual a un
espacio de funciones con ``regularidad fraccionaria $\frac12-H$'' en un sentido preciso (vía el operador
$K_H^*$).

\paragraph{Representación de Volterra.} Existe un núcleo $K_H(t,s)$, con $K_H(t,s)=0$ para $s>t$, tal que
$B^H_t=\int_0^tK_H(t,s)\,\dd W_s$ para un browniano $W$. El operador $(K_H^*\varphi)(s)=\int_s^T\varphi(t)\partial_tK_H(t,s)\,\dd t$
(con interpretación adecuada) es una isometría de $\Hf$ en $L^2([0,T])$, y la derivada de Malliavin
respecto a $B^H$ se relaciona con la de $W$ por $D^WF=K_H^*D^{B^H}F$. Esto permite transportar
resultados del browniano.

\section{Integrales respecto al fBm}
Para $u$ un proceso adecuado, hay tres candidatos a ``$\int_0^Tu_t\,\dd B^H_t$'':
\begin{enumerate}
  \item La \emph{integral de divergencia} (o de Skorokhod) $\delta^{B^H}(u)$, con $\E\delta(u)=0$.
  \item La \emph{integral de Young} trayectorial, cuando $u$ es $\alpha$-Hölder con $\alpha+H>1$:
  \[
    \int_0^Tu_t\,\dd B^H_t=\lim_{\abs\pi\to0}\sum_iu_{t_i}\big(B^H_{t_{i+1}}-B^H_{t_i}\big).
  \]
  \item La \emph{integral simétrica o de Stratonovich}, definida por regularización: $\lim_{\varepsilon\to0}\frac1{2\varepsilon}\int_0^Tu_t(B^H_{t+\varepsilon}-
  B^H_{t-\varepsilon})\,\dd t$.
\end{enumerate}

\begin{teorema}[Relación entre integrales, $H>1/2$]
Si $H>1/2$ y $u\in\D^{1,2}(\Hf)$ es
suficientemente regular, la integral trayectorial (Young o Stratonovich) existe y
\[
  \int_0^Tu_t\circ\dd B^H_t=\delta^{B^H}(u)+\alpha_H\int_0^T\int_0^TD_su_t\,\abs{t-s}^{2H-2}\,\dd s\,\dd t.
\]
\end{teorema}

\begin{proof}[Esbozo]
Para la suma de Riemann, por la Proposición~\ref{prop:reglas-delta}(ii),
\[
  u_{t_i}\big(B^H_{t_{i+1}}-B^H_{t_i}\big)=\delta\big(u_{t_i}\ind_{(t_i,t_{i+1}]}\big)+\ip{Du_{t_i}}{\ind_{(t_i,t_{i+1}]}}_\Hf.
\]
El primer término converge a $\delta(u)$ por cerrabilidad. El segundo es
\[
  \alpha_H\sum_i\int_0^T\int_{t_i}^{t_{i+1}}D_su_{t_i}\abs{t-s}^{2H-2}\,\dd t\,\dd s,
\]
que converge al término de traza.
\end{proof}

Para $H=1/2$, $\alpha_H\abs{t-s}^{2H-2}$ se sustituye por $\delta(t-s)$ y el término de traza es $\frac12\int_0^TD_tu_t\,\dd t$,
la corrección de Itô--Stratonovich, recuperando $\int u\circ\dd W=\int u\,\dd W+\frac12\langle u,W\rangle$.

\begin{teorema}[Fórmula de Itô para la divergencia]
Si $f\in C^2$ con derivadas de crecimiento
moderado y $H>1/4$,
\[
  f(B^H_t)=f(0)+\delta^{B^H}\big(f'(B^H_\cdot)\ind_{[0,t]}\big)+H\int_0^tf''(B^H_s)s^{2H-1}\,\dd s.
\]
\end{teorema}

Para $H\le1/4$, el proceso $f'(B^H)\ind_{[0,t]}$ deja de estar en el dominio de $\delta$ en general, y se
necesita una noción extendida de divergencia; este umbral coincide con el de la teoría de rough
paths para procesos gaussianos (Capítulo~\ref{cap:rough-paths}).

\section{EDEs dirigidas por fBm}
Para $H>1/2$, la ecuación $\dd X_t=b(X_t)\,\dd t+\sigma(X_t)\,\dd B^H_t$, interpretada como integral de Young,
tiene solución única trayectorial (Lyons, Nualart--Răşcanu), y es Malliavin-diferenciable con
$D_sX_t=Y_tZ_s\sigma(X_s)$, la misma fórmula que en el Teorema~\ref{teo:derivada-ede}. La existencia de densidad regular
bajo condición de Hörmander fue probada por Baudoin y Hairer ($H>1/2$), y extendida a
$H>1/4$ con rough paths por Cass, Friz, Hairer, Litterer y Tindel~\cite{CassHairerLittererTindel15}. La dificultad es que la
matriz de Malliavin ya no se controla con la isometría de Itô, porque la norma de $\Hf$ no es local;
se necesitan \emph{desigualdades de interpolación} entre la norma de $\Hf$ y normas de Hölder.

\section{Procesos de Volterra y volatilidad rugosa}
El proceso de Riemann--Liouville $\hat W^H_t=\sqrt{2H}\int_0^t(t-s)^{H-1/2}\,\dd W_s$ es una variante del fBm
(sin incrementos estacionarios pero con la misma regularidad) que es \emph{adaptado a la filtración
browniana}. Por eso es el proceso preferido en volatilidad rugosa: la derivada de Malliavin
respecto a $W$ es explícita, $D_s\hat W^H_t=\sqrt{2H}(t-s)^{H-1/2}\ind_{\{s<t\}}$, y todas las fórmulas del Capítulo~\ref{cap:esperanzas-condicionales}
se aplican con derivadas respecto a un browniano estándar. La singularidad del núcleo en $s=t$
es exactamente lo que produce la explosión del skew en tiempo corto.

\chapter{Cálculo de Malliavin en el espacio de Poisson}\label{cap:poisson}

\begin{panorama}
Para procesos con saltos, el ruido gaussiano se sustituye por una medida aleatoria de
Poisson. Todavía hay una descomposición en caos, pero la derivada pasa a ser un \emph{operador
diferencia}: $D_zF$ es el cambio de $F$ al añadir un punto en $z$. Desarrollamos la teoría (Last,
Penrose, Peccati, Nualart--Vives) y señalamos las diferencias fundamentales con el caso
gaussiano: no hay regla de la cadena, y la desigualdad de Poincaré y las representaciones
de covarianza siguen valiendo, lo que basta para el método de Malliavin--Stein.
\end{panorama}

\section{Medidas aleatorias de Poisson}
Sea $(\mathbb X,\mathcal X,\lambda)$ un espacio de medida $\sigma$-finito. Una \emph{medida aleatoria de Poisson} $\eta$ con intensidad
$\lambda$ es una medida aleatoria con valores en $\N\cup\{\infty\}$ tal que, para conjuntos disjuntos $A_1,\dots,A_n$
con $\lambda(A_i)<\infty$, las variables $\eta(A_i)$ son independientes y $\eta(A_i)\sim\mathrm{Poisson}(\lambda(A_i))$. La medida
compensada es $\hat\eta=\eta-\lambda$. Para un proceso de Lévy puro salto, $\mathbb X=\R_+\times(\R\setminus\{0\})$ y $\lambda=\dd t\otimes\nu$,
con $\nu$ la medida de Lévy.

Sea $F=f(\eta)$ un funcional de la configuración de puntos. Escribimos $\eta+\delta_z$ para la
configuración con un punto añadido en $z$.

\begin{definicion}
El \emph{operador diferencia} (o derivada de adición) es
\[
  D_zF=f(\eta+\delta_z)-f(\eta),\qquad z\in\mathbb X.
\]
Las diferencias iteradas $D^n_{z_1,\dots,z_n}F$ se definen recursivamente y son simétricas.
\end{definicion}

\section{Caos de Poisson y fórmula de Stroock}

\begin{teorema}[Last--Penrose~\cite{LastPenrose17}]
Toda $F\in L^2(\P_\eta)$ se descompone como $F=\sum_{n\ge0}I_n(f_n)$,
donde $I_n$ es la integral múltiple respecto a $\hat\eta$ (fuera de las diagonales), con
\[
  f_n(z_1,\dots,z_n)=\frac1{n!}\E\big[D^n_{z_1,\dots,z_n}F\big],\qquad\E F^2=\sum_nn!\norm{f_n}^2_{L^2(\lambda^n)}.
\]
\end{teorema}

\begin{proof}[Esbozo]
La fórmula de Mecke, $\E\int h(\eta,z)\eta(\dd z)=\E\int h(\eta+\delta_z,z)\lambda(\dd z)$, que caracteriza a la
medida de Poisson, implica una fórmula de integración por partes. Para funcionales exponenciales
$F=\exp(-\int g\,\dd\eta)$ el desarrollo se calcula explícitamente, y se concluye por densidad. La fórmula
de Stroock es exacta y no requiere ninguna regularidad: $D^n$ siempre está definido.
\end{proof}

Como en el caso gaussiano, $F\in\Dom D$ si y solo si $\sum_nn\,n!\norm{f_n}^2<\infty$, y entonces
$D_zF=\sum_nnI_{n-1}(f_n(\cdot,z))$. Se definen análogamente $\delta=D^*$ (que extiende la integral respecto
a $\hat\eta$ de integrandos predecibles) y $L=-\delta D=-\sum nJ_n$.

\section{Diferencias con el caso gaussiano}
\begin{itemize}
  \item \textbf{No hay regla de la cadena.} $D_z\varphi(F)=\varphi(F+D_zF)-\varphi(F)$, que no es $\varphi'(F)D_zF$.
  Un desarrollo de Taylor da $D_z\varphi(F)=\varphi'(F)D_zF+R$, con $\abs R\le\frac12\norm{\varphi''}_\infty(D_zF)^2$: los
  términos de resto son la principal dificultad técnica.
  \item \textbf{$D$ no es una derivación}: $D_z(FG)=FD_zG+GD_zF+D_zFD_zG$.
  \item \textbf{Las leyes pueden tener átomos.} Todo funcional de un número finito de puntos tiene
  átomos, de modo que no hay un análogo directo del criterio de Bouleau--Hirsch. Para
  densidades de funcionales de procesos de Lévy se usan otras derivadas: las de Bismut,
  que mueven los tiempos o las alturas de los saltos (cálculo de Malliavin ``por perturbación
  de saltos''; Bichteler--Gravereaux--Jacod, Bally--Clément).
  \item \textbf{La hipercontractividad falla} en su forma gaussiana: el semigrupo de Ornstein--Uhlenbeck
  de Poisson (vía la fórmula de Mehler de Last--Peccati--Schulte, que ``adelgaza''
  los puntos con probabilidad $e^{-t}$ y añade puntos independientes con intensidad $(1-e^{-t})\lambda$)
  no es hipercontractivo. Los momentos dentro de un caos fijo no son equivalentes en
  general.
\end{itemize}

\section{Lo que sí funciona}

\begin{teorema}[Poincaré y covarianzas de Poisson]
Para $F,G\in L^2$,
\[
  \Var F\le\E\int\abs{D_zF}^2\,\lambda(\dd z),\qquad\Cov(F,G)=\E\int D_zF\,(-D_zL^{-1}G)\,\lambda(\dd z).
\]
\end{teorema}

\begin{proof}
La primera es la misma prueba por caos que el Teorema~\ref{teo:poincare}. La segunda, la de
la Proposición~\ref{prop:covarianzas}, usando solo la dualidad.
\end{proof}

La representación de covarianzas, combinada con la expansión de Taylor anterior, basta
para desarrollar el método de Malliavin--Stein en el espacio de Poisson (Peccati, Solé, Taqqu,
Utzet, 2010). Last, Peccati y Schulte~\cite{LastPeccatiSchulte16} obtuvieron las \emph{desigualdades de Poincaré de segundo
orden}, que acotan la distancia a la normal solo en términos de $\E\abs{D_zF}^4$ y $\E\abs{D^2_{z,w}F}^4$. Su gran
ventaja es que se aplican a funcionales \emph{geométricos} de procesos de puntos: longitud del árbol
recubridor mínimo, volumen intrínseco de la envolvente convexa, número de aristas en grafos
aleatorios geométricos. Para todos estos se obtienen TCLs con la tasa óptima $\lambda(\mathbb X)^{-1/2}$.

\begin{ejemplo}[Estabilización]
Sea $\eta_t$ un proceso de Poisson de intensidad $t$ en $[0,1]^d$ y
$F_t=\sum_{x\in\eta_t}\xi(x,\eta_t)$, donde $\xi$ es una \emph{puntuación local} (por ejemplo, la distancia al vecino
más próximo). Si $\xi$ se ``estabiliza'', es decir, $D_zF_t$ depende solo de los puntos en una bola
aleatoria de radio con colas exponenciales, entonces $D_zF_t$ y $D^2F_t$ tienen momentos acotados
uniformemente, $D^2_{z,w}F_t$ decae rápidamente en $\abs{z-w}$, y la desigualdad de Last--Peccati--Schulte
da $d_W\big(\frac{F_t-\E F_t}{\sqrt{\Var F_t}},N\big)\le Ct^{-1/2}$.
\end{ejemplo}


\part{Estado del arte}
\chapter{El método de Malliavin--Stein}\label{cap:malliavin-stein}

\begin{panorama}
En 2005, Nualart y Peccati~\cite{NualartPeccati05} hallaron un fenómeno sorprendente: una sucesión de
variables en un caos de Wiener fijo, normalizada, converge a la normal \emph{si y solo si su
cuarto momento converge a} 3. Poco después, Nourdin y Peccati (2009) mostraron que
combinando el cálculo de Malliavin con el método de Stein se obtienen cotas \emph{cuantitativas}
explícitas. El ``método de Malliavin--Stein'' es hoy una de las áreas más activas del cálculo
de Malliavin, con aplicaciones a la estadística de procesos de memoria larga, la geometría
aleatoria, las EDPEs (Capítulo~\ref{cap:anderson}), las ondas aleatorias y las redes neuronales anchas.
\end{panorama}

\section{El método de Stein}

\begin{lema}[Stein]\label{lem:stein-ecuacion}
Sea $N\sim N(0,1)$ y $h:\R\to\R$ medible acotada. La ecuación de Stein
\[
  f'(x)-xf(x)=h(x)-\E h(N)
\]
tiene una solución acotada $f_h(x)=e^{x^2/2}\int_{-\infty}^x(h(y)-\E h(N))e^{-y^2/2}\,\dd y$. Si $h=\ind_A$ con $A$
boreliano, entonces $\norm{f_h}_\infty\le\sqrt{\pi/2}$ y $\norm{f_h'}_\infty\le2$. Si $h$ es 1-Lipschitz, $\norm{f_h'}_\infty\le\sqrt{2/\pi}$ y
$\norm{f_h''}_\infty\le2$. En consecuencia, para toda variable $X$,
\[
  d_{TV}(X,N)=\sup_A\abs{\P(X\in A)-\P(N\in A)}\le\sup_{f\in\F_{TV}}\abs{\E[f'(X)-Xf(X)]},
\]
con $\F_{TV}=\{f\in C^1:\ \norm f_\infty\le\sqrt{\pi/2},\ \norm{f'}_\infty\le2\}$, y análogamente para la distancia de
Wasserstein $d_W$, con $\F_W=\{\norm{f'}_\infty\le\sqrt{2/\pi}\}$.
\end{lema}

\begin{proof}[Esbozo]
Que $f_h$ resuelve la ecuación es una verificación directa: $(e^{-x^2/2}f_h)'=(h-\E h(N))e^{-x^2/2}$.
Las cotas se obtienen con las estimaciones de Mills, $\int_x^\infty e^{-y^2/2}\,\dd y\le\min(\sqrt{\pi/2},1/x)e^{-x^2/2}$,
usando que $\int_\R(h-\E h(N))e^{-y^2/2}=0$ para escribir $f_h$ con la cola izquierda o la derecha según
el signo de $x$. Para $h$ indicador no continuo, $f_h$ es solo Lipschitz; se aproxima. Finalmente,
evaluando la ecuación en $X$ y tomando esperanzas, $\P(X\in A)-\P(N\in A)=\E[f_h'(X)-
Xf_h(X)]$.
\end{proof}

\begin{intuicion}
El método de Stein convierte la pregunta ``¿está $X$ cerca de $N$?'' en ``¿satisface $X$
aproximadamente la identidad $\E[Xf(X)]=\E[f'(X)]$?''. Y el cálculo de Malliavin es
precisamente una máquina para calcular $\E[Xf(X)]$: por la representación de covarianzas,
\[
  \E[Xf(X)]=\E\big[f'(X)\ip{DX}{-DL^{-1}X}_\Hf\big].
\]
Así, $X$ es casi normal si $\ip{DX}{-DL^{-1}X}\approx1$. Los dos métodos encajan como una llave
en su cerradura.
\end{intuicion}

\section{La cota fundamental}

\begin{teorema}[Nourdin--Peccati, 2009]\label{teo:nourdin-peccati}
Sea $F\in\D^{1,2}$ con $\E F=0$ y $\E F^2=1$. Entonces
\[
  d_{TV}(F,N)\le2\,\E\Big|1-\ip{DF}{-DL^{-1}F}_\Hf\Big|\le2\sqrt{\Var\big(\ip{DF}{-DL^{-1}F}_\Hf\big)}.
\]
\end{teorema}

\begin{proof}
Para $f\in\F_{TV}$, por la Proposición~\ref{prop:covarianzas} y la regla de la cadena,
\[
  \E[Ff(F)]=\E\ip{Df(F)}{-DL^{-1}F}=\E\big[f'(F)\ip{DF}{-DL^{-1}F}\big].
\]
Así $\abs{\E[f'(F)-Ff(F)]}=\abs{\E[f'(F)(1-\ip{DF}{-DL^{-1}F})]}\le2\E\abs{1-\ip{DF}{-DL^{-1}F}}$. La segunda
desigualdad es Cauchy--Schwarz, pues $\E\ip{DF}{-DL^{-1}F}=\E F^2=1$. (Estrictamente, la
regla de la cadena se aplica a $f$ Lipschitz con la versión de la Proposición~\ref{prop:cadena}, lo que requiere
que $F$ tenga densidad, o un argumento de aproximación.)
\end{proof}

\section{El teorema del cuarto momento}
Si $F=I_q(f)$, entonces $L^{-1}F=-\frac1qF$, y por tanto $\ip{DF}{-DL^{-1}F}=\frac1q\norm{DF}^2_\Hf$.

\begin{lema}\label{lem:cuarto-momento}
Sea $F=I_q(f)$ con $f\in\Hf^{\odot q}$, $q\ge2$, y $\E F^2=q!\norm f^2=1$. Entonces
\begin{align*}
  \Var\Big(\frac1q\norm{DF}^2_\Hf\Big)&=\frac1{q^2}\sum_{r=1}^{q-1}r^2(r!)^2\binom qr^4(2q-2r)!\,\norm{f\tilde\otimes_rf}^2,\\
  \E F^4-3&=\frac3q\sum_{r=1}^{q-1}r\,(r!)^2\binom qr^4(2q-2r)!\,\norm{f\tilde\otimes_rf}^2.
\end{align*}
En consecuencia,
\[
  \Var\Big(\frac1q\norm{DF}^2_\Hf\Big)\le\frac{q-1}{3q}\big(\E F^4-3\big).
\]
\end{lema}

\begin{proof}[Esbozo]
Por el Teorema~\ref{teo:derivada-caos}, $D_tF=qI_{q-1}(f(\cdot,t))$. Por la fórmula del producto (Teorema~\ref{teo:producto})
aplicada a $I_{q-1}(f(\cdot,t))^2$ e integrando en $t$ (lo que añade una contracción),
\[
  \frac1q\norm{DF}^2=q\sum_{r=1}^q(r-1)!\binom{q-1}{r-1}^2I_{2q-2r}(f\tilde\otimes_rf).
\]
El término $r=q$ es la constante $q!\norm f^2=1$, y la varianza se calcula con la isometría;
reescribiendo $q(r-1)!\binom{q-1}{r-1}^2=\frac rq\cdot r!\binom qr^2$ se obtiene la primera fórmula. Para el cuarto momento,
$\E F^4=\E[F\cdot F^3]=3\E[F^2\cdot\frac1q\norm{DF}^2]$ (dualidad y regla de la cadena: $\E[F\cdot F^3]=
\E\ip{D(F^3)}{-DL^{-1}F}=3\E[F^2\frac1q\norm{DF}^2]$), y se desarrolla $F^2$ con la fórmula del producto; el
resultado es la segunda fórmula. Comparando término a término, el cociente entre el coeficiente
de $\norm{f\tilde\otimes_rf}^2$ en la varianza y en $\E F^4-3$ es $\frac{r^2/q^2}{3r/q}=\frac r{3q}\le\frac{q-1}{3q}$, lo que da la desigualdad.
\end{proof}

\begin{teorema}[Teorema del cuarto momento]\label{teo:cuarto-momento}
Sea $q\ge2$ y $F_n=I_q(f_n)$ con $\E F_n^2\to1$. Son
equivalentes:
\begin{enumerate}[label=(\roman*)]
  \item $F_n\tolaw N(0,1)$;
  \item $\E F_n^4\to3$;
  \item $\norm{f_n\otimes_rf_n}_{\Hf^{\otimes(2q-2r)}}\to0$ para $r=1,\dots,q-1$;
  \item $\norm{DF_n}^2_\Hf\to q$ en $L^2(\Omega)$.
\end{enumerate}
Además, si $\E F^2=1$, se tiene la cota cuantitativa
\[
  d_{TV}(F,N)\le2\sqrt{\frac{q-1}{3q}\big(\E F^4-3\big)}.
\]
\end{teorema}

\begin{proof}
La cota es el Teorema~\ref{teo:nourdin-peccati} combinado con el Lema~\ref{lem:cuarto-momento}; da (ii)$\Rightarrow$(i). (i)$\Rightarrow$(ii): por
el Teorema~\ref{teo:normas-caos} (hipercontractividad), $\sup_n\E\abs{F_n}^p<\infty$ para todo $p$, luego hay integrabilidad
uniforme de $F_n^4$ y $\E F_n^4\to\E N^4=3$. (ii)$\Leftrightarrow$(iii): por la fórmula del cuarto cumulante del
Lema~\ref{lem:cuarto-momento}, y porque $\norm{f\tilde\otimes_rf}$ y $\norm{f\otimes_rf}$ son comparables (identidad combinatoria). (iii)$\Leftrightarrow$(iv):
por la fórmula de la varianza (véase también~\cite{NualartOrtizLatorre08}).
\end{proof}

\begin{intuicion}
En el caos $q\ge2$, el cuarto cumulante $\kappa_4(F)=\E F^4-3(\E F^2)^2$ es siempre $\ge0$ (Lema~\ref{lem:cuarto-momento}):
las variables de caos tienen colas más pesadas que la normal. La normal es el único
``mínimo'' posible, y alcanzarlo asintóticamente solo es posible si el núcleo $f_n$ se ``dispersa'',
es decir, si sus contracciones se anulan. Es un \emph{principio de rigidez}. Notablemente, la
condición (iv) dice que la ``varianza local'' $\frac1q\norm{DF}^2$ se vuelve determinista.
\end{intuicion}

\begin{teorema}[Peccati--Tudor~\cite{PeccatiTudor05}]
Sea $F_n=(F^1_n,\dots,F^d_n)$ con $F^i_n=I_{q_i}(f^i_n)$ y $\E[F^i_nF^j_n]\to
C_{ij}$. Entonces $F_n\tolaw N(0,C)$ si y solo si cada componente converge a su marginal
gaussiana.
\end{teorema}

Es decir: en el caos de Wiener, \emph{la convergencia componente a componente implica la
convergencia conjunta}. Esto falla para variables generales.

\section{Aplicación: variación cuadrática del browniano fraccionario}
Sea $B^H$ un fBm y $X_k=n^H(B^H_{(k+1)/n}-B^H_{k/n})$, $k=0,\dots,n-1$, que forman una sucesión
gaussiana estacionaria con $\E X_k^2=1$ y covarianza
\[
  \rho(k)=\tfrac12\big(\abs{k+1}^{2H}-2\abs k^{2H}+\abs{k-1}^{2H}\big)\sim H(2H-1)\abs k^{2H-2}.
\]
Sea $V_n=\frac1{\sqrt n}\sum_{k<n}(X_k^2-1)=\frac1{\sqrt n}\sum_kH_2(X_k)$, que vive en el segundo caos.

\begin{teorema}[Breuer--Major; Nourdin--Peccati]
Si $H<3/4$, $V_n\tolaw N(0,\sigma_H^2)$ con
$\sigma_H^2=2\sum_{k\in\Z}\rho(k)^2$, y
\[
  d_{TV}\Big(\frac{V_n}{\sqrt{\Var V_n}},N\Big)\le C\cdot
  \begin{cases}
    n^{-1/2}, & H\in(0,\frac58),\\
    n^{-1/2}\log^{3/2}n, & H=\frac58,\\
    n^{4H-3}, & H\in(\frac58,\frac34).
  \end{cases}
\]
Si $H>3/4$, $n^{1-2H}V_n$ converge a una variable de \textup{Rosenblatt} (no gaussiana, en el segundo
caos).
\end{teorema}

\begin{proof}[Esbozo]
$V_n=I_2(f_n)$ con $f_n=\frac1{\sqrt n}\sum_ke_k\otimes e_k$, donde $e_k\in\Hf$ con $\ip{e_k}{e_l}=\rho(k-l)$. Entonces
$\Var V_n=\frac2n\sum_{k,l<n}\rho(k-l)^2\to\sigma_H^2$, finito si y solo si $\sum\rho^2<\infty$, es decir, $H<3/4$. Además
\[
  \norm{f_n\otimes_1f_n}^2=\frac1{n^2}\sum_{i,j,k,l<n}\rho(i-j)\rho(j-k)\rho(k-l)\rho(l-i)\le\frac Cn\Big(\sum_{\abs k<n}\abs{\rho(k)}^{4/3}\Big)^3
\]
por la desigualdad de Young para convoluciones. Con $\rho(k)\sim k^{2H-2}$, la suma es acotada si
$H<5/8$, crece como $\log n$ si $H=5/8$ y como $n^{(8H-5)/3}$ si $H>5/8$. Por el Teorema~\ref{teo:cuarto-momento} (en
$q=2$, $\kappa_4=48\norm{f\otimes_1f}^2$ y la cota de distancia es $\propto\sqrt{\kappa_4}$) se obtienen las tasas. Para $H>3/4$,
la renormalización $n^{1-2H}V_n$ hace que $f_n$ converja en $\Hf^{\otimes2}$ a un núcleo no nulo, cuya integral
doble es la variable de Rosenblatt.
\end{proof}

Este resultado tiene una consecuencia estadística inmediata: el estimador de $H$ basado en
la variación cuadrática es asintóticamente normal si y solo si $H<3/4$.

\section{Refinamientos}

\paragraph{Tasas óptimas.} Nourdin y Peccati (2015) probaron que, para $F_n$ en un caos fijo, $d_{TV}(F_n,N)\asymp
\max\{\abs{\E F_n^3},\ \E F_n^4-3\}$: la cota $\sqrt{\kappa_4}$ del Teorema~\ref{teo:cuarto-momento} no es óptima, y la tasa correcta es lineal
en los cumulantes.

\paragraph{Entropía.} Nourdin, Peccati y Swan (2014) mostraron que el teorema del cuarto momento
vale también en \emph{entropía relativa}, una noción de convergencia más fuerte, y con ello en total
variación de las densidades.

\paragraph{Desigualdades de Poincaré de segundo orden.} Para $F$ general (no en un caos fijo),
Nourdin, Peccati y Reinert~\cite{NourdinPeccatiReinert09} probaron
\[
  d_{TV}(F,N)\le\frac{\sqrt{10}}{2\sigma^2}\Big(\E\norm{D^2F}^4_{\mathrm{op}}\Big)^{1/4}\Big(\E\norm{DF}^4_\Hf\Big)^{1/4},
\]
donde $\norm{D^2F}_{\mathrm{op}}$ es la norma de operador de la segunda derivada como forma bilineal en $\Hf$. Esta
cota solo requiere estimar $DF$ y $D^2F$, y es la herramienta estándar para funcionales no lineales
(EDPEs, redes neuronales con anchura grande, estadísticos de campos gaussianos).

\paragraph{Universalidad.} Nourdin, Peccati y Reinert (2010) probaron que el caos de Wiener es \emph{universal}:
si una sucesión de sumas homogéneas multilineales $Q_n(\mathbf G)=\sum f_n(i_1,\dots,i_q)G_{i_1}\cdots G_{i_q}$ en
gaussianas satisface un TCL, entonces lo satisface $Q_n(\mathbf X)$ para cualesquiera $X_i$ independientes
centradas de varianza unidad (bajo una condición de influencia pequeña). La prueba combina
el teorema del cuarto momento con un principio de Lindeberg.

\paragraph{Más allá de la normal.} El método se extiende a aproximaciones por la ley Gamma (segundo
caos), por leyes en el segundo caos en general (Azmoodeh, Peccati, Poly), y al marco de \emph{triples
de Markov} (Ledoux; Azmoodeh, Campese, Poly), donde el operador $L$ es el generador de un
semigrupo de difusión con espectro discreto (por ejemplo, los polinomios de Jacobi o Laguerre),
lo que produce teoremas del cuarto momento para ``caos'' de otras medidas. Existe también un
teorema del cuarto momento en \emph{probabilidad libre} (Kemp, Nourdin, Peccati, Speicher, 2012),
con la ley semicircular como límite. Véase la reseña~\cite{AzmoodehPeccatiYang21}.

\paragraph{Geometría aleatoria.} El método de Malliavin--Stein ha sido decisivo en el estudio de los
conjuntos nodales de ondas aleatorias (longitud de las líneas nodales de las ondas planas de
Berry, autofunciones aleatorias en la esfera y el toro): Marinucci, Peccati, Rossi y Wigman
obtuvieron TCLs y descubrieron la \emph{cancelación de Berry}, en la que la proyección en el segundo
caos se anula y el comportamiento lo dicta el cuarto caos.

\chapter{Cálculo de Malliavin y rough paths}\label{cap:rough-paths}

\begin{panorama}
La teoría de rough paths de Lyons (1998) resuelve ecuaciones diferenciales dirigidas por
señales irregulares \emph{trayectoria a trayectoria}, sin martingalas. Para procesos gaussianos, la
solución es entonces un funcional continuo de un \emph{levantamiento} del ruido, y el cálculo
de Malliavin se puede hacer en sentido fuerte: la solución es Fréchet-diferenciable en las
direcciones de Cameron--Martin. Presentamos los ingredientes básicos, el teorema de
Hörmander para EDRs gaussianas (Cass--Friz, Cass--Hairer--Litterer--Tindel) y algunos
desarrollos recientes.
\end{panorama}

\section{Rough paths en pocas palabras}
Sea $\alpha\in(\frac13,\frac12]$. Un \emph{rough path} $\alpha$-Hölder en $\R^m$ es un par $\mathbf X=(X,\mathbb X)$ con $X:[0,T]\to\R^m$ y
$\mathbb X:[0,T]^2\to\R^{m\times m}$ tales que
\[
  \abs{X_t-X_s}\le C\abs{t-s}^\alpha,\qquad\abs{\mathbb X_{s,t}}\le C\abs{t-s}^{2\alpha},\qquad\mathbb X_{s,t}-\mathbb X_{s,u}-\mathbb X_{u,t}=X_{s,u}\otimes X_{u,t}
\]
(\emph{relación de Chen}), con $X_{s,t}=X_t-X_s$. Heurísticamente, $\mathbb X_{s,t}=\int_s^tX_{s,r}\otimes\dd X_r$: es la información
de segundo orden que la trayectoria por sí sola no determina cuando $\alpha\le\frac12$.

\begin{lema}[Lema de cosido, \textup{sewing lemma}]\label{lem:sewing}
Sea $\Xi_{s,t}$ una función de dos parámetros con
$\abs{\Xi_{s,t}-\Xi_{s,u}-\Xi_{u,t}}\le C\abs{t-s}^\theta$ para $s\le u\le t$, con $\theta>1$. Entonces existe una única función
aditiva $\mathcal I_t$ con $\abs{\mathcal I_t-\mathcal I_s-\Xi_{s,t}}\le C'\abs{t-s}^\theta$, dada por $\mathcal I_t-\mathcal I_s=\lim_{\abs\pi\to0}\sum_{[u,v]\in\pi}\Xi_{u,v}$.
\end{lema}

\begin{proof}[Esbozo]
Para particiones diádicas, quitar un punto de una partición de $n$ intervalos cambia
la suma en a lo sumo $C(\frac{2\abs{t-s}}{n})^\theta$. Eliminando puntos sucesivamente hasta llegar a $\{s,t\}$, la
diferencia total está acotada por $C\abs{t-s}^\theta\sum_nn^{-\theta}<\infty$, porque $\theta>1$. Esto da existencia del
límite y la estimación; la unicidad se sigue de que una función aditiva con $\abs{f_{s,t}}\le C\abs{t-s}^\theta$,
$\theta>1$, es nula.
\end{proof}

Con $\Xi_{s,t}=Y_sX_{s,t}+Y_s'\mathbb X_{s,t}$, para $Y$ \emph{controlado} por $X$ (es decir, $Y_{s,t}=Y_s'X_{s,t}+O(\abs{t-s}^{2\alpha})$),
se tiene $\theta=3\alpha>1$ y el lema define la \emph{integral rugosa} $\int Y\,\dd\mathbf X$. Esto permite resolver
\begin{equation}\label{eq:edr}
  \dd Y_t=V_0(Y_t)\,\dd t+V(Y_t)\,\dd\mathbf X_t
\end{equation}
por punto fijo, y la solución depende de manera \emph{localmente Lipschitz} de $\mathbf X$ en la métrica de
rough paths (\emph{teorema de continuidad de Lyons}, o teorema del límite universal).

\begin{ejemplo}[Browniano]
$\mathbb X^{\text{Itô}}_{s,t}=\int_s^tW_{s,r}\otimes\dd W_r$ y $\mathbb X^{\text{Strat}}_{s,t}=\mathbb X^{\text{Itô}}_{s,t}+\frac12(t-s)I$ son dos rough
paths distintos sobre la misma trayectoria; con el primero,~\eqref{eq:edr} es la EDE de Itô y, con el
segundo, la de Stratonovich. Ambos existen c.s.\ (por Kolmogorov en dos parámetros) para todo
$\alpha<\frac12$.
\end{ejemplo}

\section{Levantamientos gaussianos}

\begin{teorema}[Friz--Victoir]
Sea $X=(X^1,\dots,X^m)$ un proceso gaussiano centrado
continuo con componentes independientes, cuya covarianza tiene $\rho$-variación bidimensional
finita con $\rho<2$. Entonces $X$ se levanta canónicamente a un rough path geométrico $\mathbf X$ (límite
de los levantamientos de aproximaciones lineales a trozos) de $p$-variación finita para todo
$p>2\rho$. Para el browniano fraccionario, $\rho=\frac1{2H}$, de modo que la condición es $H>\frac14$.
\end{teorema}

Para $H\le1/4$, las aproximaciones lineales a trozos de $\int B^{H,1}\,\dd B^{H,2}$ divergen (la varianza del
área de Lévy explota), en consonancia con la fórmula de Itô para la divergencia del Capítulo~\ref{cap:fraccionario}.

\section{Diferenciabilidad de Malliavin trayectorial}

\begin{teorema}
Sea $X$ gaussiano como antes, con espacio de Cameron--Martin $\Hc$, y
supongamos que $\Hc$ se inyecta en el espacio de trayectorias de $q$-variación finita con $\frac1p+\frac1q>1$
\textup{(}regularidad complementaria de Young\textup{)}. Sea $Y=\Psi(\mathbf X)$ la solución de~\eqref{eq:edr} con $V\in C_b^\infty$.
Entonces, para cada trayectoria, $h\mapsto\Psi(T_h\mathbf X)$ es Fréchet-diferenciable en $\Hc$, donde $T_h\mathbf X$ es
el rough path trasladado, y
\[
  \ip{DY_t}{h}_\Hc=\int_0^tJ_{t\leftarrow s}V(Y_s)\,\dd h_s,
\]
con $J_{t\leftarrow s}$ el jacobiano del flujo, solución de la ecuación linealizada. En particular, $Y_t\in\D^{1,2}_{\mathrm{loc}}$
y, si $X$ es el fBm con $H>1/4$, la condición de Young complementaria se satisface.
\end{teorema}

Es la misma fórmula que en el Teorema~\ref{teo:derivada-ede}, pero ahora con una integral de Young respecto
a $h$ y válida $\omega$ a $\omega$. Para el fBm, $\Hc\subset C^{q\text{-var}}$ para todo $q>(H+\frac12)^{-1}$, y $\frac1p+\frac1q>1$ se reduce a
$H>\frac14$.

\section{El teorema de Hörmander para EDRs gaussianas}

\begin{teorema}[Cass--Friz 2010; Cass--Hairer--Litterer--Tindel 2015]
Sea $X$ un proceso
gaussiano que satisface las hipótesis anteriores más una condición de no degeneración \textup{(}no
determinismo local\textup{)}, por ejemplo el fBm con $H\in(\frac14,\frac12]$. Si los campos $V_0,V_1,\dots,V_m$
satisfacen la condición de Hörmander en $y_0$, entonces para todo $t>0$:
\begin{enumerate}[label=(\roman*)]
  \item la ley de $Y_t$ es absolutamente continua (Cass--Friz~\cite{CassFriz10});
  \item su densidad es $C^\infty$ (Cass--Hairer--Litterer--Tindel~\cite{CassHairerLittererTindel15}).
\end{enumerate}
\end{teorema}

\begin{proof}[Ideas de la prueba]
Hay tres obstáculos respecto al caso browniano, y cada uno requirió una
idea nueva.

\emph{(a) Sustituto del lema de Norris.} No hay martingalas, así que el argumento probabilístico
del Lema~\ref{lem:norris} no sirve. Se sustituye por un \emph{lema de Norris determinista} basado en la \emph{rugosidad
verdadera} (Hairer--Pillai): una trayectoria $X$ es $\theta$-rugosa si, en todo intervalo y en toda dirección,
no puede ser ``demasiado regular''. Para trayectorias rugosas, si $\int\phi\,\dd X+\int\psi\,\dd t$ es pequeño,
entonces $\phi$ y $\psi$ son pequeños (unicidad cuantitativa de la descomposición tipo Doob--Meyer).
La rugosidad se verifica c.s.\ para procesos gaussianos con no determinismo local.

\emph{(b) Integrabilidad del jacobiano.} Para una EDE, $J$ tiene todos los momentos por BDG--Gronwall.
Para una EDR, la cota natural es $\abs J\le\exp(C\norm{\mathbf X}^p_{p\text{-var}})$, y $\norm{\mathbf X}^p$ con $p>2$ \emph{no}
tiene colas gaussianas: la exponencial no es integrable. Cass, Litterer y Lyons resolvieron este
problema con las \emph{particiones voraces}: el jacobiano está controlado por $\exp(CN(\mathbf X))$, donde $N$
es el número de intervalos en que la $p$-variación del rough path supera un umbral, y $N$ sí tiene
colas gaussianas por el teorema de Borell aplicado a la traslación de Cameron--Martin.

\emph{(c) La matriz de Malliavin.} La norma de $\Hc$ no es local en tiempo, de modo que $\ip{DY}{DY}_\Hc$
no se escribe como una integral en $t$. Se necesitan \emph{desigualdades de interpolación} que acoten
inferiormente $\norm f_\Hc$ en términos de normas locales de $f$; para el fBm con $H<1/2$, $\norm f_\Hc\gtrsim\norm f_{L^2}$
en cierto sentido, lo que permite reducir al argumento de los corchetes.
\end{proof}

\section{Desarrollos recientes}
\begin{itemize}
  \item \textbf{Ecuaciones diferenciales estocásticas rugosas.} Friz, Hocquet y Lê (2021) desarrollaron
  una teoría para ecuaciones dirigidas simultáneamente por un browniano (tratado
  con Itô) y un rough path (tratado trayectorialmente), $\dd Y=b\,\dd t+\sigma\,\dd W+f\,\dd\mathbf X$. Son
  el lenguaje natural para filtrado no lineal robusto, volatilidad local-estocástica y control
  estocástico con ruido común. Su análisis de Malliavin combina ambos mundos.
  \item \textbf{TCL trayectoriales y desviaciones moderadas.} Se han obtenido TCLs y principios
  de desviaciones moderadas para EDPEs con ruido multiplicativo vía rough paths, donde
  el cálculo de Malliavin pathwise permite linealizar alrededor del límite determinista.
  \item \textbf{Densidades en régimen rugoso.} Inahama~\cite{Inahama14} y otros estudiaron la diferenciabilidad
  de Malliavin y los desarrollos asintóticos (Laplace, Watanabe) de soluciones de EDRs, con
  aplicaciones a núcleos de calor en geometría sub-riemanniana.
  \item \textbf{Signaturas.} La signatura de un rough path, la sucesión de todas sus integrales iteradas,
  determina la trayectoria salvo reparametrizaciones (Hambly--Lyons; Boedihardjo, Geng,
  Lyons, Yang). Es la base de métodos de aprendizaje automático para series temporales, y
  su esperanza bajo la medida de Wiener se calcula con las mismas herramientas de caos de
  Wiener del Capítulo~\ref{cap:caos}.
\end{itemize}

\chapter{Malliavin en estructuras de regularidad}\label{cap:estructuras}

\begin{panorama}
La teoría de estructuras de regularidad de Hairer~\cite{Hairer14} da sentido a EDPEs \emph{singulares},
como $\Phi^4_3$, KPZ o el PAM en dimensión 2 y 3, en las que aparecen productos de distribuciones
que solo existen tras una \emph{renormalización}. La parte analítica de la teoría es
determinista; la probabilística consiste en construir el \emph{modelo} a partir del ruido y probar
sus estimaciones estocásticas. El cálculo de Malliavin interviene de tres maneras: (1)
la hipercontractividad reduce las estimaciones a segundos momentos; (2) desde 2021, el
programa de Otto y colaboradores sustituye los diagramas de Feynman por la \emph{desigualdad
de spectral gap} y la derivada de Malliavin, con una prueba que funciona en todo el régimen
subcrítico; (3) el cálculo de Malliavin da densidades para soluciones de EDPEs singulares.
\end{panorama}

\section{EDPEs singulares}
Consideramos, en el toro $\mathbb T^d$, el \emph{modelo parabólico de Anderson generalizado} y la ecuación $\Phi^4$:
\[
  \partial_tu=\Delta u+g(u)\xi\quad\text{(gPAM, $\xi$ ruido blanco espacial en $d=2$)},\qquad\partial_t\phi=\Delta\phi-\phi^3+\xi\quad(\Phi^4_d).
\]
En gPAM con $d=2$, $\xi\in C^{-1-\kappa}$ para todo $\kappa>0$, y la teoría de Schauder sugiere $u\in C^{1-\kappa}$. El
producto $g(u)\xi$ involucra una función $C^{1-\kappa}$ por una distribución $C^{-1-\kappa}$: la suma de regularidades
es $-2\kappa<0$ y el producto no está definido (criterio de Young). En $\Phi^4_3$, $\phi\in C^{-1/2-\kappa}$ y $\phi^3$ es aún
peor. Si se regulariza el ruido, $\xi_\varepsilon=\xi*\rho_\varepsilon$, las soluciones $u_\varepsilon$ \emph{no convergen}; sí convergen las de
las ecuaciones \emph{renormalizadas}
\[
  \partial_tu_\varepsilon=\Delta u_\varepsilon+g(u_\varepsilon)\big(\xi_\varepsilon-C_\varepsilon g'(u_\varepsilon)\big),\qquad C_\varepsilon\sim\frac1{2\pi}\abs{\log\varepsilon},
\]
y el límite no depende del regularizador. Una ecuación es \emph{subcrítica} si, a escalas pequeñas, la
no linealidad es perturbativa respecto a la parte lineal; las tres ecuaciones anteriores lo son, y
$\Phi^4_4$, KPZ en $d\ge2$ o el PAM en $d=4$ no.

\section{La arquitectura de la teoría}
Muy esquemáticamente:
\begin{enumerate}
  \item Se construye un conjunto finito de \emph{símbolos} $\tau$ (árboles), que representan los objetos que
  aparecen al iterar la ecuación: $\Xi$ (el ruido), $\mathcal I(\Xi)$ (la solución de la ecuación del calor con
  ese ruido), $\mathcal I(\Xi)\Xi$, etc. Cada uno tiene una \emph{homogeneidad} $\abs\tau\in\R$.
  \item Un \emph{modelo} asigna a cada símbolo y cada punto base $x$ una distribución $\Pi_x\tau$, que es el
  ``objeto $\tau$ recentrado en $x$'', y que debe satisfacer
  \begin{equation}\label{eq:modelo}
    \Big|\big\langle\Pi_x\tau,\varphi^\lambda_x\big\rangle\Big|\lesssim\lambda^{\abs\tau},\qquad\varphi^\lambda_x(y)=\lambda^{-d}\varphi((y-x)/\lambda),
  \end{equation}
  junto con aplicaciones $\Gamma_{xy}$ que relacionan los recentrados en puntos distintos.
  \item El \emph{teorema de reconstrucción} permite construir distribuciones globales a partir de desarrollos
  locales en los símbolos. La EDPE se reformula como un punto fijo en un espacio de
  \emph{distribuciones modeladas}, que se resuelve de manera determinista y continua en el modelo.
  \item Lo probabilístico: dado el ruido regularizado, se construye el modelo canónico, se renormaliza
  (con el grupo de renormalización de Bruned--Hairer--Zambotti~\cite{BrunedHairerZambotti19}, que actúa sobre los
  modelos mediante álgebras de Hopf de árboles) y se prueba que el modelo renormalizado
  converge y satisface~\eqref{eq:modelo} en momentos:
  \begin{equation}\label{eq:modelo-momentos}
    \E\Big|\big\langle\Pi_x\tau,\varphi^\lambda_x\big\rangle\Big|^p\lesssim\lambda^{p\abs\tau}\qquad\text{uniformemente en el regularizador.}
  \end{equation}
\end{enumerate}

\section{Primer papel del cálculo de Malliavin: hipercontractividad}
Si el ruido es gaussiano, $\big\langle\Pi_x\tau,\varphi^\lambda_x\big\rangle$ es un polinomio en el ruido de grado igual al número de
hojas $\Xi$ de $\tau$, y por tanto vive en una suma finita de caos de Wiener. Por el Teorema~\ref{teo:normas-caos},
\emph{basta probar}~\eqref{eq:modelo-momentos} \emph{para $p=2$}. El segundo momento se calcula con la isometría de las integrales
múltiples, y la estructura de las contracciones (Teorema~\ref{teo:producto}) se codifica en \emph{diagramas de Feynman}.
El teorema BPHZ general de Chandra y Hairer (2016) prueba que, con la renormalización
de Bogoliubov--Parasiuk--Hepp--Zimmermann, todos los diagramas tienen la cota correcta; su
demostración es un tour de force combinatorio de análisis multiescala.

\section{Segundo papel: spectral gap y derivada de Malliavin}

\begin{proposicion}[Desigualdad de spectral gap en $L^p$]
Para $F\in\D^{1,p}$, $p\ge2$, existe $C_p$
(independiente de la dimensión) tal que
\[
  \E\abs{F-\E F}^p\le C_p\,\E\norm{DF}^p_\Hf.
\]
\end{proposicion}

\begin{proof}[Esbozo]
Por la fórmula de Mehler, en forma ``rotacional'' (\emph{desigualdad de Pisier}): para $\Phi$
convexa,
\[
  \E\,\Phi(F-\E F)\le\E\E'\,\Phi\Big(\frac\pi2W'(DF)\Big),
\]
con $W'$ una copia independiente. Se prueba interpolando $F(W)-F(W')=\int_0^{\pi/2}\frac{\dd}{\dd\theta}F(W\sin\theta+
W'\cos\theta)\,\dd\theta$, aplicando Jensen y observando que $(W\sin\theta+W'\cos\theta,W\cos\theta-W'\sin\theta)$ tiene
la misma ley que $(W,W')$. Tomando $\Phi(x)=\abs x^p$, $\E'\abs{W'(DF)}^p=c_p\norm{DF}^p_\Hf$ por ser gaussiana
condicionalmente.
\end{proof}

La idea de Otto, Sauer, Smith y Weber~\cite{OttoSauerSmithWeber25} y de Linares, Otto, Tempelmayr y Tsatsoulis~\cite{LinaresOttoTempelmayrTsatsoulis24}
es aplicar esta desigualdad a $F=\big\langle\Pi_x\tau,\varphi^\lambda_x\big\rangle$ y estimar $\norm{DF}_\Hf$ \emph{por inducción en la homogeneidad},
sin desarrollar en caos. Concretamente:
\begin{itemize}
  \item La derivada de Malliavin $\delta\Pi_x$ del modelo en una dirección $\delta\xi$ del ruido satisface una
  versión linealizada de la ecuación que define $\Pi_x$. Así, $\delta\Pi_x$ es un objeto de la misma
  naturaleza que $\Pi_x$, pero de homogeneidad mayor (porque $\delta\xi\in\Hf$ es más regular que $\xi$).
  \item Por el lema de Hölder--Kolmogorov y los argumentos de recentrado, una cota de $\delta\Pi_x$ en
  una norma adecuada (dual de $\Hf$, es decir, la norma de Cameron--Martin) se obtiene a
  partir de cotas de los $\Pi_x\tau'$ con $\abs{\tau'}<\abs\tau$.
  \item El spectral gap transforma la cota de la derivada en cota de las fluctuaciones de $\Pi_x\tau$,
  y la esperanza se fija con la renormalización (que, en este enfoque, se elige imponiendo
  condiciones de simetría o de anulación de la esperanza en ciertos símbolos).
\end{itemize}

\begin{intuicion}
La derivada de Malliavin del modelo es un \emph{vector tangente} a la ``variedad de soluciones'':
perturbar el ruido en una dirección de Cameron--Martin perturba la solución de manera
más regular que el ruido mismo. Esta ganancia de regularidad es el motor de la inducción.
Frente a los diagramas, el enfoque tiene dos ventajas: funciona automáticamente en todo
el régimen subcrítico (sin análisis caso a caso) y se extiende a ruidos \emph{no gaussianos} que
satisfacen una desigualdad de spectral gap, como los ruidos de ciertos modelos de la
mecánica estadística.
\end{intuicion}

El programa se ha desarrollado en una serie de trabajos: Linares, Otto y Tempelmayr~\cite{LinaresOttoTempelmayr23}
construyeron el grupo de estructura mediante álgebras envolventes universales de álgebras de
Lie pre-Lie; Otto, Sauer, Smith y Weber obtuvieron cotas a priori para EDPEs cuasilineales en
el régimen subcrítico completo; Hairer y Steele reformularon el teorema BPHZ en el lenguaje
de árboles usando el spectral gap. Las notas de Broux, Otto y Tempelmayr~\cite{BrouxOttoTempelmayr25}, publicadas
en 2025, son la introducción pedagógica de referencia: muestran que la derivada de Malliavin
proporciona una caracterización del modelo centrado \emph{estable al eliminar el regularizador} y, junto
con el spectral gap, sus estimaciones estocásticas en todo el régimen subcrítico. Trabajos de
2026 exploran el \emph{recentrado} del modelo directamente con la derivada de Malliavin.

\section{Tercer papel: densidades para EDPEs singulares}
¿Tiene densidad la solución $u(t,x)$ de una EDPE singular? La dificultad es que $u$ no es un
funcional regular del ruido en el sentido clásico, porque la aplicación ``ruido $\mapsto$ solución'' solo es
continua tras levantar el ruido a un modelo, y la renormalización introduce constantes divergentes.
Cannizzaro, Friz y Gassiat~\cite{CannizzaroFrizGassiat17} desarrollaron un cálculo de Malliavin para estructuras de
regularidad en el caso de gPAM en $d=2$: definen las \emph{traslaciones de Cameron--Martin de
modelos}, prueban que el modelo renormalizado trasladado coincide con el modelo del ruido
trasladado (sin renormalización adicional, porque las direcciones de $\Hf$ son regulares) y deducen
la diferenciabilidad de Malliavin de la solución y la existencia de densidad. Gassiat y Labbé
probaron la existencia de densidades para el modelo dinámico $\Phi^4_3$, y Schönbauer desarrolló un
marco general de cálculo de Malliavin para EDPEs singulares que cubre una amplia clase de
ecuaciones. La pregunta de la regularidad $C^\infty$ de las densidades en todo el régimen subcrítico
está esencialmente abierta.

\chapter{Fronteras y aplicaciones recientes}\label{cap:fronteras}

\begin{panorama}
Cerramos con un recorrido por líneas de investigación activas en las que el cálculo de
Malliavin desempeña un papel central: aprendizaje automático (modelos de difusión
generativos y redes neuronales anchas), métodos numéricos (error débil, multinivel), ecuaciones
estocásticas hacia atrás, ecuaciones de McKean--Vlasov, estadística de difusiones, y
una lista de problemas abiertos.
\end{panorama}

\section{Modelos de difusión generativos}
Los modelos generativos basados en difusiones (\emph{score-based diffusion models}) generan muestras
de una ley de datos $p_0$ invirtiendo en el tiempo una difusión que la convierte en ruido gaussiano.
Si $\dd X_t=b(t,X_t)\,\dd t+\sigma(t)\,\dd W_t$ con $X_0\sim p_0$, el proceso invertido $\bar X_s=X_{T-s}$ satisface
\[
  \dd\bar X_s=\big(-b(T-s,\bar X_s)+\sigma\sigma^\top(T-s)\nabla\log p_{T-s}(\bar X_s)\big)\,\dd s+\sigma(T-s)\,\dd\bar W_s,
\]
y la clave es aprender la \emph{función score} $\nabla\log p_t$. El cálculo de Malliavin da una representación
de la score como esperanza condicional.

\begin{proposicion}[Score por integración por partes]
Sea $X_t\in(\D^\infty)^d$ no degenerado, con
densidad $p_t>0$. Entonces
\[
  \nabla\log p_t(x)=-\E\big[H(X_t)\,\big|\,X_t=x\big],\qquad H_i(X_t)=\sum_j\delta\big((\gamma_{X_t}^{-1})_{ij}DX^j_t\big).
\]
\end{proposicion}

\begin{proof}
Por la Proposición~\ref{prop:ibp-iterada} con $G=1$, $\E[\partial_i\varphi(X_t)]=\E[\varphi(X_t)H_i(X_t)]$. El lado
izquierdo es $\int\partial_i\varphi\,p_t=-\int\varphi\,\partial_ip_t$, y el derecho es $\int\varphi(x)\E[H_i\mid X_t=x]p_t(x)\,\dd x$. Igualando para
toda $\varphi$, $\partial_ip_t=-p_t\E[H_i\mid X_t=x]$.
\end{proof}

Cuando $X_0$ es determinista, esta es la fórmula de Bismut; cuando $X_0\sim p_0$, se aplica
condicionalmente. La consecuencia práctica es que la score es la mejor aproximación en $L^2$ de
$-H(X_t)$ por funciones de $X_t$, de modo que puede aprenderse por regresión (es el análogo de la
\emph{denoising score matching}, que corresponde a elegir pesos de Malliavin particulares). Trabajos de
2025 desarrollan esta conexión para difusiones con coeficientes no lineales y ruido degenerado,
donde las fórmulas de tipo gaussiano no se aplican.

\section{Redes neuronales anchas}
Una red neuronal totalmente conectada con pesos gaussianos aleatorios converge, cuando la
anchura $n\to\infty$, a un proceso gaussiano (Neal; Lee et al.; Matthews et al.). Cada salida de la
red es un funcional no lineal de los pesos gaussianos. Basteri y Trevisan (2022) cuantificaron esta
convergencia con técnicas de transporte óptimo, y Favaro, Hanin, Marinucci, Nourdin y Peccati
(2023) obtuvieron tasas del orden de $n^{-1/2}$ en distancias de Wasserstein y total variación con las
desigualdades de Poincaré de segundo orden del Capítulo~\ref{cap:malliavin-stein}. La derivada de Malliavin respecto
a los pesos de la última capa es, esencialmente, el gradiente que usa la retropropagación.

\section{Métodos numéricos}

\paragraph{Error débil.} Para el esquema de Euler $X^n$ de una EDE, el error débil $\E\Phi(X_T)-\E\Phi(X^n_T)$ es
$O(1/n)$ si $\Phi$ es regular (Talay--Tubaro). Bally y Talay probaron, con integración por partes de
Malliavin y la condición de Hörmander, que la tasa $O(1/n)$ se mantiene para $\Phi$ \emph{solo medible y
acotada}, y que la densidad de $X^n_T$ converge a la de $X_T$ a la misma tasa.

\paragraph{Multinivel.} En el método Monte Carlo multinivel de Giles, los estimadores de griegas con
pesos de Malliavin tienen varianza de nivel $O(1/n)$ incluso para pagos discontinuos, manteniendo
la complejidad óptima $O(\varepsilon^{-2})$ en casos donde la diferenciación trayectorial falla.

\paragraph{Cubatura en el espacio de Wiener.} Kusuoka y Lyons--Victoir construyen medidas discretas
en el espacio de trayectorias que reproducen los momentos de las integrales iteradas del browniano
hasta cierto orden. El análisis del error usa la regularización del semigrupo bajo la condición de
Hörmander, probada con cálculo de Malliavin (condición UFG de Kusuoka).

\section{Ecuaciones estocásticas hacia atrás}
Una EDEA es un par adaptado $(Y,Z)$ con $Y_t=\xi+\int_t^Tf(s,Y_s,Z_s)\,\dd s-\int_t^TZ_s\,\dd W_s$. El resultado
clásico de El Karoui, Peng y Quenez es que, si $\xi$ y $f$ son Malliavin-diferenciables,
\[
  Z_t=D_tY_t\qquad\text{(versión adecuada)},
\]
y $D_rY_t$ satisface a su vez una EDEA lineal. Esto da la regularidad de $Z$, crucial para el análisis
de esquemas de discretización (Zhang; Bouchard--Touzi) y para los métodos de aprendizaje
profundo para EDPs parabólicas no lineales en alta dimensión (Han, Jentzen y E), en los que
$Z$ se aproxima con una red neuronal y las fórmulas de tipo Bismut proporcionan estimadores
alternativos.

\section{Ecuaciones de McKean--Vlasov}
En una ecuación de McKean--Vlasov, $\dd X_t=b(X_t,\mathcal L(X_t))\,\dd t+\sigma(X_t,\mathcal L(X_t))\,\dd W_t$, los coeficientes
dependen de la ley de la solución. El cálculo de Malliavin se combina con la \emph{derivada de Lions}
respecto a la medida: Crisan y McMurray obtuvieron fórmulas de integración por partes y la
regularidad de las densidades bajo condiciones de tipo Hörmander, y se han obtenido fórmulas
de Bismut para el semigrupo sobre el espacio de medidas, útiles en juegos de campo medio.

\section{Estadística de difusiones}
Gobet usó el cálculo de Malliavin para probar la propiedad LAN (normalidad asintótica local)
para difusiones elípticas observadas en tiempo discreto, expresando la derivada del logaritmo de
la verosimilitud de transición, que no es explícita, como una esperanza condicional de pesos
de Malliavin. La misma idea da estimadores de máxima verosimilitud aproximados y cotas de
eficiencia.

\section{Problemas abiertos}
\begin{enumerate}
  \item \textbf{Hörmander para $H\le1/4$.} Para el browniano fraccionario con $H\le1/4$ (o procesos
  gaussianos aún más irregulares), la teoría de rough paths requiere levantamientos de orden
  superior que existen pero no son canónicos, y no se sabe probar la existencia de densidad
  regular en general bajo la condición de Hörmander.
  \item \textbf{Densidades $C^\infty$ para EDPEs singulares} en todo el régimen subcrítico: falta un
  análogo infinito-dimensional robusto de la no degeneración de la matriz de Malliavin
  compatible con la renormalización.
  \item \textbf{Hipoelipticidad en dimensión infinita.} Más allá de los resultados de Hairer--Mattingly
  para EDPEs semilineales con ruido degenerado, el caso de ecuaciones cuasilineales o
  singulares está abierto.
  \item \textbf{Tasas óptimas fuera de un caos fijo.} Para funcionales con infinitos caos (como las
  soluciones de EDPEs), las tasas de las desigualdades de Poincaré de segundo orden no
  son siempre óptimas, y caracterizar la tasa óptima es un problema abierto en general.
  \item \textbf{Malliavin--Stein para límites no gaussianos.} Fuera del segundo caos y de las
  leyes Gamma, la teoría para límites en caos superiores o límites mezcla es incompleta:
  por ejemplo, el teorema de convergencia a variables del segundo caos tiene condiciones
  necesarias y suficientes solo en casos particulares.
  \item \textbf{Ruidos no gaussianos.} Extender el programa de spectral gap del Capítulo~\ref{cap:estructuras} a ruidos
  de Poisson o de Lévy, donde la desigualdad de Poincaré vale pero la hipercontractividad
  no, es un objetivo natural.
\end{enumerate}

\section{Lecturas recomendadas}
\begin{itemize}
  \item \emph{Para la teoría general:} Nualart~\cite{Nualart06} (la referencia estándar), Nualart--Nualart~\cite{NualartNualart18} (más
  accesible).
  \item \emph{Para Malliavin--Stein:} Nourdin--Peccati~\cite{NourdinPeccati12} y la reseña~\cite{AzmoodehPeccatiYang21}.
  \item \emph{Para finanzas:} Alòs--García-Lorite~\cite{AlosGarciaLorite21}, Malliavin--Thalmaier~\cite{MalliavinThalmaier06}, Di Nunno--Øksendal--Proske~\cite{DiNunnoOksendalProske09}.
  \item \emph{Para EDPEs:} Sanz-Solé~\cite{SanzSole05}, Khoshnevisan~\cite{Khoshnevisan14}, Walsh~\cite{Walsh86}.
  \item \emph{Para Hörmander:} Hairer~\cite{Hairer11}, una prueba completa en 15 páginas.
  \item \emph{Para rough paths:} Friz--Hairer~\cite{FrizHairer20}.
  \item \emph{Para estructuras de regularidad:} Broux--Otto--Tempelmayr~\cite{BrouxOttoTempelmayr25}.
\end{itemize}


\appendix
\chapter{Formulario}\label{ap:formulario}

\section*{Operadores básicos}
\begin{center}
\renewcommand{\arraystretch}{1.25}
\begin{tabular}{p{4.2cm}p{9.5cm}}
\toprule
Objeto & Fórmula\\
\midrule
Proceso isonormal & $\E[W(h)W(g)]=\ip{h}{g}_\Hf$\\
Derivada (suaves) & $Df(W(h_1),\dots,W(h_n))=\sum_i\partial_if(\cdot)h_i$\\
Dualidad & $\E[F\delta(u)]=\E\ip{DF}{u}_\Hf$\\
Divergencia de $Fh$ & $\delta(Fh)=FW(h)-\ip{DF}{h}$\\
Producto & $\delta(Fu)=F\delta(u)-\ip{DF}{u}$\\
Heisenberg & $D_t\delta(u)=u_t+\delta(D_tu)$\\
Isometría de Skorokhod & $\E\delta(u)^2=\E\norm u^2+\E\iint D_su_tD_tu_s$\\
Ornstein--Uhlenbeck & $L=-\delta D=-\sum_nnJ_n$, $P_t=\sum_ne^{-nt}J_n$\\
Mehler & $P_tF=\E'[F(e^{-t}W+\sqrt{1-e^{-2t}}W')]$\\
Conmutación & $DP_t=e^{-t}P_tD$\\
Covarianza & $\Cov(F,G)=\E\ip{DF}{-DL^{-1}G}$\\
Poincaré & $\Var F\le\E\norm{DF}^2$\\
Log-Sobolev & $\Ent(F^2)\le2\E\norm{DF}^2$\\
Hipercontractividad & $\norm{P_tF}_q\le\norm F_p$ si $e^{2t}\ge\frac{q-1}{p-1}$\\
Clark--Ocone & $F=\E F+\int_0^T\cond{D_tF}{\F_t}\,\dd W_t$\\
\bottomrule
\end{tabular}
\end{center}

\section*{Caos de Wiener}
\begin{center}
\renewcommand{\arraystretch}{1.25}
\begin{tabular}{p{4.2cm}p{9.5cm}}
\toprule
Hermite & $e^{tx-t^2/2}=\sum_n\frac{t^n}{n!}H_n(x)$, $H_n'=nH_{n-1}$\\
Integrales múltiples & $I_n(h^{\otimes n})=H_n(W(h))$, $\norm h=1$\\
Isometría & $\E[I_n(f)I_m(g)]=\delta_{nm}\,n!\,\ip{\tilde f}{\tilde g}$\\
Derivada & $D_tI_n(f)=nI_{n-1}(f(\cdot,t))$\\
Divergencia & $\delta(I_n(f(\cdot,t)))=I_{n+1}(\tilde f)$\\
Stroock & $f_n=\frac1{n!}\E[D^nF]$\\
Producto & $I_p(f)I_q(g)=\sum_rr!\binom pr\binom qrI_{p+q-2r}(f\otimes_rg)$\\
Momentos en un caos & $\norm F_p\le(p-1)^{n/2}\norm F_2$, $F\in\Hc_n$\\
\bottomrule
\end{tabular}
\end{center}

\section*{Densidades y finanzas}
\begin{center}
\renewcommand{\arraystretch}{1.35}
\begin{tabular}{p{4.2cm}p{9.5cm}}
\toprule
Densidad, $d=1$ & $p_F(x)=\E\big[\ind_{\{F>x\}}\delta(DF/\norm{DF}^2)\big]$\\
Nourdin--Viens & $p_F(x)=\frac{\E\abs F}{2g_F(x)}\exp(-\int_0^x\frac{y}{g_F(y)}\,\dd y)$\\
Derivada de una EDE & $D_rX_t=Y_tZ_r\sigma(X_r)$\\
Matriz de Malliavin & $\gamma_{X_t}=Y_tC_tY_t^\top$, $C_t=\int_0^tZ_r\sigma\sigma^\top(X_r)Z_r^\top\,\dd r$\\
Delta BS & $\frac{e^{-rT}}{S_0\sigma T}\E[\Phi(S_T)W_T]$\\
Vega BS & $e^{-rT}\E[\Phi(S_T)(\frac{W_T^2}{\sigma T}-W_T-\frac1\sigma)]$\\
Bismut--Elworthy--Li & $\nabla\E\Phi(X_T^x)=\E[\Phi(X_T)\frac1T\int_0^T(\sigma^{-1}(X_t)Y_t)^\top\,\dd W_t]$\\
Cobertura & $\phi_t=\E_\Q[\tilde D_t\tilde H\mid\F_t]/(\sigma_t\tilde S_t)$\\
Alòs & $V_t=\E_t[BS(t,X_t,v_t)]+\frac\rho2\E_t\int_t^Te^{-r(s-t)}H\Lambda_s\,\dd s$\\
\bottomrule
\end{tabular}
\end{center}

\chapter{Guía de notación}\label{ap:notacion}

\begin{center}
\renewcommand{\arraystretch}{1.25}
\begin{tabular}{p{3.5cm}p{10cm}}
\toprule
$\Hf$ & espacio de Hilbert separable subyacente al proceso isonormal\\
$W(h)$ & proceso isonormal; integral de Wiener en el caso browniano\\
$\Sc$, $\Sc_b$, $\mathcal P$ & variables suaves, suaves acotadas, polinomiales\\
$\D^{k,p}$, $\D^\infty$ & espacios de Sobolev--Watanabe\\
$D$, $\delta$, $L$ & derivada, divergencia, generador de Ornstein--Uhlenbeck\\
$J_n$, $\Hc_n$ & proyección sobre el $n$-ésimo caos, $n$-ésimo caos\\
$I_n(f)$ & integral múltiple de orden $n$\\
$f\otimes_rg$, $\tilde f$ & contracción de orden $r$, simetrización\\
$\gamma_F$ & matriz de Malliavin $(\ip{DF^i}{DF^j})_{ij}$\\
$Y_t$, $Z_t$ & proceso tangente $\nabla_xX_t$ y su inversa\\
$G(t,x)$ & solución fundamental (núcleo del calor o de ondas)\\
$\mu$, $\gamma$ & medida espectral y covarianza espacial del ruido\\
$d_{TV}$, $d_W$ & distancias de variación total y de Wasserstein\\
$\law$, $\tolaw$ & igualdad y convergencia en ley\\
\bottomrule
\end{tabular}
\end{center}


\backmatter
\begin{thebibliography}{99}
\addcontentsline{toc}{chapter}{Bibliografía}
\markboth{BIBLIOGRAFÍA}{BIBLIOGRAFÍA}

\bibitem{Alos06} E. Alòs. A generalization of the Hull and White formula with applications to option pricing
approximation. \emph{Finance Stoch.} 10 (2006), 353--365.

\bibitem{AlosLeonVives07} E. Alòs, J. A. León, J. Vives. On the short-time behavior of the implied volatility for jump-diffusion
models with stochastic volatility. \emph{Finance Stoch.} 11 (2007), 571--589.

\bibitem{AlosGarciaLorite21} E. Alòs, D. García-Lorite. \emph{Malliavin Calculus in Finance: Theory and Practice}. Chapman \&
Hall/CRC, 2021.

\bibitem{AzmoodehPeccatiYang21} E. Azmoodeh, G. Peccati, X. Yang. Malliavin--Stein method: a survey of some recent developments.
\emph{Mod. Stoch. Theory Appl.} 8 (2021), 141--177.

\bibitem{BallyCaramellino11} V. Bally, L. Caramellino. Riesz transform and integration by parts formulas for random variables.
\emph{Stoch. Proc. Appl.} 121 (2011), 1332--1355.

\bibitem{BallyCaramellinoZanette05} V. Bally, L. Caramellino, A. Zanette. Pricing and hedging American options by Monte Carlo
methods using a Malliavin calculus approach. \emph{Monte Carlo Methods Appl.} 11 (2005), 97--133.

\bibitem{BallyPardoux98} V. Bally, E. Pardoux. Malliavin calculus for white noise driven parabolic SPDEs. \emph{Potential Anal.} 9
(1998), 27--64.

\bibitem{BayerFrizGatheral16} C. Bayer, P. Friz, J. Gatheral. Pricing under rough volatility. \emph{Quant. Finance} 16 (2016), 887--904.

\bibitem{Bismut81} J.-M. Bismut. Martingales, the Malliavin calculus and hypoellipticity under general Hörmander's
conditions. \emph{Z. Wahrsch. Verw. Gebiete} 56 (1981), 469--505.

\bibitem{BouleauHirsch91} N. Bouleau, F. Hirsch. \emph{Dirichlet Forms and Analysis on Wiener Space}. De Gruyter, 1991.

\bibitem{BrouxOttoTempelmayr25} L. Broux, F. Otto, M. Tempelmayr. Lecture notes on Malliavin calculus in regularity structures.
\emph{Stoch. PDE: Anal. Comp.} (2025); arXiv:2401.05935.

\bibitem{BrunedHairerZambotti19} Y. Bruned, M. Hairer, L. Zambotti. Algebraic renormalisation of regularity structures. \emph{Invent.
Math.} 215 (2019), 1039--1156.

\bibitem{CannizzaroFrizGassiat17} G. Cannizzaro, P. K. Friz, P. Gassiat. Malliavin calculus for regularity structures: the case of gPAM.
\emph{J. Funct. Anal.} 272 (2017), 363--419.

\bibitem{CassFriz10} T. Cass, P. Friz. Densities for rough differential equations under Hörmander's condition. \emph{Ann. of
Math.} 171 (2010), 2115--2141.

\bibitem{CassHairerLittererTindel15} T. Cass, M. Hairer, C. Litterer, S. Tindel. Smoothness of the density for solutions to Gaussian
rough differential equations. \emph{Ann. Probab.} 43 (2015), 188--239.

\bibitem{ChenKhoshnevisanNualartPu21} L. Chen, D. Khoshnevisan, D. Nualart, F. Pu. Spatial ergodicity for SPDEs via Poincaré-type
inequalities. \emph{Electron. J. Probab.} 26 (2021), 1--37.

\bibitem{ChenHuNualart21} L. Chen, Y. Hu, D. Nualart. Regularity and strict positivity of densities for the nonlinear stochastic
heat equation. \emph{Mem. Amer. Math. Soc.} 273 (2021).

\bibitem{DaPratoZabczyk14} G. Da Prato, J. Zabczyk. \emph{Stochastic Equations in Infinite Dimensions}, 2.ª ed. Cambridge University
Press, 2014.

\bibitem{Dalang99} R. C. Dalang. Extending the martingale measure stochastic integral with applications to spatially
homogeneous s.p.d.e.'s. \emph{Electron. J. Probab.} 4 (1999), 1--29.

\bibitem{DalangKhoshnevisanNualart07} R. Dalang, D. Khoshnevisan, E. Nualart. Hitting probabilities for systems of non-linear stochastic
heat equations with additive noise. \emph{ALEA} 3 (2007), 231--271.

\bibitem{DiNunnoOksendalProske09} G. Di Nunno, B. Øksendal, F. Proske. \emph{Malliavin Calculus for Lévy Processes with Applications to
Finance}. Springer, 2009.

\bibitem{ElEuchRosenbaum19} O. El Euch, M. Rosenbaum. The characteristic function of rough Heston models. \emph{Math. Finance} 29
(2019), 3--38.

\bibitem{Fournie99} E. Fournié, J.-M. Lasry, J. Lebuchoux, P.-L. Lions, N. Touzi. Applications of Malliavin calculus to
Monte Carlo methods in finance. \emph{Finance Stoch.} 3 (1999), 391--412.

\bibitem{Fournie01} E. Fournié, J.-M. Lasry, J. Lebuchoux, P.-L. Lions. Applications of Malliavin calculus to Monte
Carlo methods in finance II. \emph{Finance Stoch.} 5 (2001), 201--236.

\bibitem{FrizHairer20} P. K. Friz, M. Hairer. \emph{A Course on Rough Paths}, 2.ª ed. Springer, 2020.

\bibitem{GatheralJaissonRosenbaum18} J. Gatheral, T. Jaisson, M. Rosenbaum. Volatility is rough. \emph{Quant. Finance} 18 (2018), 933--949.

\bibitem{Gross75} L. Gross. Logarithmic Sobolev inequalities. \emph{Amer. J. Math.} 97 (1975), 1061--1083.

\bibitem{Hairer14} M. Hairer. A theory of regularity structures. \emph{Invent. Math.} 198 (2014), 269--504.

\bibitem{Hairer11} M. Hairer. On Malliavin's proof of Hörmander's theorem. \emph{Bull. Sci. Math.} 135 (2011), 650--666.

\bibitem{HairerMattingly11} M. Hairer, J. C. Mattingly. A theory of hypoellipticity and unique ergodicity for semilinear stochastic
PDEs. \emph{Electron. J. Probab.} 16 (2011), 658--738.

\bibitem{HuNualartSong11} Y. Hu, D. Nualart, J. Song. Feynman--Kac formula for heat equation driven by fractional white
noise. \emph{Ann. Probab.} 39 (2011), 291--326.

\bibitem{HuangNualartViitasaari20} J. Huang, D. Nualart, L. Viitasaari. A central limit theorem for the stochastic heat equation. \emph{Stoch.
Proc. Appl.} 130 (2020), 7170--7184.

\bibitem{HuangNualartViitasaariZheng20} J. Huang, D. Nualart, L. Viitasaari, G. Zheng. Gaussian fluctuations for the stochastic heat equation
with colored noise. \emph{Stoch. PDE: Anal. Comp.} 8 (2020), 402--421.

\bibitem{Imkeller03} P. Imkeller. Malliavin's calculus in insider models: additional utility and free lunches. \emph{Math. Finance}
13 (2003), 153--169.

\bibitem{Inahama14} Y. Inahama. Malliavin differentiability of solutions of rough differential equations. \emph{J. Funct. Anal.}
267 (2014), 1566--1584.

\bibitem{KaratzasShreve91} I. Karatzas, S. Shreve. \emph{Brownian Motion and Stochastic Calculus}, 2.ª ed. Springer, 1991.

\bibitem{Khoshnevisan14} D. Khoshnevisan. \emph{Analysis of Stochastic Partial Differential Equations}. CBMS 119, AMS, 2014.

\bibitem{Kunita90} H. Kunita. \emph{Stochastic Flows and Stochastic Differential Equations}. Cambridge University Press,
1990.

\bibitem{KusuokaStroock} S. Kusuoka, D. Stroock. Applications of the Malliavin calculus I, II, III. \emph{Stochastic Analysis}
(Katata/Kyoto, 1982); \emph{J. Fac. Sci. Univ. Tokyo} 32 (1985); 34 (1987).

\bibitem{LastPenrose17} G. Last, M. Penrose. \emph{Lectures on the Poisson Process}. Cambridge University Press, 2017.

\bibitem{LastPeccatiSchulte16} G. Last, G. Peccati, M. Schulte. Normal approximation on Poisson spaces: Mehler's formula, second
order Poincaré inequalities and stabilization. \emph{Probab. Theory Relat. Fields} 165 (2016), 667--723.

\bibitem{Malliavin78} P. Malliavin. Stochastic calculus of variations and hypoelliptic operators. \emph{Proc. Intern. Symp. SDE}
(Kyoto 1976), Wiley, 1978, 195--263.

\bibitem{MalliavinThalmaier06} P. Malliavin, A. Thalmaier. \emph{Stochastic Calculus of Variations in Mathematical Finance}. Springer,
2006.

\bibitem{MuellerNualart08} C. Mueller, D. Nualart. Regularity of the density for the stochastic heat equation. \emph{Electron. J.
Probab.} 13 (2008), 2248--2258.

\bibitem{Norris86} J. Norris. Simplified Malliavin calculus. \emph{Séminaire de Probabilités XX}, LNM 1204, Springer, 1986,
101--130.

\bibitem{NourdinPeccati12} I. Nourdin, G. Peccati. \emph{Normal Approximations with Malliavin Calculus: From Stein's Method to
Universality}. Cambridge University Press, 2012.

\bibitem{NourdinPeccatiReinert09} I. Nourdin, G. Peccati, G. Reinert. Second order Poincaré inequalities and CLTs on Wiener space.
\emph{J. Funct. Anal.} 257 (2009), 593--609.

\bibitem{NourdinViens09} I. Nourdin, F. Viens. Density formula and concentration inequalities with Malliavin calculus.
\emph{Electron. J. Probab.} 14 (2009), 2287--2309.

\bibitem{Nualart06} D. Nualart. \emph{The Malliavin Calculus and Related Topics}, 2.ª ed. Springer, 2006.

\bibitem{NualartNualart18} D. Nualart, E. Nualart. \emph{Introduction to Malliavin Calculus}. Cambridge University Press, 2018.

\bibitem{NualartPeccati05} D. Nualart, G. Peccati. Central limit theorems for sequences of multiple stochastic integrals. \emph{Ann.
Probab.} 33 (2005), 177--193.

\bibitem{NualartOrtizLatorre08} D. Nualart, S. Ortiz-Latorre. Central limit theorems for multiple stochastic integrals and Malliavin
calculus. \emph{Stoch. Proc. Appl.} 118 (2008), 614--628.

\bibitem{NualartQuerSardanyons07} D. Nualart, L. Quer-Sardanyons. Existence and smoothness of the density for spatially homogeneous
SPDEs. \emph{Potential Anal.} 27 (2007), 281--299.

\bibitem{NualartZheng20} D. Nualart, G. Zheng. Averaging Gaussian functionals. \emph{Electron. J. Probab.} 25 (2020), 1--54.

\bibitem{OconeKaratzas91} D. Ocone, I. Karatzas. A generalized Clark representation formula, with application to optimal
portfolios. \emph{Stochastics} 34 (1991), 187--220.

\bibitem{OttoSauerSmithWeber25} F. Otto, J. Sauer, S. Smith, H. Weber. A priori bounds for quasi-linear SPDEs in the full sub-critical
regime. \emph{J. Eur. Math. Soc.} (2025); arXiv:2103.11039.

\bibitem{LinaresOttoTempelmayr23} P. Linares, F. Otto, M. Tempelmayr. The structure group for quasi-linear equations via universal
enveloping algebras. \emph{Comm. Amer. Math. Soc.} 3 (2023), 1--64.

\bibitem{LinaresOttoTempelmayrTsatsoulis24} P. Linares, F. Otto, M. Tempelmayr, P. Tsatsoulis. A diagram-free approach to the stochastic
estimates in regularity structures. \emph{Invent. Math.} (2024); arXiv:2112.10739.

\bibitem{PeccatiTudor05} G. Peccati, C. Tudor. Gaussian limits for vector-valued multiple stochastic integrals. \emph{Sém. Probab.}
XXXVIII, LNM 1857, 2005, 247--262.

\bibitem{Pisier88} G. Pisier. Riesz transforms: a simpler analytic proof of P.-A. Meyer's inequality. \emph{Sém. Probab.} XXII,
LNM 1321, 1988, 485--501.

\bibitem{RevuzYor99} D. Revuz, M. Yor. \emph{Continuous Martingales and Brownian Motion}, 3.ª ed. Springer, 1999.

\bibitem{SanzSole05} M. Sanz-Solé. \emph{Malliavin Calculus with Applications to Stochastic Partial Differential Equations}.
EPFL Press, 2005.

\bibitem{Walsh86} J. B. Walsh. An introduction to stochastic partial differential equations. \emph{École d'Été de Probabilités
de Saint-Flour XIV}, LNM 1180, Springer, 1986, 265--439.
\end{thebibliography}

\end{document}
